\documentclass[12pt,a4paper]{report}

\usepackage[utf8]{inputenc}
\usepackage[T1]{fontenc}
\usepackage{underscore}
\usepackage{listings}
\lstdefinelanguage{Solidity}{
  keywords={pragma,solidity,contract,function,returns,external,internal,public,private,
             view,pure,payable,mapping,address,uint256,uint,bool,bytes32,string,
             event,emit,modifier,require,revert,import,is,memory,storage,indexed,
             bytes,keccak256,struct,enum,constructor,fallback,receive},
  sensitive=true,
  comment=[l]{//},
  morecomment=[s]{/*}{*/},
  morestring=[b]",
}
\lstset{
  basicstyle=\ttfamily\scriptsize,
  keywordstyle=\bfseries,
  commentstyle=\color{black!55}\itshape,
  stringstyle=\itshape,
  breaklines=true,
  frame=tb,
  rulecolor=\color{black!55},
  captionpos=b,
  numbers=left,
  numberstyle=\tiny\color{black!55},
  xleftmargin=2em,
}
\IfFileExists{libertine.sty}{%
  \usepackage[mono=false]{libertine}%
  \IfFileExists{zi4.sty}{\usepackage[varqu]{zi4}}{}%
}{%
  \usepackage{lmodern}%
}
\usepackage[protrusion=true,expansion=false]{microtype}
\usepackage[a4paper,width=155mm,top=25mm,bottom=25mm,headheight=15pt]{geometry}
\usepackage{setspace}
\onehalfspacing
\emergencystretch=3em
\hbadness=2000

\usepackage{amsmath,amssymb}
\usepackage{fontawesome5}
\newcommand{\tmarkDone}{\textbf{\faCheck}}
\newcommand{\tmarkNo}{\textbf{\faTimes}}
\newcommand{\tmarkPartial}{\faHourglassHalf}
\newcommand{\tmarkPlanned}{\textcolor{black!55}{\faCircle}}

\usepackage{graphicx}
\usepackage{pdflscape} % full-page landscape for wide ER diagrams
% Self-contained build: figure assets in Diagrams/ (SVG sources; LaTeX includes generated PDF).
% Rebuild PDFs:  ./Diagrams/build-pdfs.sh
% Compile from this directory:  pdflatex Pre-thesis_v33.tex (twice)
\newcommand*{\ThesisDiagDir}{Diagrams/}
\newcommand*{\ThesisNewDiagDir}{mermaid-diagrams/pdf/}
\graphicspath{{\ThesisNewDiagDir}{./\ThesisNewDiagDir}{\ThesisDiagDir}{./\ThesisDiagDir}{./}{Tables/}}
\pdfimageresolution=300
\makeatletter
\newcommand*{\@thesis@doinclude}[2]{%
  \if\relax\detokenize{#1}\relax
    \includegraphics{#2}%
  \else
    \includegraphics[#1]{#2}%
  \fi
}
\newcommand*{\@thesismissingimg}[1]{%
  \fbox{\parbox{0.92\linewidth}{\scriptsize\centering\ttfamily Missing diagram: \detokenize{\ThesisDiagDir#1}.\\[2pt]Add \texttt{#1} (from matching \texttt{.svg} via \texttt{build-pdfs.sh}) to \texttt{Diagrams/} and recompile.}}%
}
\newcommand*{\@thesis@loadfig}[2]{%
  \IfFileExists{\ThesisDiagDir#2}{%
    \@thesis@doinclude{#1}{\ThesisDiagDir#2}}{%
  \IfFileExists{#2}{%
    \@thesis@doinclude{#1}{#2}}{%
  \IfFileExists{Tables/#2}{%
    \@thesis@doinclude{#1}{Tables/#2}}{%
    \@thesismissingimg{#2}%
  }}}%
}
\newcommand{\ThesisIncludeGraphics}[2][]{%
  \IfFileExists{\ThesisNewDiagDir#2.pdf}{%
    \@thesis@doinclude{#1}{\ThesisNewDiagDir#2.pdf}}{%
  \IfFileExists{\ThesisDiagDir#2.pdf}{%
    \@thesis@doinclude{#1}{\ThesisDiagDir#2.pdf}}{%
  \@thesis@loadfig{#1}{#2}%
  }}%
}
% Shorthand: \Diagram{stem} -> Diagrams/stem.pdf
\newcommand{\Diagram}[2][]{%
  \if\relax\detokenize{#1}\relax
    \ThesisIncludeGraphics{#2}%
  \else
    \ThesisIncludeGraphics[#1]{#2}%
  \fi
}
\makeatother
\newcommand{\TableScreenshotEnd}[2][]{%
  \refstepcounter{table}%
  \if\relax\detokenize{#1}\relax\else\label{#1}\fi%
  \phantomsection%
  \addcontentsline{lot}{table}{\protect\numberline{\thetable}#2}%
}
\usepackage[table]{xcolor}
\usepackage{tikz}
\usetikzlibrary{arrows.meta,positioning,fit,shapes.geometric,calc}
\usepackage{pgfplots}
\usepgfplotslibrary{statistics}
\pgfplotsset{compat=1.18, minor tick num=0}

\usepackage{booktabs}
\usepackage{array}
\usepackage{longtable}
\usepackage{colortbl}
\usepackage{multirow}
\usepackage{float}
\usepackage{etoolbox}
\usepackage{caption}
\captionsetup{
    labelfont={bf,small},
    textfont={small},
    skip=6pt,
    format=hang
}
\makeatletter
% Compact, professional List of Figures / List of Tables (no tocloft dependency)
\renewcommand*\l@figure{\@dottedtocline{1}{0em}{3.2em}}
\renewcommand*\l@table{\@dottedtocline{1}{0em}{3.8em}}
\makeatother

\newlength{\ThesisTableImgReserve}
\newlength{\ThesisFigImgReserve}
\setlength{\ThesisTableImgReserve}{\dimexpr2\intextsep+0.75\baselineskip\relax}
\setlength{\ThesisFigImgReserve}{\dimexpr2\intextsep+\abovecaptionskip+\belowcaptionskip+4\baselineskip\relax}
\newcommand{\TableImageMaxFit}[1]{%
  \ThesisIncludeGraphics[
    width=0.7\linewidth,
    height=\dimexpr(\textheight-\ThesisTableImgReserve)*7/10\relax,
    keepaspectratio
  ]{#1}%
}
\newcommand{\FigureImageMaxFit}[1]{%
  \ThesisIncludeGraphics[width=\linewidth,height=\dimexpr\textheight-\ThesisFigImgReserve\relax,keepaspectratio]{#1}%
}
\newcommand{\OnePageDiagram}[2][]{%
  \if\relax\detokenize{#1}\relax
    \ThesisIncludeGraphics[width=\linewidth,height=\dimexpr\textheight-\ThesisFigImgReserve\relax,keepaspectratio]{#2}%
  \else
    \ThesisIncludeGraphics[#1,width=\linewidth,height=\dimexpr\textheight-\ThesisFigImgReserve\relax,keepaspectratio]{#2}%
  \fi
}
\newcommand{\HalfWidthDiagram}[1]{%
  \ThesisIncludeGraphics[width=0.5\linewidth,keepaspectratio]{#1}%
}
\newcommand{\WideDiagram}[2][]{%
  \if\relax\detokenize{#1}\relax
    \ThesisIncludeGraphics[width=\linewidth,height=0.78\textheight,keepaspectratio]{#2}%
  \else
    \ThesisIncludeGraphics[#1,width=\linewidth,height=0.78\textheight,keepaspectratio]{#2}%
  \fi
}
\newcommand{\TallDiagram}[2][]{%
  \if\relax\detokenize{#1}\relax
    \ThesisIncludeGraphics[width=\linewidth,height=0.96\textheight,keepaspectratio]{#2}%
  \else
    \ThesisIncludeGraphics[#1,width=\linewidth,height=0.96\textheight,keepaspectratio]{#2}%
  \fi
}
% Balanced layout slot (reference: Figure 4.3 ml-explainability) — wide banded diagrams.
\newcommand{\BalancedDiagram}[2][]{%
  \if\relax\detokenize{#1}\relax
    \ThesisIncludeGraphics[width=\linewidth,height=0.72\textheight,keepaspectratio]{#2}%
  \else
    \ThesisIncludeGraphics[#1,height=0.72\textheight,keepaspectratio]{#2}%
  \fi
}
\newcommand{\LandscapeDiagram}[2][]{%
  \begin{landscape}%
  \begin{figure}[H]%
  \centering%
  \if\relax\detokenize{#1}\relax
    \ThesisIncludeGraphics[width=\linewidth,height=0.88\textheight,keepaspectratio]{#2}%
  \else
    \ThesisIncludeGraphics[#1,width=\linewidth,height=0.88\textheight,keepaspectratio]{#2}%
  \fi
}
\newcommand{\LandscapeDiagramEnd}{%
  \end{figure}%
  \end{landscape}%
}
% Sequence diagrams — uniform 2x scale (text + arrows); landscape slot avoids vertical overflow.
\newcommand{\SequenceDiagram}[1]{%
  \begin{landscape}%
  \begin{figure}[H]%
  \centering%
  \scalebox{2}{\ThesisIncludeGraphics[width=0.5\linewidth,keepaspectratio]{#1}}%
}
\newcommand{\SequenceDiagramEnd}{%
  \end{figure}%
  \end{landscape}%
}
% Activity diagrams — full-height slot for 3× label size.
\newcommand{\ActivityDiagram}[2][]{%
  \if\relax\detokenize{#1}\relax
    \ThesisIncludeGraphics[width=\linewidth,height=0.88\textheight,keepaspectratio]{#2}%
  \else
    \ThesisIncludeGraphics[#1,height=0.88\textheight,keepaspectratio]{#2}%
  \fi
}
% Half-size slot (50\% of BalancedDiagram width and height).
\newcommand{\CompactDiagram}[2][]{%
  \if\relax\detokenize{#1}\relax
    \ThesisIncludeGraphics[width=0.5\linewidth,height=0.36\textheight,keepaspectratio]{#2}%
  \else
    \ThesisIncludeGraphics[#1,width=0.5\linewidth,height=0.36\textheight,keepaspectratio]{#2}%
  \fi
}

\usepackage[skins]{tcolorbox}

\usepackage{titlesec}
\usepackage{parskip}
\usepackage{fancyhdr}
\usepackage{enumitem}
\PassOptionsToPackage{hyphens,spaces,obeyspaces}{url}
\usepackage{hyperref}
\makeatletter
\g@addto@macro\UrlBreaks{\do\/\do\-\do\.\do\:\do\&\do\=\do\#}
\makeatother

\definecolor{PrimaryBlue}{HTML}{000000}
\definecolor{AccentBlue}{gray}{0.55}
\definecolor{LightBlue}{gray}{0.96}
\definecolor{MidBlue}{gray}{0.90}
\definecolor{DarkText}{HTML}{000000}
\definecolor{GrayText}{gray}{0.45}
\definecolor{RowShade}{gray}{0.965}
\definecolor{RowShadeAlt}{gray}{0.92}
\definecolor{TableHeadBg}{gray}{0.88}
\definecolor{TableRule}{gray}{0.55}
\definecolor{GreenAccent}{HTML}{000000}
\definecolor{AmberAccent}{HTML}{000000}
\definecolor{RedAccent}{HTML}{000000}
\definecolor{VioletAccent}{HTML}{000000}

\hypersetup{
    colorlinks=true,
    linkcolor=black,
    filecolor=black,
    urlcolor=black,
    citecolor=black,
    pdfborder={0 0 0},
}

\titleformat{\chapter}[display]
    {\normalfont\LARGE\bfseries}
    {\chaptertitlename\ \thechapter}{20pt}{\LARGE}
    [\vspace{4pt}{\color{black!40}\titlerule[0.6pt]}]

\titleformat{\section}
    {\large\bfseries}
    {\thesection}{0.8em}{}

\titleformat{\subsection}
    {\normalsize\bfseries}
    {\thesubsection}{0.6em}{}

\titleformat{\subsubsection}
    {\normalsize\bfseries\itshape}
    {\thesubsubsection}{0.6em}{}

\titlespacing*{\chapter}{0pt}{-20pt}{20pt}
\titlespacing*{\section}{0pt}{16pt plus 3pt minus 2pt}{8pt plus 2pt minus 1pt}
\titlespacing*{\subsection}{0pt}{12pt plus 2pt minus 1pt}{6pt plus 1pt minus 1pt}
\titlespacing*{\subsubsection}{0pt}{10pt plus 2pt minus 1pt}{4pt plus 1pt minus 1pt}

\newcommand{\tableheadcolor}{\rowcolor{TableHeadBg}}
\newcommand{\tableheadfont}{\bfseries}
\newcommand{\ThesisTableBodyFont}{\small}

% Consistent greyscale table polish (font size unchanged; spacing only).
\newcommand{\ThesisTablePrep}{%
  \ThesisTableBodyFont
  \renewcommand{\arraystretch}{1.22}%
  \setlength{\tabcolsep}{6pt}%
  \arrayrulecolor{TableRule}%
  \setlength{\aboverulesep}{0.55ex}%
  \setlength{\belowrulesep}{0.55ex}%
}
\AtBeginEnvironment{table}{\ThesisTablePrep}
\AtBeginEnvironment{tabular}{%
  \rowcolors{2}{}{RowShade}%
  \arrayrulecolor{TableRule}%
}

\newenvironment{moderntable}[1][H]{%
    \begin{table}[#1]\centering\ThesisTablePrep
}{%
    \end{table}%
}

\newtcolorbox{layerbox}[1]{
    colback=white,
    colframe=black!55,
    coltitle=black,
    fonttitle=\bfseries\small,
    title=#1,
    boxrule=0.6pt,
    arc=2pt,
    left=6pt,right=6pt,top=4pt,bottom=4pt,
}

\newtcolorbox{highlightbox}[1][]{
    colback=black!4,
    colframe=black!45,
    boxrule=0.5pt,
    arc=2pt,
    left=8pt,right=8pt,top=6pt,bottom=6pt,
    #1
}

\pagestyle{fancy}
\fancyhf{}
\renewcommand{\headrulewidth}{0.4pt}
\fancyhead[L]{\small\color{black!55}Crypto World Bank}
\fancyhead[R]{\small\color{black!55}BRAC University}
\fancyfoot[C]{\small\color{black!55}\thepage}

\newcommand{\manualpie}[6]{%
  \fill[#6] (#1,#2) -- ++({#4}:{#3}) arc ({#4}:{#5}:{#3}) -- cycle;
  \draw[white, line width=0.8pt] (#1,#2) -- ++({#4}:{#3}) arc ({#4}:{#5}:{#3}) -- cycle;
}
\newcommand{\ddcat}[3]{\node[draw=black!55,fill=black!5,rounded corners=2pt,
  minimum width=2.2cm,minimum height=0.65cm,align=center,font=\scriptsize]
  (#1) at (#2) {#3};}
\newcommand{\ddone}[3]{\node[draw=black,fill=black!22,
  minimum width=2.2cm,minimum height=0.65cm,align=center,
  font=\scriptsize\bfseries,rounded corners=2pt] (#1) at #2 {#3};}
\newcommand{\ddtwo}[3]{\node[draw=black!55,fill=white,
  minimum width=2.2cm,minimum height=0.65cm,align=center,
  font=\scriptsize,rounded corners=2pt] (#1) at #2 {#3};}

\begin{document}

%% ============================================================
%% TITLE PAGE
%% ============================================================
\thispagestyle{empty}

\begin{center}

\vspace*{0.3cm}

\ThesisIncludeGraphics[width=3.5cm]{bracu\_logo\_12-0-2022}

\vspace{0.8cm}

{\LARGE\bfseries\color{PrimaryBlue} Decentralized Crypto World Bank}\\[0.4em]
{\large\color{PrimaryBlue!80} A Blockchain-Based Banking Platform\\[0.15em]with AI-Enhanced Security and Assistance}

\vspace{1.2cm}

{\large by}

\vspace{0.5cm}

{\Large\bfseries\color{PrimaryBlue} Md.\ Bokhtiar Rahman Juboraz}\\[0.2em]
{\color{GrayText}20301138}\\[0.8em]
{\Large\bfseries\color{PrimaryBlue} Md.\ Mahir Ahnaf Ahmed}\\[0.2em]
{\color{GrayText}20301083}

\vspace{1.2cm}

A pre-thesis 1 report submitted to the Department of Computer Science and Engineering\\
in partial fulfillment of the requirements for the degree of\\
\textbf{B.Sc.\ in Computer Science}

\vfill

{\color{PrimaryBlue}\rule{0.4\textwidth}{0.6pt}}\\[0.6em]
Department of Computer Science and Engineering\\
\textbf{BRAC University}\\
February 2026

\vspace{0.5cm}

{\small\color{GrayText}\copyright\ 2026.\ BRAC University --- All rights reserved.}

\end{center}

%% ============================================================
%% DECLARATION
%% ============================================================
\newpage
\thispagestyle{empty}

\chapter*{Declaration}
\phantomsection
\addcontentsline{toc}{chapter}{Declaration}

It is hereby declared that

\begin{enumerate}
\item The report submitted is our own original work while completing degree at BRAC University.

\item The report does not contain material previously published or written by a third party, except where this is appropriately cited through full and accurate referencing.

\item The report does not contain material which has been accepted, or submitted, for any other degree or diploma at a university or other institution.

\item We have acknowledged all main sources of help.
\end{enumerate}

\vspace{2cm}

\textbf{Student's Full Name \& Signature:}

\vspace{2cm}

\begin{tabular}{p{6cm}p{1cm}p{6cm}}
\hrulefill & & \hrulefill \\
\centering Md. Bokhtiar Rahman Juboraz & & \centering Md. Mahir Ahnaf Ahmed \tabularnewline
\centering 20301138 & & \centering 20301083 \tabularnewline
\end{tabular}

%% ============================================================
%% APPROVAL
%% ============================================================
\newpage
\thispagestyle{empty}

\chapter*{Approval}
\phantomsection
\addcontentsline{toc}{chapter}{Approval}

The pre-thesis 1 report of final year project titled ``Decentralized Crypto World Bank: A Blockchain-Based Banking Platform with AI-Enhanced Security and Assistance'' submitted by

\begin{enumerate}
\item Md. Bokhtiar Rahman Juboraz (20301138)
\item Md. Mahir Ahnaf Ahmed (20301083)
\end{enumerate}

Of Spring, 2026 has been accepted as satisfactory in partial fulfillment of the requirement for the degree of B.Sc. in Computer Science on February 20, 2026.

\vspace{1cm}

\textbf{Examining Committee:}

\vspace{1.5cm}

\begin{flushleft}
Supervisor:\\
(Member)
\end{flushleft}

\begin{flushright}
\vspace{-0.5cm}
\hrulefill\\
Mr. Annajiat Alim Rasel\\
Senior Lecturer\\
Department of Computer Science and Engineering\\
BRAC University
\end{flushright}

\vspace{1.5cm}

\begin{flushleft}
Project Coordinator:\\
(Member)
\end{flushleft}

\begin{flushright}
\vspace{-0.5cm}
\hrulefill\\
Dr. Md. Golam Rabiul Alam\\
Professor\\
Department of Computer Science and Engineering\\
BRAC University
\end{flushright}

\newpage
\thispagestyle{empty}

\vspace*{\fill}

\begin{flushleft}
Head of Department:\\
(Chair)
\end{flushleft}

\begin{flushright}
\vspace{-0.5cm}
\hrulefill\\
Dr. Sadia Hamid Kazi\\
Chairperson and Associate Professor\\
Department of Computer Science and Engineering\\
BRAC University
\end{flushright}

\vspace*{\fill}

%% ============================================================
%% ETHICS STATEMENT
%% ============================================================
\newpage
\thispagestyle{empty}

\chapter*{Ethics Statement}
\phantomsection
\addcontentsline{toc}{chapter}{Ethics Statement}

This project is designed to be functional on public blockchain testnets, with no intention for monetary profit, nor to create institutional leverage over decentralized financing. No personal banking data, or private financial data, or codes of existing projects, are used in the development of the prototype platform and training AI/ML data. The full project prototype build is to be demonstrated in the final defence of the project; the platform is expected to be ready for testing and troubleshooting by pre-thesis~2. All the development workflow can be traced by the official GitHub repository for the project, and the planning, coding, testing, and usage of different dependencies for this project is considered open-source with MIT licensing, and will be open for contribution after the submission of the final build.

For the training of the AI models, AI-assisted financial decision-making, risk assessment, privacy features, on-chain and off-chain data processing and data storing, we acknowledge the ethical considerations while making decisions for the development of the project.

%% ============================================================
%% ABSTRACT
%% ============================================================
\newpage
\thispagestyle{empty}

\chapter*{Abstract}
\phantomsection
\addcontentsline{toc}{chapter}{Abstract}

Modern Institutional Banking is established in layers to obtain better control of global financial flows to create intentional restrictions for equal access to financing capabilities and centralized system. To address the complex layered or tier-based architecture that the global market follows for better governance and balance in distribution of financing, and to demonstrate that a decentralized crypto currency system can be established as a more viable solution to global financing, more fair and equal distribution with advanced security system, \textit{Crypto World Bank} project is built with the intention to show a connecting bridge between Decentralized banking system for transparency in traditional institutional layer system for adaptability of Blockchain technology for the global market. The project reveals that it is possible to build a programmable lending platform with better transparency in global financial flow with better security implementation and how Artificial Intelligence can enhance the experience for better adaptability with the growing demand of AI assistance and security. The integration of a combined stack of blockchain technology into the layered-based banking system with AI enhanced security and assistance has not been recorded or attempted before and the purpose of the project is to show that it is possible to make this combined complex system sustainable and functional while making it adaptable with the growing global financial development demands.

\vspace{0.5em}
\noindent\textbf{Keywords:} Blockchain, Decentralized Finance, Deposit Mobilization, Explainable AI, Financial Sustainability, Group Lending, Hierarchical Lending, Institutional DeFi, Multi-Entity Co-Lending, Oracle-Mediated Risk Analytics, Reserve Management, Smart Contracts

%% ============================================================
%% DEDICATION
%% ============================================================
\newpage
\thispagestyle{empty}

\chapter*{Dedication}
\phantomsection
\addcontentsline{toc}{chapter}{Dedication}

\vspace{3cm}
\begin{center}
\textit{Dedicated to our families for their unwavering support throughout our academic journey.}
\end{center}

%% ============================================================
%% ACKNOWLEDGMENT
%% ============================================================
\newpage
\thispagestyle{empty}

\chapter*{Acknowledgment}
\phantomsection
\addcontentsline{toc}{chapter}{Acknowledgment}

We would like to express our sincere gratitude to our panel members for their guidance and support throughout this project. We also thank our supervisor, Mr.\ Annajiat Alim Rasel, Senior Lecturer at the Department of Computer Science and Engineering, BRAC University, for his guidance and support. Finally, we thank the Department of Computer Science and Engineering at BRAC University for providing us with the resources and academic environment to pursue this work.

%% ============================================================
%% TABLE OF CONTENTS, LIST OF TABLES, LIST OF FIGURES
%% ============================================================
\newpage
\singlespacing
\setcounter{page}{1}
\pagenumbering{roman}
\tableofcontents
\newpage
\phantomsection
\addcontentsline{toc}{chapter}{List of Tables}
\listoftables
\newpage
\phantomsection
\addcontentsline{toc}{chapter}{List of Figures}
\listoffigures
\newpage
\chapter*{List of Formulas}
\phantomsection
\addcontentsline{toc}{chapter}{List of Formulas}
\begin{description}[leftmargin=0pt]
    \item[Utilization ($U$):] \(U = \frac{L}{L + B}\)
    \item[Kinked rate ($U \leq U^*$):] \(r_b(U) = r_0 + \frac{U}{U^*} \cdot r_1\)
    \item[Kinked rate ($U > U^*$):] \(r_b(U) = r_0 + r_1 + \frac{U - U^*}{1 - U^*} \cdot r_2\)
    \item[Collateral Ratio:] \(CR = \frac{C}{D}\)
    \item[Loan-to-Value:] \(LTV = \frac{\text{Loan Amount}}{\text{Collateral Value}}\)
    \item[APR (simple):] \(APR = r_{\text{period}} \times N\)
    \item[Compound interest:] \(A = P\left(1 + \frac{r}{n}\right)^{nt}\)
    \item[EMI:] \(EMI = \frac{P \cdot r \cdot (1+r)^n}{(1+r)^n - 1}\)
    \item[Health Factor:] \(HF = \frac{C \times LT}{D}\)
    \item[Group Health Factor:] \(HF_{\text{group}} = \frac{\text{total\_group\_collateral} \times LT}{\text{total\_group\_outstanding\_debt}}\)
    \item[Institutional Health Factor:] \(HF_{\text{inst}} = \frac{\text{on\_chain\_reserve} - \text{min\_reserve}}{\text{outstanding\_borrowing}}\)
    \item[Reserve Ratio:] \(RR = \frac{\text{Reserves}}{\text{Total Deposits}}\)
    \item[Platform Solvency Ratio:] \(PSR = \frac{\text{TotalReserveBalance}}{\text{TotalOutstandingLoans}}\)
    \item[Credit Velocity:] \(CV = \frac{\text{Capital Recycled per Period}}{\text{Total Reserve}}\)
    \item[Interbank rate cap:] \(r_{\text{IB}}(U) \leq r_{\text{downward}}(U) - \delta\)
    \item[Upward deposit rate cap:] \(r_{\text{up}} < r_{\text{down}} - \delta\)
    \item[FX conversion:] \(\text{Amount}_B = \text{Amount}_A \times \frac{P_A}{P_B}\)
    \item[Treasury swap:] \(\text{Amount}_{\text{to}} = \text{Amount}_{\text{from}} \times \frac{P_{\text{from}}}{P_{\text{to}}} \times (1 - \text{spread})\)
    \item[Oracle commit--reveal:] \(h = \text{keccak256}(s \,\|\, \text{nonce})\)
    \item[Group loan share:] \(\text{Share}_i = \frac{\text{LoanTotal}}{N_{\text{members}}}\)
    \item[On-chain credit score:] \(CS = \sum_i w_i \cdot \text{feature}_i\)
    \item[Isolation Forest score:] \(s(x, n) = 2^{-E(h(x))/c(n)}\)
    \item[Random Forest fraud probability:] \(P(\text{fraud}\mid x) = \frac{1}{T}\sum_{t=1}^{T} f_t(x)\)
    \item[Net interest split:] \(\text{NetInterest} = \text{InterestCollected} - \text{DepositorYield} - \text{InsuranceAllocation}\)
    \item[SHAP (Shapley value):] \(\phi_i = \sum_{S \subseteq F \setminus \{i\}} \frac{|S|!(|F|-|S|-1)!}{|F|!}\left[f(S \cup \{i\}) - f(S)\right]\)
\end{description}
\newpage
\chapter*{List of Abbreviations}
\phantomsection
\addcontentsline{toc}{chapter}{List of Abbreviations}
\begin{description}[leftmargin=0pt]
    \item[AA:] Account Abstraction (ERC-4337)
    \item[ALM:] Asset-Liability Management
    \item[AI:] Artificial Intelligence
    \item[AI/ML:] Artificial Intelligence / Machine Learning
    \item[AML:] Anti-Money Laundering
    \item[APR:] Annual Percentage Rate
    \item[API:] Application Programming Interface
    \item[BRAC:] Bangladesh Rural Advancement Committee
    \item[CAR:] Capital Adequacy Ratio (Basel~III; see PSR for this platform)
    \item[CBDC:] Central Bank Digital Currency
    \item[CCIP:] Cross-Chain Interoperability Protocol (Chainlink)
    \item[CCS:] ACM Computing Classification System
    \item[CR:] Collateral Ratio
    \item[CVL:] Certora Verification Language
    \item[DeFi:] Decentralized Finance
    \item[DID:] Decentralized Identifier (W3C)
    \item[DON:] Decentralised Oracle Network (Chainlink)
    \item[EIP:] Ethereum Improvement Proposal
    \item[EIP-7702:] Ethereum Improvement Proposal 7702 --- Session Key standard enabling scoped, time-bound wallet authorisations for an autonomous agent
    \item[EMI:] Equated Monthly Installment
    \item[ERC:] Ethereum Request for Comments (token / interface standard)
    \item[EVM:] Ethereum Virtual Machine
    \item[FL:] Federated Learning
    \item[FATF:] Financial Action Task Force
    \item[FX:] Foreign Exchange
    \item[GNN:] Graph Neural Network
    \item[HF:] Health Factor
    \item[KYC:] Know Your Customer
    \item[LLM:] Large Language Model
    \item[L2:] Layer 2
    \item[LTV:] Loan-to-Value Ratio
    \item[MFI:] Microfinance Institution
    \item[MiCA:] Markets in Crypto-Assets Regulation (EU)
    \item[ML:] Machine Learning
    \item[MCP:] Model Context Protocol --- tool-calling interface used by the AI agent to interact with CWB banking operations
    \item[NIM:] Net Interest Margin
    \item[PRISMA:] Preferred Reporting Items for Systematic Reviews and Meta-Analyses
    \item[PoS:] Proof of Stake
    \item[PoR:] Proof of Reserve (Chainlink on-chain reserve verification)
    \item[PSR:] Platform Solvency Ratio
    \item[QRC:] QR code
    \item[RBAC:] Role-Based Access Control
    \item[RLS:] Row-Level Security (PostgreSQL policy feature enforcing append-only access on audit tables)
    \item[RR:] Reserve Ratio
    \item[SAR:] Suspicious Activity Report
    \item[SBT:] Soulbound Token (non-transferable ERC standard)
    \item[SHAP:] SHapley Additive exPlanations
    \item[SOFR:] Secured Overnight Financing Rate
    \item[SSE:] Server-Sent Events
    \item[SSI:] Self-Sovereign Identity
    \item[SoK:] Systematization of Knowledge
    \item[SWC:] Smart Contract Weakness Classification
    \item[TVL:] Total Value Locked
    \item[U:] Utilization
    \item[UUPS:] Universal Upgradeable Proxy Standard (ERC-1822)
    \item[VC:] Verifiable Credential (W3C)
    \item[XAI:] Explainable Artificial Intelligence
    \item[ZKP:] Zero-Knowledge Proof
    \item[zkAML:] Zero-Knowledge Anti-Money-Laundering proof
    \item[zkEVM:] Zero-Knowledge Ethereum Virtual Machine --- a ZK rollup that inherits Ethereum L1 security via cryptographic validity proofs
    \item[zk-SNARK:] Zero-Knowledge Succinct Non-Interactive Argument of Knowledge
\end{description}
\newpage
\pagenumbering{arabic}
\setcounter{page}{1}


%% CHAPTER 1: INTRODUCTION
\chapter{Introduction}
\label{sec:intro}

The \textit{Crypto World Bank} is a research-based project for exploring how blockchain technology can be shaped to fit into institutional financial systems with all the added benefits that come with a decentralized financing system. Based on our findings that had similar motivation to our study, existing DeFi protocols show isolated functionalities, such as Aave for lending, Uniswap for exchange, and MakerDAO for stablecoins. This project specifies a tier-based, hierarchically governed platform that includes a mixture of institutional financing with DeFi, accommodating deposit mobilization, credit allocation, interbank liquidity, installment repayment, FX facilitation, solidarity group lending, reserve enforcement, syndicated co-lending, and risk-based governance covering all four tiers that systematically mirrors real-world institutional financing. At the top, Tier~1 represents the World Bank reserve; the core banking hierarchy is specified in Solidity and being improved until final deployment. Another significant feature---the multi-entity contract suite---is fully specified in Chapter~\ref{ch:design-alternatives}, which is scheduled for practical development in Phases~II--III.

The system is designed for three categories of participant: a programmable World Bank reserve as Tier~1, national and local development banks serving in Tiers~2--3, and retail clients in Tier~4. Retail clients get access to savings, group loans, and installment lending features. This hierarchical structure in this project was motivated by unbanked and MSME financing gaps~[9, 13]. The project is not built with any motivation to deploy only in regions with access to advanced technologies; rather, proper distribution of access to finance and fairness is the leading cause of this attempt to find suitable ground to adapt DeFi among the general population. The whole project is open source and available in our \href{https://github.com/Yuvraajrahman/Cryto-World-Bank}{GitHub repository}. The added benefits of blockchain: it provides state visibility, tamper-evident audit trails, and programmable rule enforcement among institutions having competition, as demonstrated by the renowned World Bank FundsChain programme~[25] and JPMorgan Kinexys~[26]. From the study of institutions, open ledger infrastructure and institutional hierarchy are not mutually exclusive.




\begin{figure}[H]
\centering
\BalancedDiagram{Ch1_system-overview-four-tier-stack}
\caption[Crypto World Bank system overview]{Crypto World Bank system overview}
\label{fig:intro-system-overview}
\end{figure}

\section{Background}
\label{sec:background}

Local lenders, or our target users, for this aspect of lending through layered institutions---supranational bodies that is provided by development finance channels have structural costs. Idle capital is usually locked by correspondent banking in pre-funded nostro/vostro accounts~[15], cross-border transactions that take around two to five days cost roughly \$45 per transaction~[101], charges for remittance corridors average 6.49\%, which is above the UN SDG target of 3\%~[17]. Due to unequal access to traditional banking systems, 1.4~billion adults remain unbanked~[9] and MSMEs face a \$4.5~trillion annual financing gap~[13].

In decentralized financing, lending, collateral, and interest can be automated by enabling rules on-chain transparently. However, institutions like Aave~v3 that holds \$26.3~billion TVL~[18] and the broader DeFi lending market exceeds \$55~billion~[8], use flat pool design that ignore institutional hierarchy, inter-tier exchange, and accessing clients via local bank lend path which is more generalized for normal people with less knowledge about DeFi.

\paragraph{Contribution of smart-contract and blockchain.}
Smart-contracts can develop lending rules as immutable, atomic, with auditable on-chain logic. This technology replaces the bilateral ledger reconciliation with a system of shared record, when a local bank initiates a loan, the World Bank's reserve ratio updates on the same transaction status~[25]. DeFi extends this feature to financial services without licensed custodians. Here in The \textit{Crypto World Bank} adds three dimensional feature that is missing from the existing institutional protocols: institutional hierarchy (four-tier capital flows), multi-entity operations (solidarity groups, syndicated loans, tranched pools), and a connection to the evolving development finance framework identified as absent in the literature~[1].

\paragraph{Optional conversational interface.}
Every user interacts with this system primarily through the React UI and wallet credential signing. For development purposes and demonstration, a locally hosted agent (Qwen3-8B + MCP, 17 banking tools) provides an alternate path to get the essential tasks done, including all the customer services that require system interaction information guidance and assistance in acquiring a loan or exchange finances with the system. In the final build it is planned to be replaced with secured API-based services.

\paragraph{Technical foundations.}
\label{sec:blockchain-fundamentals}

\textbf{Consensus.} The chain is an append-only ledger: nodes agree on history, and no one can rewrite past blocks without controlling most of the network~[50].

\textbf{EVM execution~[104].} Loan logic runs as deterministic bytecode on Ethereum-compatible chains. Storage writes are the main cost driver (Section~\ref{sec:gas-cost}).

\textbf{State machines.} Each loan moves through clear stages---\texttt{PENDING} $\to$ \texttt{APPROVED} $\to$ \texttt{ACTIVE} $\to$ \texttt{REPAYING} $\to$ \texttt{COMPLETED/DEFAULTED}---with only authorised roles allowed to advance it.

\textbf{Oracles~[45, 96].} Contracts cannot read off-chain data on their own. ML risk scores reach the chain via a commit-reveal oracle (Section~\ref{sec:oracle-architecture}); Chainlink Feeds, Automation, and Proof of Reserve handle prices, due-date checks, and reserve attestation. Live ML wiring starts after pre-thesis~2.

\textbf{Polygon zkEVM.} We deploy on Cardona testnet (Section~\ref{sec:high-level-arch}) for low fees and L1-backed validity proofs.

\section{Rationale and Motivation}
\label{sec:rationale}

Real development finance depends on layered institutions, yet settlement friction, opaque reserves, and correspondent-banking idle capital leave both intermediaries and unbanked clients underserved~[17, 90]. Flat DeFi protocols automate lending at scale but ignore institutional hierarchy, cross-tier flows, and the Local-Bank-to-client path that microfinance and development banks use in practice~[1]. This study is motivated by the opportunity to encode that hierarchy on a programmable ledger---with enforced reserves, role-based governance, and auditable capital movement---while supporting (not replacing) human credit judgment through lightweight analytics and an optional conversational interface. The thesis specifies a reusable architecture; it does not claim a particular national deployment target.

\section{Problem Statement}

Three primary forces of motivation for this study: development-finance has restricting transparency and settlement; the absence of multi-tiered banking functions in DeFi (deposits, solidarity and lending); the keen interest in combining blockchain auditability with ML computational analytics. There are six structural inefficiencies noticeable in development finance routes for individual borrowers from national and local banks:

\begin{enumerate}[noitemsep]
\item \textbf{Lack of transparency:} there is very little chance of the lower tier users to verify capital management, reserve levels or lending decisions are made by the higher tiers, very less transparency on how money flows through the different levels of systems. Only the top tiers get the information and records of reserves and audited at most quarterly.
\item \textbf{Sluggish settlements and idle capital:} cross-border transactions are quite expensive, average \$45 per transaction~[101], and time consuming taking two to five days~[23].
\item \textbf{Inconsistent risk evaluation:} fraud and credit assessment are completely relied on human judgement, borrowers are generally re-evaluated at each tier or institutions based on information available at that tier, no unified credit history.
\item \textbf{Access and trust barriers:} the cost of inter-institutional trust is quite expensive for legal and audit process. It creates congestions in developing countries where investors earn returns roughly 3\% below compared to developed markets~[103].
\item \textbf{No programmable savings:} depositors in developing economies have no real-time access to visibility into how savings are processed and deployed.
\item \textbf{No on-chain group credit:} without individual collateral, borrowers cannot access programmable on-chain mutual group lending with shared responsibility. Real-world organizations like BRAC exist but does not function like DeFi.
\end{enumerate}

DeFi enabled automatic lending transparency at scale such as Aave~v3 with \$26.3~billion TVL, \$17.7~billion borrows~[18], Compound~v3 (\$1.4~billion)~[19], MakerDAO with \$10.5~billion stablecoin issuance~[20]. Despite large market coverage they exhibit four limitations for institutional development finance:

\begin{itemize}[noitemsep]
\item \textbf{Flat lending models.} Aave/Compound uses one protocol for many asset pools (USDC, ETH, etc.). There is no cross-tier allocation, no institutional hierarchy, no difference in role-based governance. Here our CWB models institutional relationships that go deep into what kind of services to provide to what kind of roles/institutions to increase liquidity in overall lending and exchange situations.
\item \textbf{No AI risk analytics.} DeFi relies on collateral ratio, utilization curves. There is no behavioral fraud detection or explainable credit support system.
\item \textbf{No tiered governance.} DeFi protocol treats tokens as a measurement of voting, not based on what is the role of the client or user or institution in the bigger to smaller picture as it lacks relation to a deeper level like institutions.
\item \textbf{Single-level institutional DeFi.} Maple \& Goldfinch~[21] with \$2.6--3.8~billion TVL and \$680~million originations in 18+ countries respectively serve institutions, but none models a full multi-tier banking system with interbank flows, lack hierarchical capital flows.
\end{itemize}

\section{Objectives}

The objectives of this project are:

\begin{enumerate}[noitemsep]
\item Design and specify all features and constraints of a four-tier lending architecture on an EVM-compatible blockchain and complete a full testnet deployment (World Bank $\to$ National Bank $\to$ Local Bank $\to$ Client). The whole development cycle is planned for Phases~I--IV.
\item Design not only pool lending, but specify an extended banking product suite on the hierarchical system (deposits, savings, and solidarity group lending).
\item Investigate a lightweight off-chain analytics layer using ML, training and testing the models in practical testnet for fraud detection, anomaly detection, explainable review, and an optional MCP-based conversational interactive agent for ease of use.
\item Specify programmable on-chain lending with borrowing limits based on financial behaviour analysis, installments, and role-based access control and information via smart contracts.
\item Plan and execute testnet deployment using Polygon zkEVM Cardona, Ethereum Sepolia, and evaluate hierarchical controls in the final thesis phase.
\item Document governance for future improvements, security analysis and feature improvements, role-based regulatory considerations, and a controlled rollout toward sandbox piloting.
\end{enumerate}

\section{Proposed Solution}
\label{sec:proposed-solution}

The Crypto World Bank will implement a four-tier on-chain lending model: Tier~1 World Bank reserve custodian $\to$ Tier~2 National Banks $\to$ Tier~3 Local Banks processing retail loans $\to$ Tier~4 clients building on-chain credit history (Figure~\ref{fig:intro-system-overview}), with six standard banking functions enforced on-chain---deposit mobilization, credit allocation, payment and settlement, risk intermediation, liquidity management, and ancillary services---rather than discretionary audit alone.
\label{sec:banking-functions}
Banks also borrow from peers, park surplus with parent tiers, and syndicate large loans; the platform replicates four-tier downward flow, same-tier pools, upward deposits, and multi-bank co-funding on-chain (Figure~\ref{fig:cross-tier-lending}).
\label{sec:cross-tier-lending}

\begin{figure}[H]
\centering
\BalancedDiagram{Ch1_cross-tier-lending-flows}
\caption[Cross-tier lending flow directions]{Cross-tier lending flow directions}
\label{fig:cross-tier-lending}
\end{figure}



\paragraph{Can a group of entities lend or fund together?}\label{par:group-entities}
Syndicated loans, solidarity groups, tranched pools, and upward surplus deposits (Figure~\ref{fig:multi-entity-cofunding}; Section~\ref{sec:multi-entity}) coexist with same-tier peer liquidity pools (Section~\ref{sec:interbank-pool}); retail clients enter through Local Banks via KYC, Credit Passport, and borrowing limits, with syndication reserved for large loans.




\begin{figure}[H]
\centering
\BalancedDiagram{Ch3_multi-entity-cross-tier-operations}
\caption[Multi-entity co-funding mechanisms]{Multi-entity co-funding mechanisms: syndicated loan, solidarity group, tranched pool, and upward deposit}
\label{fig:multi-entity-cofunding}
\end{figure}

\section{Methodology in Brief}

There are four Agile phases to this methodology for this project. This report is designed with the architecture specification. The development phase will start from Phase~II after pre-thesis~1. In this report there is architecture, contracts, database, and ML pipeline. The Phase~II development phase will demonstrate the loan cycle. Phase~III will start the implementation of ML scoring and the Chainlink oracle (Section~\ref{sec:oracle-architecture}). Phase~IV will show the evaluation, metrics, and the overall feasibility of the project in the real world. Stack: Solidity, React, Express/FastAPI, scikit-learn on Polygon zkEVM Cardona. There is also an AI chat agent that will be built as an optional feature that will run alongside the project to help the clients navigate better with the system.

\section{Scopes and Challenges}
\label{sec:scopes-challenges}

\textbf{In scope now.} In our development phases, we have calculated that we can implement a limited amount of our full vision: the four-tier lending architecture, loan lifecycle workflow, borrowing limits for the clients and the banks, the design of Oracle fraud analytics, and the training and the live demonstration of the AI/ML workflow. It can be shown to a limited amount on the prototype builds only.

\textbf{Out of scope for now.} Mainnet deployment, live retail payments in the real world, production KYC/AML at the banking level and client interaction, stress testing the whole project, and providing an MVT checklist before the product is ready is not quite possible yet.

\textbf{Main risks.} There are a couple of risks that we have noticed for our project, starting with the fraud labels may not be completely accurate based on the limited amount of data we have for our system architecture. We need more synthetic data that will be collected over time. If regulations have changed---mitigated by testnet-only builds---then we might have to change the architectural elements. The gas toll on micro-loans may not be sustainable. The ML layer with SHAP may keep borrowers in the loop in the system since it is not 100\% accurate and accuracy of the system depends on the training process. Rest of the economics and trade-offs are discussed in Section~\ref{sec:economic-analysis}.

\paragraph{Economic sustainability.} On Cardona, 10--15 txs on a typical loan cost under \$0.15---small next to interest earned. Mainnet micro-loans would need tighter gas work.

\paragraph{Institutional capture.} Another huge issue that needs to be addressed is that Tier~1 and the National Banks, which have huge reserves and influence over the rates, reserves, and who gets credit, might influence the overall system and cost-to-performance ratio. It might temper the trust and relationship with the retail customers and smaller banks.

\paragraph{Stablecoin-first lending.}
\label{sec:stablecoin-first}
Retail clients are not to be carrying high-priced positions like ETH, since the prices of Ethereum might fluctuate. MiCA and the U.S. GENIUS Act support the idea that stablecoins should be more generalized for authorized cryptocurrency. There are many settlement projects like \textit{mBridge} and \textit{Agora}~[61, 62], moving the value of cryptocurrencies across borders, but it is not a replacement for our lending system. The architectural system for stablecoin pools and L2 deployment (Section~\ref{sec:bridge}) gives the liquidity move fluidly from the top tier to the bottom, and the surplus to go up from the bottom.

\section{Thesis Organization}
\label{sec:thesis-organization}

\begin{itemize}[noitemsep]
\item \textbf{Ch.~1 --- Introduction:} problem, objectives, approach (\textbf{CO1}).
\item \textbf{Ch.~2 --- Literature:} PRISMA review and protocol comparison (\textbf{CO9}).
\item \textbf{Ch.~3 --- Specifications and impact:} requirements, societal/environmental impact, ethics, project management, economics (\textbf{CO2--4}, \textbf{CO10--12}).
\item \textbf{Ch.~4 --- Design alternatives:} design process, preliminary architecture, technology choices (\textbf{CO5}).
\item \textbf{Ch.~5 --- Evaluation and analysis:} performance evaluation, design analysis, discussion (\textbf{CO6--7}).
\item \textbf{Ch.~6 --- Conclusion:} summary, contributions, and future work; appendices hold contracts and deployment notes.
\end{itemize}

%% CHAPTER 2: LITERATURE REVIEW
\chapter{Literature Review}

\section{Preliminaries}
\label{sec:preliminaries}

Important terms---XAI, PoS, CBDC, solidarity/group lending, DeFi, smart contracts, over-collateralization, ALM, flash loans, ZKP, and Health Factor---have been introduced in Glossary appendix. This section has been expanded in appendix~D that will be presented in the final versions of the complete project file.

\subsection{Review Methodology (PRISMA-Style Frame)}
\label{sec:lit-prisma}

We did our research to find the appropriate papers and highly qualified papers that goes without project from ACM, IEEE, Elsevier, MDPI, arXiv, and BIS/World Bank (2017--2026) with PRISMA-style screening~[52]. Primary topics for this search: DeFi lending, fraud ML, tiered finance, and smart-contract security.



\section{Review of Existing Research}
\label{sec:review-existing}

We explain our literature review, and our findings from the papers and the work of the authors, in the subsections shown below.

\subsection{Decentralized Lending and DeFi Protocol Design}

For a big picture, Werner et al.~[1] show that DeFi lenders use flat pools with no bank-style tiers; we have found a gap that can be filled by introducing our project. According to Bastankhah et al.~[2], adaptive borrow rates work better than Aave-style curves for financial flows, which we introduced in our kinked-rate module. Dao et al.~[39] show how on-chain limits can get off-chain credit scores that in general are applied at our Local Bank level.

\subsection{Machine Learning, Explainable AI, and Extended AI Security}

For fraud detection, Palaiokrassas et al.~[3] engineer behavioral features that do a really good job in DeFi transaction fraud detection, not just hashes. We plan to use Random Forest in Phase~III that will further strengthen the fraud detection level with the mixture of hash functions that is built into the DeFi protocol. Moreover, for the situations when the fraud detection level is scarce, we will use Liu et al.~[6] Isolation Forest. We also studied a very in-depth comparison of SHAP versus LIME by Adom et al.~[4], which shows a significant advantage by using SHAP for loan detection; approvers in bank and authority level will see SHAP screening results before signing.

\subsection{Hierarchical Distribution, Cross-Tier Flows, and Group Lending}
\label{sec:lit-cross-tier}

Tan~[5] demonstrated how CBDC flows in bank-to-bank and bank-to-client financial transactions. We extended this idea in our project by showing tier-based reserve checks. For layered reserve contracts, \textit{Project Pine}~[107] has shown a prototype that is similar to our \texttt{InterBankLendingPool} and the syndicated-style loan flow. We also integrated the style of microcredit admin cost with smart contracts based on the study of Asaduzzaman et al.~[36]. The structure of \texttt{GroupLendingPool} function is the on-chain version, without claiming that we are trying to replace BRAC-style social microfinancing enforcement.

\subsection{Smart Contract Security and Layer~2 Scalability}

To improve the security system of the project, we are following Atzei et al.~[7], which lists the main Ethereum contract attacks. Since the tier-based system is very complex, many transactions on Cardona might be triggered. Gas prices can be cut to 30--50\% by transaction batching~[102], with further DeFi gas optimisation~[37].

\subsection{Monetary Policy, Tokenization, and Institutional Landscape}
\label{sec:lit-institutional}

Very large institutions like World Bank FundsChain~[25] and similar institutions show that development finance is already moving on-chain for auditability; we are further improving this by encoding rules, not just managing disbursement logs.

\subsection{Graph ML, Federated Learning, Regulation, and Identity Primitives}

We noticed the usability and rise of stablecoin and went further deep into it and try to explore how stablecoin pools for retail clients can be strengthened by using MiCA and the U.S.\ GENIUS Act~[58, 60]. For our project, soulbound tokens will be used that fit into our Credit Passport idea, which was extended by the idea of Buterin et al.~[94]. This introduces a secured identity on-chain, a portable reputation, without the fear of exposing personal data.

\subsection{Oracle-Mediated ML and Optional LLM Interfaces}
\label{sec:lit-llm}

Lastly, for an optional feature to our project, future studies can be made and features can be added to the system. Caldarelli~[96] describes how ML scores still have some issues that go beyond price oracles; we focus on this part and extend our commit-reveal or Chainlink path targets. We collected some reputable databases, and more to be added such as the BCCC fraud dataset~[99] that goes with our system. This section will be continued and extended after pre-thesis~2.

\begin{table}[H]
\centering
\footnotesize
\setlength{\tabcolsep}{3pt}
\renewcommand{\arraystretch}{1.08}
\caption[Literature review summary (Part 1 of 2)]{Literature review summary (Part 1 of 2)}
\label{tab:lit-summary}
\resizebox{\textwidth}{!}{%
\begin{tabular}{@{}>{\raggedright\arraybackslash}p{3.0cm}>{\raggedright\arraybackslash}p{2.6cm}>{\raggedright\arraybackslash}p{6.8cm}@{}}
\toprule
\tableheadcolor
\tableheadfont Author \& Focus & \tableheadfont Methodology & \tableheadfont What they found \& how we use it \\
\midrule
Palaiokrassas et al.\ (2024) [3] · DeFi fraud & XGBoost, NN on 54M+ txs & We plan to use behavioural features in Phase~II for Random Forest, as it lifted F1 to 0.76--0.85 \\
Adom et al.\ (2025) [4] · XAI lending & LIME / SHAP comparison & Approvers get a SHAP brief before signing, as it gives a clearer result than LIME on loan data \\
Tan (2023) [5] · CBDC hierarchy & Developing-nation model & CBDC backs the four-tier layout and works as a bank-to-user tier \\
Liu et al.\ (2008) [6] · Anomaly detection & Isolation Forest & Flags odd wallets without many labels; backup when fraud data is thin \\
Atzei et al.\ (2017) [7] · SC security & SoK of Ethereum attacks & Lists re-entrancy and access-control risks; guides OpenZeppelin guards \\
Werner et al.\ (2022) [1] · DeFi SoK & Survey of 12+ protocol types & Main gap CWB targets is DeFi lending having no bank hierarchy \\
Bastankhah et al.\ (2024) [2] · Adaptive lending & Dual fast/slow control & Dynamic rates support our kinked-rate modules as it beats static curves \\
Hartmann \& Hasan (2021) [33] · P2P lending & Smart-contract origination & Gas is cut in on-chain P2P, which is a minor input for loan-facing clients \\
Hasan et al.\ (2022) [34] · Federated fraud & FL on encrypted features & Lack of raw-data sharing features by banks leaves future cross-bank options \\
Wang et al.\ (2020) [105] · SC vulnerabilities & ML on bytecode & Common bytecode bugs spotting by ML helps provide extra security over manual audit \\
Liao et al.\ (2019) [106] · Audit automation & Classification + fuzzing & Workflow idea for Phase~V is inspired from ML prioritizing fuzzing speed \\
\bottomrule
\end{tabular}%
}
\end{table}

\begin{table}[H]
\centering
\footnotesize
\setlength{\tabcolsep}{3pt}
\renewcommand{\arraystretch}{1.08}
\caption[Literature review summary (Part 2 of 2)]{Literature review summary (Part 2 of 2)}
\label{tab:lit-summary-b}
\resizebox{\textwidth}{!}{%
\begin{tabular}{@{}>{\raggedright\arraybackslash}p{3.0cm}>{\raggedright\arraybackslash}p{2.6cm}>{\raggedright\arraybackslash}p{6.8cm}@{}}
\toprule
\tableheadcolor
\tableheadfont Author \& Focus & \tableheadfont Methodology & \tableheadfont What they found \& how we use it \\
\midrule
BIS et al.\ (2020) [35] · CBDC design & Retail distribution analysis & Standard CBDC two-tier is extended to four-tier with on-chain reserves \\
Kiff et al.\ (2020) [80] · CBDC models & IMF architecture survey & Matching National/Local roles; retail CBDC flow implemented through supervised banks \\
Bindseil (2020) [79] · CBDC rates & ECB tiered remuneration & Upward deposit spread idea inspired from tiered rates that protect bank intermediation \\
Dong et al.\ (2023) [78] · Supply-chain finance & Three-tier game model & Capital can move down a chain, parallel to World Bank $\to$ client allocation \\
Project Pine (2025) [107] · Hierarchical SCs & NY Fed / BIS prototype & IBLP and syndicated loans used on-chain layered reserve contracts prototype \\
CPMI/FSB [15, 23] · Cross-border settlement & Nostro/vostro cost analysis & Tiered on-chain flow is better than legacy settlement as it is slow and costly \\
World Bank (2025) [25] · Dev.\ finance blockchain & 13-project pilot & Programmable lending rules added with dev-banks system \\
Chen et al.\ (2024) [31] · Multi-chain lending & Cross-chain design & Scaling by splitting loads throughout chain is later option, not for Phase~I \\
Yang \& Cui (2023) [32] · Protocol evaluation & Quantitative risk metrics & Gives utilisation and liquidation benchmarks for health checks tier \\
Asaduzzaman et al.\ (2020) [36] · Smart micro-credit & Smart-contract MFI & Model for \texttt{GroupLendingPool} is smart contracts, which cuts the admin costs in Bangladesh \\
Wang et al.\ (2021) [102] · Tx batching & iBatch middleware & Per call saved 15--59\% in gas price by which we deploy on Cardona L2 \\
Kim \& Kim (2024) [37] · Gas optimisation & DeFi parameter tuning & $\sim$18\% gas cut on Ethereum pools helps our L2 deployment \\
Yuan (2024) [38] · SHAP credit default & RF, GBM + SHAP & Deploying on Cardona L2 as batching saved 30--50\% gas; SHAP is also used for our decline path \\
\bottomrule
\end{tabular}%
}
\end{table}

\begin{table}[H]
\centering
\footnotesize
\setlength{\tabcolsep}{3pt}
\resizebox{\textwidth}{!}{%
\begin{tabular}{@{}p{3.4cm}*{6}{>{\centering\arraybackslash}p{1.35cm}}@{}}
\toprule
\tableheadcolor
\tableheadfont Feature & \tableheadfont Aave & \tableheadfont Comp. & \tableheadfont Maker & \tableheadfont Maple & \tableheadfont Gfinch & \tableheadfont CWB \\
\midrule
Institutional hierarchy & \tmarkNo{} & \tmarkNo{} & \tmarkNo{} & \tmarkNo{} & \tmarkNo{} & \tmarkPlanned{} \\
\rowcolor{RowShade}
Cross-tier capital flow & \tmarkNo{} & \tmarkNo{} & \tmarkNo{} & \tmarkNo{} & \tmarkNo{} & \tmarkPartial{} \\
Same-tier interbank lending & \tmarkNo{} & \tmarkNo{} & \tmarkNo{} & \tmarkNo{} & \tmarkNo{} & \tmarkPartial{} \\
\rowcolor{RowShade}
Syndicated / co-lending & \tmarkNo{} & \tmarkNo{} & \tmarkNo{} & Part. & Part. & \tmarkPartial{} \\
Upward capital repatriation & \tmarkNo{} & \tmarkNo{} & \tmarkNo{} & \tmarkNo{} & \tmarkNo{} & \tmarkPartial{} \\
\rowcolor{RowShade}
Solidarity group lending & \tmarkNo{} & \tmarkNo{} & \tmarkNo{} & \tmarkNo{} & Part. & \tmarkPlanned{} \\
AI/ML fraud detection & \tmarkNo{} & \tmarkNo{} & \tmarkNo{} & Man. & Man. & \tmarkPlanned{} \\
\rowcolor{RowShade}
SHAP explainability & \tmarkNo{} & \tmarkNo{} & \tmarkNo{} & \tmarkNo{} & \tmarkNo{} & \tmarkPlanned{} \\
ZKP KYC compliance & \tmarkNo{} & \tmarkNo{} & \tmarkNo{} & \tmarkNo{} & \tmarkNo{} & \tmarkPlanned{} \\
\rowcolor{RowShade}
Kinked interest rate curve & \tmarkDone{} & \tmarkDone{} & Gov. & \tmarkDone{} & \tmarkNo{} & \tmarkPartial{} \\
Role-based access control & Part. & Part. & Gov. & \tmarkDone{} & Part. & \tmarkPlanned{} \\
\rowcolor{RowShade}
Institutional / L2 focus & \tmarkNo{} & \tmarkNo{} & \tmarkNo{} & \tmarkDone{} & Part. & \tmarkPlanned{} \\
TVL / Status (2026) & \$26B & \$1.4B & \$10B & \$2.6B & \$680M & Planning \\
\bottomrule
\end{tabular}%
}
\caption[Protocol comparison on eleven architectural dimensions]{Protocol comparison on eleven architectural dimensions}
\label{tab:protocol-comparison}
\label{sec:comparative-protocol}
\end{table}

\section{Summary of Key Findings}
\label{sec:lit-key-findings}

Four points stand out after reviewing and Table~\ref{tab:protocol-comparison}:

\begin{itemize}[noitemsep]
\item \textbf{Flat DeFi vs.\ tiers.} Large pool lenders~[1, 8] have no hierarchical bank currently, and it opens a door to explore how a hierarchical system would behave with DeFi. CBDC and Project Pine work~[5, 107] prove that nobody combines a four reserve-enforced tier with group lending and cross-tier financial flows and an ML oracle in a singular architecture.
\item \textbf{ML and security.} ML and security layers combined with behavioural fraud features and SHAP~[3, 4, 6] further extend the usability, feasibility, and adaptability of our project. Based on these literature reviews, we further improve our Phase~III pipeline; contract attack catalogues~[7] and L2 gas batching~[102, 37] justify OpenZeppelin guards to be added to our project.
\item \textbf{Groups and institutions.} We explored smart-contract microcredit in Bangladesh~[36] and World Bank on-chain pilots~[25] that gave us the idea and improved the implementation of \texttt{GroupLendingPool} function and institutional lending system.
\item \textbf{Regulation.} For regulation improvement and securities, MiCA and GENIUS~[58, 60] are accepted by stablecoin-first pools. Soulbound credit tokens~[94] fit the Credit Passport without putting KYC on-chain.
\end{itemize}

%% CHAPTER 3: SPECIFICATIONS, IMPACT, ETHICS, AND PROJECT MANAGEMENT
\chapter{Specifications, Impact, Ethics, and Project Management}
\label{ch:specifications-impact}

\section{Final Specifications and Requirements}
\label{sec:specifications-requirements}
The functional and non-functional requirements are shown in Table~\ref{tab:functional-requirements} and Table~\ref{tab:nonfunctional-requirements}; based on them, the software engineering with in-depth architecture and system analysis is specified in Chapter~\ref{ch:design-alternatives}. The implementation plan has been specified in Section~\ref{sec:impl-phase-plan}. The MVT checklist has been specified in Table~\ref{tab:mvt-checklist}, and the complete features list is in Table~\ref{tab:feature-list}.

\begin{table}[H]
\centering
\caption[Functional requirements (representative)]{Functional requirements (representative subset). Full feature inventory in Table~\ref{tab:feature-list}.}
\label{tab:functional-requirements}
\small
\begin{tabular}{@{}p{1.2cm}p{4.2cm}p{8.8cm}@{}}
\toprule
\tableheadcolor
\textbf{ID} & \textbf{Requirement} & \textbf{Description} \\
\midrule
FR-01 & Four-tier hierarchy & World, National, and Local Bank contracts with role-gated downward capital allocation and client onboarding at Local Bank tier. \\
FR-02 & Loan lifecycle & Request $\to$ ML score commit--reveal $\to$ approver decision $\to$ disbursement $\to$ installment schedule $\to$ repayment. \\
FR-03 & Reserve enforcement & Each tier retains a minimum reserve ratio before allocating capital downward; on-chain revert on breach. \\
FR-04 & Borrowing limits & Rolling six-month and one-year limits enforced on-chain; off-chain projection in \texttt{BORROWING\_LIMIT}. \\
FR-05 & Risk scoring & Off-chain RF + IF + stacking; SHAP explainability; immutable \texttt{LOAN\_RISK\_ASSESSMENT} audit row per inference. \\
FR-06 & Group lending & \texttt{GroupLendingPool} with mutual liability, consent, and per-member share disbursement (Section~\ref{sec:banking-products}). \\
FR-07 & Event synchronisation & Blockchain event listener projects on-chain state into PostgreSQL for dashboards and agent context. \\
FR-08 & Dual client interface & React dashboard and MCP conversational agent call the same Express.js API and contracts. \\
\bottomrule
\end{tabular}
\end{table}

\begin{table}[H]
\centering
\caption[Non-functional requirements]{Non-functional requirements.}
\label{tab:nonfunctional-requirements}
\small
\begin{tabular}{@{}p{1.2cm}p{2.4cm}p{10.6cm}@{}}
\toprule
\tableheadcolor
\textbf{ID} & \textbf{Category} & \textbf{Constraint / target} \\
\midrule
NFR-01 & Performance & ML inference $\leq$50\,ms on laptop hardware; API p95 $\leq$500\,ms for read endpoints at MVT scale. \\
NFR-02 & Security & RBAC on contracts and API; \texttt{ReentrancyGuard}; human approver gate before disbursement; agent write-path confirmation. \\
NFR-03 & Scalability & Modular contract deployment; PostgreSQL indexing (Appendix~\ref{app:db-schema-reference}); L2/testnet gas target sub-cent per retail lifecycle step. \\
NFR-04 & Maintainability & TypeScript/React/Express monorepo at repository root; OpenZeppelin primitives; documented phase gates G1--G4. \\
NFR-05 & Usability & Mobile-first React UI; five-stage onboarding funnel; Authority Brief for approvers (Section~\ref{sec:ui-design}). \\
NFR-06 & Compliance (design) & Tiered KYC ladder; GDPR data minimisation (hash-only on-chain); EU AI Act Art.~86 explainability artifacts. \\
NFR-07 & Reliability & Pause mechanism on contracts; zero-cost demo substitutes with documented upgrade path (Section~\ref{sec:design-limitations}). \\
NFR-08 & Interoperability & EVM-standard contracts; WalletConnect; ERC-20/ERC-4626 patterns; MCP tool schema for agent integration. \\
\bottomrule
\end{tabular}
\end{table}

\subsection{Minimum Viable Thesis (MVT) Checklist}
\label{sec:mvt-checklist}

\begin{table}[H]
\centering
\caption[MVT deliverable checklist]{MVT deliverable checklist.}
\label{tab:contributions-vs-future}
\label{tab:mvt-checklist}
\small
\begin{tabular}{p{0.5cm}p{9.2cm}p{1.8cm}p{1.8cm}}
\toprule
\tableheadcolor
\textbf{\#} & \textbf{Deliverable} & \textbf{Priority} & \textbf{Phase} \\
\midrule
1 & Four-tier bank contracts on testnet & Must & I \\
2 & Loan lifecycle E2E (request, approve, disburse, repay) & Must & II \\
3 & On-chain reserve and borrowing limits & Must & II \\
4 & ML risk scoring (RF + IF + stacking) & Must & III \\
5 & SHAP explainability and Authority Brief UI & Must & III \\
6 & Risk assessment audit row per inference & Must & III \\
7 & Commit--reveal score gate before approval & Must & III \\
8 & BCCC fraud-detection benchmark & Must & III \\
9 & 300-client simulation or gas-cost table & Must & IV \\
10 & README with reproduction instructions & Must & IV \\
11 & MCP agent (3--5 tools + confirmation gate) & Must & III \\
12 & LLM evaluation protocol & Should & IV \\
13 & Contract audit reports (Slither + Mythril) & Must & IV \\
14 & Foundry fuzz and invariant tests & Must & IV \\
15 & Safe multisig for World Bank admin & Must & I \\
16 & Tenderly monitoring alerts & Should & IV \\
\bottomrule
\end{tabular}
\end{table}



\subsection{List of Features}
\label{sec:feature-list}

Table~\ref{tab:feature-list} lists the complete list of features of the full project. The Must features will be completed during the development phases, Should and Stretch features are extensions of the main Must features, and Specified features are documented but not yet implemented, Future Work features are to show the potential improvements of the project and some of these features might be deliverable by the end of the project as special features.

\begingroup
\footnotesize
\setlength{\tabcolsep}{3.5pt}
\renewcommand{\arraystretch}{0.88}
\begin{longtable}{@{}>{\centering\arraybackslash}p{0.5cm}>{\raggedright\arraybackslash}p{9.0cm}>{\raggedright\arraybackslash}p{2.2cm}>{\centering\arraybackslash}p{1.6cm}@{}}
\caption[Complete feature inventory]{Complete feature inventory by category, priority, and implementation phase.}
\label{tab:feature-list} \\
\toprule
\tableheadcolor
\tableheadfont \# & \tableheadfont Feature & \tableheadfont Priority & \tableheadfont Phase \\
\midrule
\endfirsthead
\multicolumn{4}{r}{\small\itshape Table~\thetable{} (continued)} \\[2pt]
\toprule
\tableheadcolor
\tableheadfont \# & \tableheadfont Feature & \tableheadfont Priority & \tableheadfont Phase \\
\midrule
\endhead
\midrule
\multicolumn{4}{r}{\small\itshape Table~\thetable{} (continued on next page)} \\
\endfoot
\bottomrule
\endlastfoot
\multicolumn{4}{@{}l@{}}{\textit{Core platform and hierarchy}} \\
\rowcolor{RowShade}
1 & Four-tier banking hierarchy (World $\to$ National $\to$ Local $\to$ Client) & Must & I \\
2 & Role-based access control (RBAC) on contracts and API & Must & I \\
\rowcolor{RowShade}
3 & World Bank Reserve --- global capital custody & Must & I \\
4 & National Bank --- jurisdiction-level allocation & Must & I \\
\rowcolor{RowShade}
5 & Local Bank --- retail operations hub & Must & I \\
6 & Downward capital allocation (\texttt{allocateCapital}) & Must & I \\
\rowcolor{RowShade}
7 & Hierarchical bank registration (National $\to$ Local) & Must & II \\
8 & Reserve-ratio enforcement before downward allocation & Must & II \\
\rowcolor{RowShade}
9 & MockUSDC / testnet stablecoin for lending & Must & I \\
10 & 2-of-3 Safe multisig for \texttt{WorldBankAdmin} & Must & I \\
\rowcolor{RowShade}
11 & Contract pause / unpause (emergency controls) & Must & I \\
12 & Single-chain deployment (Sepolia / Amoy) & Must & I--IV \\
\midrule
\multicolumn{4}{@{}l@{}}{\textit{Retail lending lifecycle}} \\
\rowcolor{RowShade}
13 & Loan request submission (on-chain + UI) & Must & II \\
14 & Approver review workflow (approve / reject / request info) & Must & II \\
\rowcolor{RowShade}
15 & Loan disbursement to client wallet & Must & II \\
16 & Installment schedule generation & Must & II \\
\rowcolor{RowShade}
17 & Installment payment processing & Must & II \\
18 & On-chain loan state machine (\texttt{PENDING\_SCORE} $\to$ \texttt{ACTIVE}) & Must & II \\
\rowcolor{RowShade}
19 & Borrowing limits (six-month / one-year rolling caps) & Must & II \\
20 & Collateral-based lending (LTV) & Must & II \\
\rowcolor{RowShade}
21 & Credit-based lending (no collateral, SBT-gated) & Should & II \\
22 & Kinked interest rate model (utilization-based) & Should & II \\
\rowcolor{RowShade}
23 & Credit tier schedule (Bronze $\to$ Diamond) & Should & II \\
24 & Overdue installment checks (cron / Chainlink Automation) & Should & II \\
\rowcolor{RowShade}
25 & Loan history and transaction tracking UI & Should & II \\
26 & Income verification document upload and review & Should & II \\
\midrule
\multicolumn{4}{@{}l@{}}{\textit{Deposits and savings products}} \\
\rowcolor{RowShade}
27 & SavingsVault --- variable-yield savings (ERC-4626) & Should & II \\
28 & FixedDeposit --- term deposits (30/90/180/365 days) & Specified & II+ \\
\rowcolor{RowShade}
29 & Current / checking account (atomic transfers) & Specified & II+ \\
30 & Deposit-to-lending pool integration & Should & II \\
\rowcolor{RowShade}
31 & Interest split (depositor / InsuranceFund / protocol) & Should & II \\
32 & Withdrawal gated by reserve ratio & Must & II \\
\rowcolor{RowShade}
33 & InsuranceFund (5\% interest capture, default coverage) & Should & II \\
\midrule
\multicolumn{4}{@{}l@{}}{\textit{Group and solidarity lending}} \\
\rowcolor{RowShade}
34 & Group formation (3--20 members) & Must (C2) & II \\
35 & Multi-signature group consent & Must (C2) & II \\
\rowcolor{RowShade}
36 & Per-member disbursement share & Must (C2) & II \\
37 & Shared collateral / mutual liability & Must (C2) & II \\
\rowcolor{RowShade}
38 & Over-indebtedness controls (DTI, cooling-off, max two group loans) & Should & II \\
39 & Group credit history on SBT & Should & II \\
\midrule
\multicolumn{4}{@{}l@{}}{\textit{Multi-entity and cross-tier operations}} \\
\rowcolor{RowShade}
40 & InterBankLendingPool --- same-tier peer lending (1/7/30 day) & Should & II \\
41 & UpwardDepositFacility --- surplus repatriation to parent tier & Should & II \\
\rowcolor{RowShade}
42 & SyndicatedLoan --- multi-bank co-funding & Stretch & III \\
43 & TranchedPool --- senior/junior tranches & Stretch & III \\
\rowcolor{RowShade}
44 & TreasurySwap --- cross-tier treasury FX & Should & III \\
45 & FXModule --- oracle-priced retail FX & Should & III \\
\rowcolor{RowShade}
46 & Dual-currency facility (fiat $\leftrightarrow$ crypto) & Specified & III \\
47 & Cascading repayment / interest across tiers & Should & II \\
\rowcolor{RowShade}
48 & NettingEngine --- multilateral settlement & Future Work & --- \\
49 & Trade finance / escrow letters of credit & Future Work & --- \\
\midrule
\multicolumn{4}{@{}l@{}}{\textit{Credit passport and identity}} \\
\rowcolor{RowShade}
50 & On-chain Credit Passport (Soulbound Token) & Must & II \\
51 & Credit score, tier, and repayment history on SBT & Must & II \\
\rowcolor{RowShade}
52 & Tier-based max loan limits and rate modifiers & Should & II \\
53 & Wallet authentication (MetaMask / WalletConnect) & Must & I \\
\rowcolor{RowShade}
54 & EIP-712 typed-data login + JWT session & Must & I--II \\
55 & Five-stage retail onboarding funnel & Should & I--II \\
\rowcolor{RowShade}
56 & Tiered KYC (Level~1 / Level~2) & Must & I--II \\
57 & Off-chain KYC document storage (hash-linked) & Must & I--II \\
\rowcolor{RowShade}
58 & ERC-4337 account abstraction (gasless onboarding) & Stretch & III \\
59 & Groth16 zkKYC / zkAML circuits & Future Work & --- \\
\midrule
\multicolumn{4}{@{}l@{}}{\textit{AI/ML risk and explainability}} \\
\rowcolor{RowShade}
60 & Random Forest fraud probability scoring & Must (C3) & III \\
61 & Isolation Forest anomaly detection & Must (C3) & III \\
\rowcolor{RowShade}
62 & Stacking meta-learner (calibrated composite score) & Must (C3) & III \\
63 & SHAP explainability (top-$k$ features) & Must (C3) & III \\
\rowcolor{RowShade}
64 & Authority Brief UI for approvers & Must (C3) & II--III \\
65 & Commit--reveal oracle score gate before approval & Must & III \\
\rowcolor{RowShade}
66 & \texttt{LOAN\_RISK\_ASSESSMENT} audit row per inference & Must & III \\
67 & BCCC-DeFiFraudTrans benchmark evaluation & Must & II--III \\
\rowcolor{RowShade}
68 & Chainlink Functions DON oracle & Future Work & --- \\
69 & GNN (GraphSAGE) ablation & Future Work & --- \\
\rowcolor{RowShade}
70 & Federated learning across banks & Future Work & --- \\
\midrule
\multicolumn{4}{@{}l@{}}{\textit{Conversational agent}} \\
\rowcolor{RowShade}
71 & Qwen3-8B local LLM chat interface & Must & III \\
72 & MCP tool server (3--5 tools at MVT) & Must & III \\
\rowcolor{RowShade}
73 & Human confirmation gate for all write tools & Must & III \\
74 & Read tools: account state, loan status, requirements & Must & III \\
\rowcolor{RowShade}
75 & Write tools: loan application, installment payment & Must & III \\
76 & RAG over policy documents (Q\&A path) & Must & III \\
\rowcolor{RowShade}
77 & Live on-chain context injection per turn & Must & III \\
78 & Three-tier prompt assembly (Stable / Context / Volatile) & Must & III \\
\rowcolor{RowShade}
79 & Prompt injection scanning & Must & III \\
80 & \texttt{AGENT\_ACTION\_LOG} audit trail & Must & III \\
\rowcolor{RowShade}
81 & Full 17-tool MCP suite & Future Work & --- \\
82 & EIP-7702 scoped session keys for agent & Stretch & III \\
\midrule
\multicolumn{4}{@{}l@{}}{\textit{Frontend, API, and data layer}} \\
\rowcolor{RowShade}
83 & React + TypeScript dashboard (mobile-first) & Must & I--II \\
84 & Reserve transparency dashboard & Should & III \\
\rowcolor{RowShade}
85 & Five-state transaction UX machine (idle $\to$ success/error) & Should & II \\
86 & PostgreSQL event listener (on-chain $\to$ DB sync) & Must & I--II \\
\rowcolor{RowShade}
87 & Real-time notifications (WebSocket / Socket.IO) & Should & III \\
88 & Client--bank chat & Should & II \\
\rowcolor{RowShade}
89 & Express.js REST API with RBAC & Must & I--II \\
90 & FastAPI ML inference service & Must & II--III \\
\rowcolor{RowShade}
91 & API rate limiting and correlation-ID tracing & Should & II \\
92 & The Graph subgraph (optional query layer) & Future Work & --- \\
\midrule
\multicolumn{4}{@{}l@{}}{\textit{Security, governance, and compliance}} \\
\rowcolor{RowShade}
93 & Five-layer defense-in-depth architecture & Specified & I--IV \\
94 & ReentrancyGuard + CEI pattern on contracts & Must & I--II \\
\rowcolor{RowShade}
95 & LiquidationEngine (health factor, liquidator bonus) & Should & III \\
96 & SAR workflow (anomaly $\to$ freeze account) & Should & III \\
\rowcolor{RowShade}
97 & Governance timelock + snapshot voting & Specified & III+ \\
98 & Slither + Mythril static analysis & Must & IV \\
\rowcolor{RowShade}
99 & Foundry fuzz + invariant test suite & Must & II--IV \\
100 & Tenderly monitoring alerts & Should & III--IV \\
\rowcolor{RowShade}
101 & Regulatory read-only audit access & Specified & III+ \\
102 & GDPR-style data minimisation (hash-only on-chain) & Specified & I--IV \\
\midrule
\multicolumn{4}{@{}l@{}}{\textit{Evaluation and operations}} \\
\rowcolor{RowShade}
103 & 300-client Foundry / Hardhat simulation & Must & IV \\
104 & LLM red-team evaluation protocol & Should & IV \\
\rowcolor{RowShade}
105 & Open-source README with reproduction steps & Must & IV \\
\end{longtable}
\endgroup



\section{Societal Impact}
\label{sec:societal-impact}
The existing global financial system does not provide equal access to financial lending capabilities; tier-based financial flows are often kept hidden from a general client who has little to no financial literacy, and bankers often try tactics to misdirect clients into taking loans since transparency is not their priority. In general, the rich have the ability to back up loans of large amounts by showing their assets and potential profit margins and are given priorities, and the poor who might need financial help as loans for critical issues are ignored for the higher profit margins of banks. Programmable smart contracts are being utilized to generalise the loan limits and loan request submitting and approval system that improves equal access globally to those who need it regardless of financial class in a society and who have the ability to repay a loan despite their reputation set by typical judgemental views of traditional bank authorities. All the retail clients are treated equally; if the client can provide necessary documents for proof of income and stability and stakes, he or she may get loan approval from the \textit{Crypto World Bank} branch that has been set for that local client. The benefit of the smart contracts in CWB is that they are completely transparent for all the tiers, so clients can understand how currency is being made sustainable and functional for all; the blockchain comes with the added benefit of decentralisation, so no single party holds the ability to stop financial flow for one section of the world as a restriction of equal access to currency. It is not just a project to show technical design; it is made to prove that blockchain technology and cryptocurrencies can be introduced to the general mass for universal usage, and a lending platform can be built to improve the traditional tier-based banks.

\section{Environmental Impact}
\label{sec:environmental-impact}
One of the widely popular major complaints of using a blockchain network to build projects and products on it is the energy consumption and its effect on the environment. Running a proof-of-work (PoW) blockchain calculation consumes a massive amount of energy by the servers or personal computer systems, and the waste of energy compared to the end results of calculations is hard to justify.

That is why our project, CWB, runs on proof-of-stake (PoS) EVM, which is 99.9\% more energy efficient compared to PoW consensus chains. The power consumption of Bitcoin's PoW network is estimated at 120 to 150 TWh annually, which is 0.4--0.5\% of total global electricity consumption, where our approach to serve a worldwide unified lending platform should not exceed electricity consumption of a few small commercial buildings to keep the servers and network operational.

During development, Ethereum Sepolia will be used for demonstration and to prepare the testnets, and in future Polygon zkEVM Cardona is preferred for more efficient gas cost optimisation and more efficient settlement. The ML layer risk score containing fraud detection, anomaly detection and explainability runs on mid-tier commodity hardware or free cloud services; these are not very expensive services as we planned with lightweight models. ML models are not processed on-chain with the smart contracts; the loan approval, financial flow and other banking services are off-chain work, which reduces stress of on-chain calculations to minimize complexities. During the pre-thesis stage, the total estimation of carbon emission and electricity usage is hard to determine; this information will be more predictable after the final execution of the prototype.

\section{Ethical Issues}
\label{sec:ethical-issues}
The ethical statements have been clarified at the beginning of the paper. The intention of CWB is to demonstrate how this technology and development stack can be used to create a blockchain-based lending platform with AI-enhanced security and assistance. There is no intention to store and use user data, monetary profit, or to build services for financial gains. The full development process and code are open source and available for use publicly in the official GitHub repository for the project.

For the banking system, client data are mostly controlled and processed for on-chain use using hashes, and the AI-assisted techniques that have been utilized here only help the clients to navigate and suggest how to use the platform; all major decisions regarding loan approval are controlled by human approval gates. The ML layers that use client data are not stored and only used for each session time when a user visits the platform; the risk scores and SHAP explainability are designed to not expose any critical data about the users and the platform adapters.

For professional-grade usage of the platform in the future, clients are managed by fair rules and regulation; the treatment for the lending system and money allocation system follows the traditional banking system. No extraordinary social, psychological, financial systems or laws have been applied here, as this build is a prototype which may evolve in future for more robust usage and more advanced adaptability strategies may be used.

\section{Project Management Plan}
\label{sec:project-management-plan}

\subsection{Implementation Phase Plan and Deliverables}
\label{sec:impl-phase-plan}

\subsubsection{Phase I: Foundation and Infrastructure (Weeks 1--4)}

Phase~I focuses on the foundation layers, such as the smart contracts of World Bank, National Bank, and Local Bank, setting the right format and design of the database and schemas, and successfully deploying the smart contracts on Sepolia to show the contracts are functional. Priority follows the MoSCoW classification:

\enlargethispage{1.1cm}
\begin{table}[H]
\centering
\caption[Phase I task register]{Phase I task register: smart contract and database foundation.}
\label{tab:phaseI-contracts}
\scriptsize
\setlength{\tabcolsep}{2.5pt}
\renewcommand{\arraystretch}{0.98}
\begin{tabular}{p{1.6cm}p{1.5cm}p{7.7cm}p{1.2cm}}
\toprule
\tableheadcolor
\textbf{Task} & \textbf{Priority} & \textbf{Development Task} & \textbf{Effort (days)} \\
\midrule
DT-I.01 & Must & World Bank contract --- reserve management, national bank registration & 5 \\
DT-I.02 & Must & National Bank contract --- borrow from World Bank, lend to Local Banks & 5 \\
DT-I.03 & Must & Local Bank contract --- borrow from National Bank, lend to users & 5 \\
DT-I.04 & Must & RBAC --- WorldBankAdmin, NationalBankAdmin, LocalBankAdmin, BankUser, RetailClient & 3 \\
DT-I.04a & Must & \texttt{MockUSDC.sol} --- mintable ERC-20 for testnet loans (Section~\ref{sec:design-limitations}) & 1 \\
DT-I.05 & Should & Gas cost management --- initiator-pays pattern; Polygon low-fee design & 2 \\
DT-I.06 & Should & \texttt{InsuranceFund} --- 5\% interest capture, claims, default coverage & 4 \\
DT-I.07 & Should & \texttt{getReserveSummary()} on each tier contract (PSR, reserve, insurance fund) & 2 \\
DT-I.08 & Should & Chainlink Price Feeds (fiat/USD, ETH/USD via \texttt{AggregatorV3Interface}) & 3 \\
DT-I.09 & Should & Chainlink Proof of Reserve --- WorldBankReserve attestation & 2 \\
DT-I.10 & Should & Append-only \texttt{audit\_logs} PostgreSQL role + RLS policies & 2 \\
DT-I.11 & Could & Polygon zkEVM Cardona migration --- RPC, chain IDs, factory addresses & 1 \\
DT-I.12 & Must & Phase~I schema migration --- 13 M1 tables (Section~\ref{sec:db-impl-tiers}); M2/M3/stub tables deferred & 4 \\
DT-I.13 & Must & Core indexes and referential integrity constraints & 2 \\
DT-I.14 & Must & React frontend shell --- Vite build, routing, layout components & 3 \\
DT-I.15 & Must & WalletConnect --- connection, address display, network switching & 2 \\
DT-I.16 & Should & Sessions and assets tables in frontend (dashboard stubs) & 2 \\
DT-I.17 & Should & Credit tier and interest rate schedule display & 2 \\
DT-I.18 & Must & Safe 2-of-3 multisig + \texttt{WorldBankAdmin} role transfer (MVT item~15) & 1 \\
\multicolumn{3}{l}{\textit{Phase~I effort: 31 Must days; 51 all priorities.}} & \\
\bottomrule
\end{tabular}
\end{table}


\subsubsection{Phase II: Core Banking and On-Chain Lifecycle (Weeks 5--9)}

For Phase~II, we focus on building the React user interface to be responsive with loan requests and other banking functionalities such as loan approval, installments, registering accounts, and setting borrowing limits. Moreover, this phase builds the building blocks for Phase~III that will incorporate the conversational agent and the tools that the agent has access to.

\begin{table}[H]
\centering
\caption[Phase II task register]{Phase II task register.}
\label{tab:phaseII-tasks}
\begin{tabular}{p{1.8cm}p{1.8cm}p{7.5cm}p{1.5cm}}
\toprule
\tableheadcolor
\textbf{Task} & \textbf{Priority} & \textbf{Development Task} & \textbf{Effort (days)} \\
\midrule
DT-II.01 & Must & Loan request submission with blockchain transaction & 5 \\
DT-II.02 & Must & Loan approval and rejection workflow with bank approver & 5 \\
DT-II.03 & Must & Installment payment system with automated schedule generation & 8 \\
DT-II.04 & Must & Borrowing limit enforcement (on-chain authoritative check per client SBT) & 5 \\
DT-II.05 & Should & Client--bank real-time chat system & 5 \\
DT-II.06 & Should & Income verification document upload and review & 5 \\
DT-II.07 & Must & Hierarchical bank registration (National $\to$ Local) & 5 \\
DT-II.08 & Should & Bank user management and approver designation & 3 \\
DT-II.09 & Should & Loan history and transaction tracking pages & 5 \\
DT-II.10 & Could & QR code generation for wallet addresses & 2 \\
DT-II.11 & Should & Responsive UI polish and error handling & 2 \\
DT-II.12 & Should & Authority Brief UI (SHAP breakdown for bank approver dashboard) & 5 \\
DT-II.13 & Should & Installment due-date checks: Node.js cron (MVT demo path) or Chainlink Automation upgrade when LINK funded & 5 \\
DT-II.14 & Should & Credit tier schedule in Credit Passport SBT (Bronze $\to$ Diamond) & 3 \\
DT-II.15 & Should & \texttt{interest\_rate\_tier} table integration with Kinked Rate Model & 2 \\
DT-II.16 & Should & \texttt{agent\_action\_log} + \texttt{SESSIONS} schema migration (agent audit trail foundation) & 3 \\
\multicolumn{3}{l}{\textit{Phase~II total estimated effort (Must-have):}} & \textit{28 days} \\
\multicolumn{3}{l}{\textit{Phase~II total estimated effort (all priorities):}} & \textit{68 days} \\
\bottomrule
\end{tabular}
\end{table}


\subsubsection{Phase III: AI/ML Oracle Pipeline and Conversational Agent (Weeks 10--13)}

Phase~III is the continuation of Phase~II, using all lending functionalities such as loan request, approval workflow, and adding the ability of the MCP agent to interact with these tasks. The ML layer is also integrated with the oracle link and a connection between the on-chain and off-chain is established. Implementation of conversational agent harnessing including MCP server, conversational UI, and confirmation gates from both client and bank authority.

\begin{table}[H]
\centering
\caption[Phase III task register]{Phase III task register.}
\label{tab:phaseIII-tasks}
\begin{tabular}{p{1.8cm}p{1.8cm}p{7.5cm}p{1.5cm}}
\toprule
\tableheadcolor
\textbf{Task} & \textbf{Priority} & \textbf{Development Task} & \textbf{Effort (days)} \\
\midrule
DT-III.01 & Future Work & Chainlink Functions oracle (production upgrade): DON-based score commit when supported; commit--reveal is MVT demonstration path & 6 \\
DT-III.02 & Must & RF + Isolation Forest + SHAP pipeline (FastAPI); Authority Brief explainability & 7 \\
DT-III.03 & Must & MCP tool server --- minimal (3--5) tool subset wired to Express.js API; one write demo E2E & 4 \\
DT-III.04 & Must & Agent chat interface (Qwen3-8B + tool calling + human-gate) & 8 \\
DT-III.05 & Should & Session-key management (EIP-7702) and scope enforcement & 5 \\
DT-III.06 & Must & Three-tier prompt constructor in Express.js & 5 \\
DT-III.07 & Must & Prompt injection scanning + confirmation audit hook (Hook~1) & 5 \\
DT-III.08 & Should & Session lineage + \texttt{get\_session\_history} MCP tool & 5 \\
DT-III.09 & Should & Context compression path + toolset scoping + Hook~2 & 9 \\
DT-III.10 & Should & (Optional) The Graph query layer + WebSocket event listener for Reserve Transparency Dashboard & 4 \\
DT-III.11 & Should & SAR workflow: Isolation Forest $\to$ aml-alert $\to$ \texttt{freezeAccount}; compliance review queue integration & 4 \\
DT-III.12 & Should & Reserve Transparency Dashboard: live queries to event-listener index (Graph optional later) & 3 \\
DT-III.13 & Could & GraphSAGE GNN ablation & 5 \\
DT-III.14 & Could & Federated learning aggregator & 7 \\
\multicolumn{3}{l}{\textit{Phase~III total estimated effort (Must-have):}} & \textit{29 days} \\
\multicolumn{3}{l}{\textit{Phase~III total estimated effort (all priorities):}} & \textit{77 days} \\
\bottomrule
\end{tabular}
\end{table}

\subsubsection{Phase IV: Verification, Evaluation, and Finalization (Weeks 14--16)}

Implementing the Foundry fuzz suite, fully functional 300-client simulation, testnet deployment (local + remote), and the LLM agent evaluation reports, and finally paper refinements and finalization.

\begin{table}[H]
\centering
\caption[Phase IV task register]{Phase IV task register.}
\label{tab:phaseIV-tasks}
\begin{tabular}{p{1.8cm}p{1.8cm}p{7.5cm}p{1.5cm}}
\toprule
\tableheadcolor
\textbf{Task} & \textbf{Priority} & \textbf{Development Task} & \textbf{Effort (days)} \\
\midrule
DT-IV.01 & Must & 300-client Foundry simulation on Anvil/Hardhat (deterministic seed); gas-cost table and reserve-dynamics manifest to \texttt{deployments/testnet/}; optional Sepolia replay for explorer evidence & 7 \\
DT-IV.02 & Should & LLM evaluation protocol: task accuracy, action accuracy, human-gate bypass (Section~\ref{sec:optional-agent}; MVT item~12) & 5 \\
DT-IV.03 & Future Work & Certora reserve proofs: CVL specs for Invariant~1 (reserve never under-collateralised) and Invariant~2 (no over-allocation downstream); specs in Appendix~B & 5 \\
DT-IV.04 & Should & Foundry invariant + fuzz suite: solvency, role segregation, capital-flow direction; 10{,}000 fuzz runs per CI build & 4 \\
DT-IV.05 & Should & AML pre-check (Hook~3): IF anomaly gate before disbursement writes; compliance queue routing; requires Phase~III IF pipeline & 3 \\
DT-IV.06 & Could & Groth16 zkKYC (Circom~2.0): ID possession proof without PII; circuit I/O specification & 5 \\
DT-IV.07 & Must & Final paper revision: tool-count and phase refs consistent; master checklist complete; bibliography verified & 4 \\
DT-IV.08 & Must & Red-team testing: 20 injection payloads; 20 bypass attempts; 10 session-scope scenarios; results in evaluation subsections & 3 \\
DT-IV.09 & Must & Slither + Mythril on World/National/Local Bank contracts; triage and archive (MVT item~13) & 2 \\
DT-IV.10 & Should & Tenderly alerts: 3 demo-critical triggers + README reproduction (MVT items~10, 16) & 2 \\
\multicolumn{3}{l}{\textit{Phase~IV total estimated effort (Must-have):}} & \textit{16 days} \\
\multicolumn{3}{l}{\textit{Phase~IV total estimated effort (all priorities):}} & \textit{40 days} \\
\bottomrule
\end{tabular}
\end{table}

\begin{table}[H]
\centering
\caption[Four-phase roadmap]{Four-phase roadmap.}
\label{tab:phase-roadmap}
\begin{tabular}{p{1.5cm}p{3.5cm}p{2cm}p{2cm}p{3cm}}
\toprule
\tableheadcolor
\textbf{Phase} & \textbf{Focus} & \textbf{Weeks} & \textbf{Tasks} & \textbf{Key deliverables} \\
\midrule
I & Foundation and infrastructure & 1--4 & DT-I.01--18 & Contract hierarchy, \texttt{MockUSDC}, core DB schema, React shell \\
II & Core banking and on-chain lifecycle & 5--9 & DT-II.01--16 & Lending workflow, SBT limits, cron/Automation triggers \\
III & AI/ML oracle pipeline + conversational agent & 10--13 & DT-III.01--14 & Commit--reveal oracle, ML pipeline, MCP agent (minimal toolset), dashboard \\
IV & Verification, evaluation, finalization & 14--16 & DT-IV.01--10 & Foundry, simulation, evaluation pack, testnet deploy \\
\bottomrule
\end{tabular}
\end{table}

\vspace{0.35\baselineskip}
\noindent\begin{minipage}{\textwidth}
\centering
\ThesisIncludeGraphics[width=\linewidth,height=0.26\textheight,keepaspectratio]{Ch4_dev-verification-toolchain}
\captionof{figure}[Planned development and verification toolchain]{Planned development and verification toolchain}
\label{fig:dev-toolchain}
\end{minipage}



\subsection{SDLC Stage Mapping}

\begin{table}[H]
\centering
\caption[SDLC stage mapping]{SDLC stage mapping.}
\label{tab:sdlc-mapping}
\begin{tabular}{p{0.5cm}p{2.5cm}p{4.5cm}p{5cm}}
\toprule
\tableheadcolor
\textbf{\#} & \textbf{SDLC Stage} & \textbf{Project Activity} & \textbf{Deliverable} \\
\midrule
1 & Planning & Feasibility studies; professor consultations; guideline review & Feasibility Report; Project Plan \\
2 & Requirements & System analysis; use case definitions; constraints identification & Use case diagrams; UC-1 through UC-11 (eleven grouped use cases) \\
3 & Design & Four-layer architecture; DB schema (51 entities, 3NF); smart contract interfaces & Architecture diagrams; ERD; DFD \\
4 & Development & Phase I--IV: smart contracts, frontend, backend, AI/ML, agent & Source code; DApp prototype \\
5 & Testing & Hardhat unit tests; Foundry fuzz; integration testing; AI/ML + agent evaluation & Test reports; model metrics \\
6 & Deployment & Ethereum Sepolia (Phases~I--IV demonstration); Certora on Sepolia (Phase~IV stretch); frontend (Vercel); backend local-first (Render optional) & Testnet prototype manifest \\
7 & Maintenance & Monitoring; model retraining; bug fixes; iteration & Updated docs; retrained models \\
\bottomrule
\end{tabular}
\TableScreenshotEnd{Technology Stack: Additional Alternative Evaluations}
\end{table}




\section{Risk Management}
\label{sec:risk-management}
\label{sec:design-limitations}

\begin{table}[H]
\centering
\small
\caption[Design reference versus demonstration path]{Design reference versus demonstration path and key mitigations.}
\label{tab:design-vs-demo}
\label{tab:zero-cost-risks}
\begin{tabular}{@{}p{2.8cm}p{4.2cm}p{4.2cm}p{3.2cm}@{}}
\toprule
\tableheadcolor
\textbf{Area} & \textbf{Full design} & \textbf{Demo path (Ph.~I--IV)} & \textbf{Risk mitigated} \\
\midrule
Network & Polygon zkEVM Cardona & Sepolia primary; Amoy secondary & Testnet instability \\
Stablecoin & USDC / USDT & \texttt{MockUSDC.sol} & No official testnet USDC \\
Oracle & Chainlink DON & Commit--reveal relay & Oracle cost / latency \\
ML delivery & Chainlink Functions & FastAPI + local relay & Integration complexity \\
Simulation & 300 wallets live & Foundry/Anvil + 5--10 live demo & Faucet funding \\
Hosting & Cloud API + ML & Local stack + Vercel frontend & Cold-start / cost \\
\bottomrule
\end{tabular}
\end{table}

\section{Economic Analysis}
\label{sec:economic-analysis}

\subsection{Prototype Cost Model}
\label{sec:economic-feasibility}

\begin{table}[H]
\centering
\caption[Economic feasibility --- zero-cost prototype]{Economic feasibility --- zero-cost prototype.}
\label{tab:economic-feasibility}
\begin{tabular}{p{4cm}p{2.5cm}p{7cm}}
\toprule
\tableheadcolor
\textbf{Cost Category} & \textbf{Estimate} & \textbf{Notes} \\
\midrule
Blockchain deployment & \$0 & Public testnets --- no real cryptocurrency required \\
Frontend hosting & \$0 & Vercel free tier or localhost for demo \\
Backend hosting & \$0 & Render free tier or localhost \\
AI/ML training & \$0 & Local machine or Google Colab free tier \\
Development tools & \$0 & Hardhat, VS Code, Git --- all open-source \\
\textbf{Total prototype} & \textbf{\$0} & Entire prototype operates at zero financial cost \\
\bottomrule
\end{tabular}
\end{table}

\subsection{Gas Cost and Break-Even}
\label{sec:gas-cost}

\begin{table}[H]
\centering
\caption[Gas cost and break-even assumptions]{Gas cost and break-even assumptions (Polygon zkEVM Cardona design preference).}
\label{tab:gas-cost}
\begin{tabular}{p{4.2cm}p{2.8cm}p{6.8cm}}
\toprule
\tableheadcolor
\textbf{Parameter} & \textbf{Estimate} & \textbf{Notes} \\
\midrule
On-chain steps per loan lifecycle & 27--32 & Full retail lifecycle \\
Gas price (planning) & ${\approx}$30~Gwei & --- \\
ETH price (planning) & \$2{,}500 & --- \\
Total gas (Cardona / L2) & ${\approx}$\$0.011 & $<$0.01\% of interest on \$25{,}000 USDC loan at 8\% APR \\
Total gas (Ethereum mainnet) & \$2.40--\$4.80 & Same lifecycle; supports L2 design choice~[37] \\
Average loan size & \$25{,}000 & Break-even base case \\
Default rate assumption & 8\% & --- \\
Infrastructure cost & \$500/month & Pilot-scale assumption \\
Break-even active loans & ${\approx}$200 & Covers default losses at base case \\
\bottomrule
\end{tabular}
\end{table}

\subsection{Net Present Value and Return on Investment}
\label{sec:npv-roi}

\begin{table}[H]
\centering
\caption[Illustrative three-year NPV and ROI (pilot scale)]{Illustrative three-year NPV and ROI at pilot scale (300 active loans, \$25{,}000 average).}
\label{tab:npv-roi}
\small
\begin{tabular}{@{}lrrrr@{}}
\toprule
\tableheadcolor
\textbf{Year} & \textbf{Benefits (\$k)} & \textbf{Costs (\$k)} & \textbf{Net (\$k)} & \textbf{PV factor @ 10\%} \\
\midrule
0 (setup) & 0 & 6.0 & $-$6.0 & 1.000 \\
1 & 48.0 & 6.0 & 42.0 & 0.909 \\
2 & 72.0 & 8.0 & 64.0 & 0.826 \\
3 & 96.0 & 10.0 & 86.0 & 0.751 \\
\midrule
\multicolumn{2}{l}{\textbf{NPV (sum of discounted net)}} & & \multicolumn{2}{r}{\textbf{\$149.6k}} \\
\multicolumn{2}{l}{\textbf{ROI (3-yr benefits $-$ costs) / costs}} & & \multicolumn{2}{r}{\textbf{1{,}247\%}} \\
\multicolumn{2}{l}{\textbf{Break-even loans (base case)}} & & \multicolumn{2}{r}{\textbf{$\approx$300}} \\
\bottomrule
\end{tabular}
\end{table}



%% CHAPTER 4: DESIGN ALTERNATIVES
\chapter{Design Alternatives}
\label{ch:methodology}
\label{ch:design-alternatives}

This project paper is constructed to match with the final year project evaluation rubrics. Full development work has been distributed to 4 phases, and priority of all the capability sections has been divided to Must, Should and Future work. The Must and Should priority sections are to be completed by the end of the development and testing phase, Future Work sections have been specified only, to show the Potential of the project by extending its functionality and development scopes.

\begin{table}[H]
\centering
\caption[Capability priority summary]{Capability priority summary.}
\label{tab:capability-priority}
\begin{tabular}{p{3.2cm}p{2.2cm}p{7.8cm}}
\toprule
\tableheadcolor
\textbf{Capability} & \textbf{Priority} & \textbf{Primary interface / phase} \\
\midrule
Four-tier smart contracts & Must (Ph.~I--II) & React UI + wallet; Hardhat/Foundry tests \\
On-chain loan lifecycle & Must (Ph.~II) & React forms; \texttt{LoanController} state machine \\
Chainlink Functions oracle & Future Work & FastAPI $\to$ DON $\to$ on-chain score commit; commit--reveal is MVT demonstration path \\
RF + Isolation Forest + SHAP & Must (Ph.~III) & Approver Authority Brief; client decline reasons \\
Authority Brief UI & Should (Ph.~II--III) & Bank approver dashboard (React) --- \textit{not} agent-only \\
MCP conversational agent & Must (Ph.~III) & Qwen3-8B + 3--5-tool MCP server; one write demo E2E \\
LLM evaluation protocol & Should (Ph.~IV) & Red-team + task accuracy (Section~\ref{sec:optional-agent}; deferred) \\
\bottomrule
\end{tabular}
\end{table}


\section{Design Process or Methodology Overview}
\label{sec:design-process-overview}

We have developed a plan to adapt lightweight Agile/Scrum Methodology for the prototype build with an estimated development window of 16 weeks. This iterative approach enables us to explore and improve the demonstrable features that include the minimum requirements and advanced research oriented features. These feature development phases are organised in four incremental stages (Section~\ref{sec:impl-phase-plan}).

Agile Methodology application process:
\begin{itemize}[noitemsep]
\item \textbf{Phase duration:} each phase is designed for 3--5 weeks, total 4 phases, total estimated time --- 16 weeks.
\item \textbf{Team Size:} 2 developers
\item \textbf{Weekly Sync:} One weekly meeting will be conducted on Tuesday. Completion of each feature, testing and next development plans will be discussed in these meetings, based on the opinion and suggestions received from the supervisor.
\item \textbf{Phase Planning and review:} After completion of each phase, a long reflective phase meeting will be arranged, to organize and plan upcoming weeks.
\end{itemize}

\subsection{Process Model Selection}
\label{sec:process-model-selection}
\label{sec:mvc-pattern}

Based on the Software engineering process for a blockchain project, a development plan of Agile/Scrum with incremental phases has been chosen. Other models such as Waterfall, Spiral, V-Model have some flaws that do not meet the requirements to successfully build all the features.

\begin{table}[H]
\centering
\caption[Process model comparison]{Process model comparison and selection rationale.}
\label{tab:process-model-comparison}
\small
\begin{tabular}{@{}p{2.4cm}p{3.2cm}p{3.2cm}p{5.4cm}@{}}
\toprule
\tableheadcolor
\textbf{Model} & \textbf{Fit for CWB} & \textbf{Limitation} & \textbf{Decision} \\
\midrule
Waterfall & Clear phase documents & Cannot absorb evolving smart-contract and ML scope mid-project & Rejected as primary \\
Incremental & Matches four phases + gates G1--G4 & Still needs within-phase flexibility & \textbf{Adopted} (phase structure) \\
Iterative / Agile & Weekly sync; phase retrospectives & Requires discipline on scope & \textbf{Adopted} (within-phase) \\
Spiral & Strong for high-risk R\&D & Overhead excessive for 16-week student project & Not adopted \\
V-Model & Maps to test levels in Ch.~\ref{ch:evaluation} & Too rigid for research prototyping & Used for \textit{verification planning} only \\
\bottomrule
\end{tabular}
\end{table}

The project is built in the main branch from the root section. It follows a conventional \textbf{Model--View--Controller (MVC)} architecture for separation in root. \textbf{Model} --- \texttt{contracts/}: on-chain state machines and PostgreSQL (off-chain projections), \textbf{View} --- \texttt{frontend/} (React components), \textbf{Controller} --- \texttt{backend/}. For version control the repository is managed with git command, with consistency and compatibility checks, validated by \texttt{npm run test:contracts} (Hardhat) before each merge into main branch.

\subsection{Implementation Feasibility and Demonstration Scope}
\label{sec:impl-feasibility}

The full project scope---phase task registers (DT-I.01--IV.10), the MVT checklist (Table~\ref{tab:mvt-checklist}), and the module inventory in Chapter~\ref{ch:design-alternatives}---remains the project record. This section states what we will actually build and demo in the 16-week window using free testnets and local tools.

Phase~I targets roughly 31 person-days of Must work for a two-person team over four weeks. That is enough for core contracts, database setup, and a basic frontend, but only if Should tasks are not treated as gate blockers. Phase~I originally listed about 30 Must person-days; the revised plan keeps the same scope but adds explicit room for testnet debugging.

The Phase~I (G1) demo must show:
\begin{itemize}[noitemsep]
\item Three tier contracts and \texttt{MockUSDC} deployed on Sepolia
\item A working capital-flow demo (\texttt{allocateCapital} across tiers)
\item Passing RBAC tests in Hardhat or Foundry
\item A React dashboard connected via WalletConnect that reads live reserve balances
\item PostgreSQL M1 tables migrated with seed data (Section~\ref{sec:db-impl-tiers})
\item A short screen recording ($\leq$5 minutes) as the MVT~\#1 artifact
\end{itemize}

Production upgrades---Chainlink DON oracles, live USDC, full 51-table schema, and mainnet deployment---stay documented in Chapter~\ref{ch:design-alternatives} and Future Work. They are not required for the graded prototype.

\section{Preliminary Design or Design (Model) Specification}
\label{sec:preliminary-design-specification}

\subsection{System Design Overview}
\label{sec:prototype-scope}
This chapter documents the technical dependencies and logical design based on in-depth system analysis of the project. The chapter demonstrates a four-tier banking system with on-chain lending flows and off-chain data and AI/ML layer integration and optional conversational AI agent support. All technical specifications and logical explanations of their interconnectivity, and how they work as a complete system, have been documented here sequentially and structurally. The tables contain concise information and the figures are provided for visualization of each complicated section.

\subsection{High-Level Architecture}
\label{sec:high-level-arch}

The system uses a 4-layer architecture: React for frontend, application services by Express.js and FastAPI, persistence by PostgreSQL, and Solidity for on-chain contracts. The layers are summarized in Table~\ref{tab:four-layer-arch} and visually demonstrated in Figure~\ref{fig:three-layer-arch}, including how they interact.

\begin{table}[H]
\centering
\caption[Four-layer architecture technical detail]{Four-layer architecture: technical components and responsibilities.}
\label{tab:four-layer-arch}
\small
\begin{tabular}{@{}p{2.4cm}p{3.2cm}p{4.2cm}p{4.6cm}@{}}
\toprule
\tableheadcolor
\textbf{Layer} & \textbf{Technology} & \textbf{Key components} & \textbf{Responsibility} \\
\midrule
Presentation &
  React, TypeScript, WalletConnect / RainbowKit &
  Web UI, wallet signing, approver dashboard &
  User interaction and transaction initiation \\
\rowcolor{RowShade}
Application services &
  Express.js, FastAPI &
  REST API, ML scoring service, MCP agent server, event listener &
  Business logic, AI/ML inference, agent tool calls, API orchestration \\
Persistence &
  PostgreSQL, Redis (cache) &
  Lending projections, audit logs, session and agent tables &
  Off-chain data storage and read models synced from chain events \\
\rowcolor{RowShade}
On-chain &
  Solidity 0.8.20 (EVM) &
  World / National / Local Bank contracts, \texttt{LoanController}, deposit and pool modules &
  Authoritative banking state, RBAC, loan lifecycle, reserve rules \\
\bottomrule
\end{tabular}
\end{table}

\begin{figure}[H]
\centering
\BalancedDiagram{fig-three-layer-arch}
\caption[Four-layer application architecture]{Four-layer application architecture}
\label{fig:three-layer-arch}
\end{figure}

Figure~\ref{fig:component-diagram} shows component-level interactions. Core lending contracts are specified in Phase~I; extended modules are documented in Appendix~\ref{app:smart-contract-capabilities}.

\begin{figure}[H]
\centering
\BalancedDiagram{fig-component-architecture}
\caption[Component diagram with architecture-pattern labels]{Component diagram (planning phase): layered client--server topology with PostgreSQL repository and external oracle feeds.}
\label{fig:component-diagram}
\end{figure}

Figures~\ref{fig:dfd-level0-lending} and~\ref{fig:dfd-level1-lending} model the core lending subsystem in Gane-Sarson notation: Level-0 (Figure~\ref{fig:dfd-level0-lending}) shows boundary flows to five external entities; Level-1 (Figure~\ref{fig:dfd-level1-lending}) decomposes Process~0 into Processes~1.0--5.0 with data stores D1--D6.

\begin{figure}[H]
\centering
\BalancedDiagram{fig-dfd-level0-lending}
\caption[Level-0 data-flow diagram (core lending)]{Level-0 data-flow diagram (Gane-Sarson notation): context view of the core lending subsystem and its external entities.}
\label{fig:dfd-level0-lending}
\label{fig:dfd-context}
\end{figure}

\LandscapeDiagram{fig-dfd-level1-lending}
\caption[Level-1 data-flow diagram (core lending)]{Level-1 data-flow diagram (Gane-Sarson notation): decomposition of Process~0 into Processes~1.0--5.0 with data stores D1--D6.}
\label{fig:dfd-level1-lending}
\label{fig:dfd-suite}
\label{fig:dfd-level1a}
\LandscapeDiagramEnd

Figure~\ref{fig:system-component-architecture} is the detailed UML-style component diagram: four layers, required interfaces (lollipop notation), and delegation between Express.js, FastAPI, MCP agent, persistence stores, and on-chain contract modules.

\begin{figure}[H]
\centering
\TallDiagram{d1_system_component_architecture}
\caption[System component architecture]{System component architecture: presentation, application services, persistence, and on-chain layers with interface dependencies and contract groupings.}
\label{fig:system-component-architecture}
\end{figure}

\begin{table}[H]
\centering
\caption[Smart contract inventory]{Smart contract inventory: seventeen modular contracts across core, product, and multi-entity modules.}
\label{tab:contract-inventory}
\footnotesize
\setlength{\tabcolsep}{3.5pt}
\renewcommand{\arraystretch}{1.12}
\begin{tabular}{@{}>{\raggedright\arraybackslash}p{2.75cm}>{\raggedright\arraybackslash}p{1.85cm}>{\raggedright\arraybackslash}p{1.65cm}>{\raggedright\arraybackslash}p{6.85cm}@{}}
\toprule
\tableheadcolor
\textbf{Contract} & \textbf{Module} & \textbf{Status} & \textbf{Primary responsibility} \\
\midrule
WorldBankReserve & Tier~1 core & Specified (Phase~I) & Global reserve custody, national-bank registration, downward capital allocation. \\
NationalBank & Tier~2 core & Specified (Phase~I) & Local-bank registration, upward borrowing, downward allocation. \\
LocalBank & Tier~3 core & Specified (Phase~I) & Client onboarding, approver management, owns \texttt{LoanController}; \texttt{freezeAccount} (Phase~II+). \\
LoanController & Lending state machine & Ph.~I shell; Ph.~II & Loan request, approval, disbursement, installment state machine (DT-II.01--03); deployed by \texttt{LocalBank}. \\
SavingsVault & Deposits & Specified (Phase~II) & Variable-yield savings, reserve-gated withdrawals. \\
FixedDeposit & Deposits & Stub (Phase II+) & Term deposits with deterministic early-withdrawal penalty. \\
GroupLendingPool & Group credit & Specified (Phase~II) & Solidarity-group formation, consent, mutual liability. \\
FXModule & Treasury / FX & Planned (Phase~III) & Oracle-priced retail FX with disclosed spread. \\
InsuranceFund & Risk buffer & Ph.~I (often II) & Default coverage buffer; Chainlink Proof of Reserve. \\
CurrentAccount & Payments & Stub (Phase II+) & Transactional accounts and atomic peer transfers. \\
Inter\-Bank\-Lending\-Pool & Multi-entity & Specified (Phase~II) & Same-tier interbank liquidity pools per tier. \\
Upward\-Deposit\-Facility & Multi-entity & Specified (Phase~II) & Upward surplus repatriation with yield spread. \\
SyndicatedLoan & Multi-entity & Planned (Phase~III) & Multi-lender syndicate consent and disbursement. \\
TranchedPool & Multi-entity & Planned (Phase~III) & Senior/junior tranching and waterfall recovery. \\
TreasurySwap & Multi-entity & Planned (Phase~III) & Institutional oracle-priced asset swaps. \\
LiquidationEngine & Risk / liquidation & Planned (Phase~III) & Health-factor monitoring; third-party liquidator bonus when $HF < 1.0$ (Section~\ref{sec:liquidation}). \\
NettingEngine & Multi-entity & Future Work & ERD stubs only; Section~\ref{sec:future-work} \\
\bottomrule
\end{tabular}
\end{table}

\subsection{Blockchain Platform Selection}

\begin{figure}[H]
\centering
\ThesisIncludeGraphics[width=\linewidth,height=0.40\textheight,keepaspectratio]{fig-blockchain-stack}
\caption[Layered blockchain and application stack]{Layered blockchain and application stack}
\label{fig:blockchain-stack}
\end{figure}


\begin{table}[H]
\centering
\caption[Blockchain platform selection criteria and justification]{Blockchain platform selection criteria and justification.}
\label{tab:blockchain-selection}
\begin{tabular}{p{3.5cm}p{3.5cm}p{6.5cm}}
\toprule
\tableheadcolor
\textbf{Criterion} & \textbf{Selection} & \textbf{Justification} \\
\midrule
Platform & Ethereum Virtual Machine (EVM) & Largest EVM ecosystem; mature tooling. \\
Network & Ethereum Sepolia (primary) and/or Polygon Amoy PoS (secondary prototype implementation) & Free public testnets; zkEVM evaluated, not default prototype chain. \\
Consensus & Public-testnet consensus (PoS/L1 testnet) & Stable testnet consensus; ZK settlement is a design preference only. \\
Smart contract language & Solidity 0.8.20 & Industry standard; overflow-safe compiler; rich libraries. \\
\bottomrule
\end{tabular}
\TableScreenshotEnd{Blockchain Platform Selection: Comparison of EVM-Compatible Networks}
\end{table}

\begin{table}[H]
\centering
\caption[Blockchain platform selection (continued)]{Blockchain platform selection (continued): operational and deployment factors.}
\label{tab:blockchain-selection-cont}
\begin{tabular}{p{3.5cm}p{3.5cm}p{6.5cm}}
\toprule
\tableheadcolor
\textbf{Criterion} & \textbf{Selection} & \textbf{Justification} \\
\midrule
Gas cost & \$0.001--\$0.01 per transaction & Far below mainnet; viable for micro-loans. \\
Block finality & $\approx$2 seconds & Fast retail UX; L1-anchored security. \\
Developer tooling & Hardhat + OpenZeppelin & Hardhat tests; audited OpenZeppelin libs. \\
Testnet availability & Amoy (Polygon PoS), Sepolia (Ethereum) & Free faucets; EVM-identical; no real funds. \\
Security model & Public-testnet security model & Testnet PoS; ZK settlement evaluated separately. \\
Migration path & EVM-compatible & Portable Solidity across EVM chains. \\
\bottomrule
\end{tabular}
\TableScreenshotEnd{Blockchain Platform Selection (Continued): Operational and Deployment Factors}
\end{table}

\subsubsection{Transaction Verification and Consensus}

The prototype deploys on \textbf{Ethereum Sepolia} (primary) with \textbf{Polygon Amoy} as a secondary fallback (Tables~\ref{tab:blockchain-selection}). Both are public PoS testnets with free faucets and no real-asset exposure. Polygon zkEVM Cardona remains a \textbf{design-preference} target for future ZK-anchored settlement but is not the default MVP chain. The retail loan lifecycle involves multiple on-chain state transitions per client; projected gas cost stays well below 0.01\% of interest on a representative micro-loan and is validated in Phase~IV (Section~\ref{sec:gas-cost}).

\subsubsection{Oracle Architecture: Off-Chain AI to On-Chain Decision}
\label{sec:oracle-architecture}

Off-chain ML risk scores must reach the on-chain loan approval workflow without trusting a single party at decision time (the oracle problem~[40, 41]). The graded MVT path is a \textbf{commit--reveal relay}: the FastAPI scoring service commits $h = \text{keccak256}(s \,\|\, \text{nonce})$ to \texttt{LoanController}, then reveals the score~$s$ and nonce within the decision window. The contract verifies the hash and stores $s$ immutably, giving an audit trail and preventing post-commit score changes. A trusted backend relay may assist integration in Phases~I--II only.

Figure~\ref{fig:oracle-architecture} summarises the authoritative MVT path and the documented Chainlink upgrade routes. Table~\ref{tab:chainlink-upgrade-paths} maps each Chainlink service to its production role and the zero-cost demonstration substitute (Section~\ref{sec:design-limitations}; Table~\ref{tab:design-vs-demo}).

\begin{table}[H]
\centering
\caption[Chainlink upgrade paths versus MVT substitutes]{Chainlink upgrade paths versus MVT demonstration substitutes.}
\label{tab:chainlink-upgrade-paths}
\small
\begin{tabular}{@{}p{2.6cm}p{4.8cm}p{5.8cm}@{}}
\toprule
\tableheadcolor
\textbf{Service} & \textbf{Production role} & \textbf{MVT / demo substitute} \\
\midrule
Chainlink Functions &
  DON-signed fetch of ML risk score from the off-chain service &
  Commit--reveal relay (authoritative at MVT scope) \\
\rowcolor{RowShade}
Chainlink Automation &
  On-chain upkeep for overdue installments and interest accrual &
  Node.js cron calling \texttt{checkOverdueLoans()} \\
Price feeds &
  \texttt{FXModule} and fiat-equivalent UI via \texttt{AggregatorV3Interface} &
  \texttt{MockPriceFeed.sol} with fixed testnet answers \\
\rowcolor{RowShade}
Proof of Reserve &
  External attestation of \texttt{WorldBankReserve} solvency &
  On-chain \texttt{getReserveSummary()} view; PoR job is Future Work \\
\bottomrule
\end{tabular}
\end{table}

Chainlink Functions adds roughly 30--60\,s latency and a per-call fee, acceptable at loan origination but unnecessary for the graded prototype. Production contract interfaces for reserve reporting and FX conversion are specified in Appendix~D; integration code remains in the repository rather than in this chapter.

\begin{figure}[H]
\centering
\BalancedDiagram{fig-oracle-architecture}
\caption[Oracle architecture]{Oracle architecture}
\label{fig:oracle-architecture}
\end{figure}


\subsection{Data Model and Database Design}
\label{sec:data-model}

This section describes the PostgreSQL data model for Crypto World Bank. The complete logical schema defines \textbf{51 entities} in Third Normal Form (3NF); Table~\ref{tab:db-entities} summarises the \textbf{16 core entities} that drive the prototype, and Table~\ref{tab:db-entities-ext2} groups the remaining extension and stub relations by phase.

The graded prototype \textbf{migrates about 32 tables} across Phases~I--III (Section~\ref{sec:db-impl-tiers}); Gate~G1 requires 13 tables only. The full 51-entity relational ERD is in Appendix~\ref{app:db-schema-reference} (Figure~\ref{fig:erd-relational}) because of its size; Figures~\ref{fig:core-system-graph} and~\ref{fig:erd-extended} below provide readable focused views of the core lending graph and banking-product extensions.

\begin{figure}[H]
\centering
\BalancedDiagram{fig-erd-core}
\caption[Core system entity graph]{Core system entity graph}
\label{fig:core-system-graph}
\end{figure}

\begin{figure}[H]
\centering
\BalancedDiagram{fig-erd-extended}
\caption[Extended ERD: banking products]{Extended ERD: banking-product entities (crow's-foot; PK shown; entity roles in Table~\ref{tab:db-entities-ext2}).}
\label{fig:erd-extended}
\end{figure}

Figure~\ref{fig:erd-extended} documents banking-product extensions; FK anchors into core lending entities are in Table~\ref{tab:ext-entity-fks} (Section~\ref{sec:multi-entity-consistency}).

\subsubsection{Entity Summary}

\begin{table}[H]
\centering
\caption[Core database entities]{Core database entities: primary keys and roles (Phase~I--III prototype scope).}
\label{tab:db-entities}
\footnotesize
\setlength{\tabcolsep}{3.5pt}
\renewcommand{\arraystretch}{1.06}
\resizebox{\textwidth}{!}{%
\begin{tabular}{@{}>{\raggedright\arraybackslash}p{2.4cm}>{\raggedright\arraybackslash}p{2.4cm}>{\raggedright\arraybackslash}p{8.5cm}@{}}
\toprule
\tableheadcolor
\textbf{Entity} & \textbf{Primary key} & \textbf{Role} \\
\midrule
\multicolumn{3}{l}{\textit{Institution hierarchy (M1)}} \\
\midrule
INSTITUTION & \texttt{institution\_id} & Shared supertype for World, National, and Local bank identities; enables single FK target for cross-tier operations. \\
COUNTRY & \texttt{country\_code} & ISO~3166-1 reference; one \texttt{NATIONAL\_BANK} per country (\texttt{UNIQUE}). \\
WORLD\_BANK & \texttt{institution\_id} & Global reserve subtype; singleton (\texttt{CHECK institution\_id = 1}). \\
NATIONAL\_BANK & \texttt{institution\_id} & National reserve subtype; FK \texttt{country\_code}, \texttt{parent\_world\_bank\_id}. \\
LOCAL\_BANK & \texttt{institution\_id} & City-level lending institution; FK \texttt{national\_bank\_id}. \\
\midrule
\multicolumn{3}{l}{\textit{Actors and lending lifecycle (M1)}} \\
\midrule
BANK\_USER & \texttt{bank\_user\_id} & Bank staff; exactly one parent institution enforced by CHECK (Table~\ref{tab:integrity-constraints}). \\
BORROWER & \texttt{borrower\_id} & Retail client (UK: \texttt{wallet\_address}); home bank via \texttt{registered\_local\_bank\_id}; hash-only KYC fields. \\
LOAN\_REQUEST & \texttt{request\_id} & Application workflow (\texttt{status}, \texttt{oracle\_state}); links borrower to local bank. \\
LOAN & \texttt{loan\_id} & Active loan after disbursement; FK \texttt{loan\_request\_id}; on-chain linkage fields. \\
INSTALLMENT & \texttt{loan\_id}, \texttt{installment\_number} & Repayment schedule (weak entity on \texttt{LOAN}). \\
\midrule
\multicolumn{3}{l}{\textit{Sync, assets, and audit (M1)}} \\
\midrule
BLOCKCHAIN\_EVENT\_LOG & \texttt{event\_id} & Canonical on-chain event cache (event-listener owned); FK target for lending tables. \\
ASSETS & \texttt{asset\_id} & Collateral and loan asset registry (UK: \texttt{symbol}). \\
AUDIT\_LOGS & \texttt{audit\_id} & Append-only governance and compliance log (Express-owned). \\
\midrule
\multicolumn{3}{l}{\textit{Phase~II--III extensions (selected)}} \\
\midrule
TRANSACTION & \texttt{transaction\_id} & Financial ledger for retail and institutional flows. \\
BORROWING\_LIMIT & \texttt{limit\_id} & Off-chain projection of on-chain rolling borrow limits. \\
CREDIT\_PASSPORT & \texttt{passport\_id} & Event-listener projection of CreditPassport SBT state. \\
LOAN\_RISK\_ASSESSMENT & \texttt{assessment\_id} & Per-loan ML inference record; FK \texttt{model\_id} $\to$ \texttt{MODEL\_REGISTRY}. \\
SESSIONS & \texttt{session\_id} & Retail agent wallet session (Phase~III); FK \texttt{borrower\_id}. \\
\bottomrule
\end{tabular}%
}
\end{table}

\begin{table}[H]
\centering
\caption[Deferred and extension entities]{Deferred and extension entities (35 relations): grouped by implementation phase. Full attribute lists are omitted; see ERD figures and Section~\ref{sec:db-impl-tiers}.}
\label{tab:db-entities-ext2}
\footnotesize
\setlength{\tabcolsep}{3.5pt}
\renewcommand{\arraystretch}{1.05}
\resizebox{\textwidth}{!}{%
\begin{tabular}{@{}>{\raggedright\arraybackslash}p{1.2cm}>{\raggedright\arraybackslash}p{5.8cm}>{\raggedright\arraybackslash}p{6.5cm}@{}}
\toprule
\tableheadcolor
\textbf{Tier} & \textbf{Entities} & \textbf{Purpose} \\
\midrule
M2 & \texttt{CREDIT\_PASSPORT\_HISTORY}, \texttt{INCOME\_PROOF}, \texttt{INTEREST\_RATE\_TIER} & SBT history, Level~2+ document hashes, normalised kinked rates \\
M2 & \texttt{GROUP\_CONSENT}, \texttt{LOAN\_GROUP}, \texttt{GROUP\_MEMBER} & Group lending (if \texttt{GroupLendingPool} demo chosen) \\
M2 & \texttt{INTERBANK\_LOAN}, \texttt{UPWARD\_DEPOSIT}, \texttt{SAVINGS\_ACCOUNT} & One multi-entity PoC path \\
M3 & \texttt{MODEL\_REGISTRY}, \texttt{SECURITY\_EVENT\_LOG} & ML lineage and system-wide security events \\
M3 & \texttt{AGENT\_ACTION\_LOG}, \texttt{AGENT\_CONVERSATION\_TURN}, \texttt{SESSION\_ANCESTOR} & Agent audit and session lineage \\
Stub & \texttt{FIXED\_DEPOSIT}, \texttt{CURRENT\_ACCOUNT}, \texttt{INSURANCE\_FUND} & Deposit and risk-buffer products \\
Stub & \texttt{SYNDICATE}, \texttt{SYNDICATE\_MEMBER}, \texttt{TRANCHED\_POOL}, \texttt{TRANCHED\_POOL\_SUBSCRIPTION} & Syndication and tranched lending \\
Stub & \texttt{TREASURY\_SWAP}, \texttt{NETTING\_BATCH}, \texttt{NETTING\_ENTRY} & Treasury swap and netting (Future Work) \\
Future & \texttt{CHAT\_MESSAGE}, \texttt{INSTITUTION\_MESSAGE}, \texttt{PROFILE\_SETTING} & UX messaging and preferences \\
Future & \texttt{MARKET\_DATA}, \texttt{CHAINLINK\_PRICE\_ROUND}, \texttt{ORACLE\_HEARTBEAT\_MONITOR} & Oracle caches (mock at MVT) \\
Future & \texttt{COLLATERAL\_POSITION}, \texttt{KYC\_VERIFICATION\_REQUEST}, \texttt{NETTING\_COORDINATOR\_STATE} & Liquidation, ZKP KYC, settlement coordinator \\
Future & \texttt{CHAIN\_REGISTRY}, \texttt{SESSION\_KEY\_PERMISSION} & Phase~IV platform configuration \\
\bottomrule
\end{tabular}%
}
\end{table}

\subsubsection{Logical Schema versus Physical Migration Tiers}
\label{sec:db-impl-tiers}

The 51-entity count describes the \textbf{complete logical data model} for architectural consistency with the seventeen-contract inventory. The graded prototype does \textit{not} require 51 physical PostgreSQL tables at Gate~G1. \textbf{Phase~I (M1, 13 tables)} covers the institution hierarchy, actors, lending lifecycle, and event sync---the Gate~G1 minimum. \textbf{Phase~II ($\sim$12 tables)} adds ledger, credit passport, and rate-tier tables plus one optional product path. \textbf{Phase~III ($\sim$7 tables)} adds ML, agent session, and security-log tables. Remaining entities stay as ERD stubs or Future Work (Table~\ref{tab:db-entities-ext2}).

\paragraph{Platform-configuration relations (outside the 51-entity count).}
\texttt{CHAIN\_REGISTRY} (deployment chain IDs, RPC endpoints, EIP-7702 delegation contract addresses) and \texttt{SESSION\_KEY\_PERMISSION} (normalised MCP tool allowlists per session) are Phase~IV platform-configuration tables. They are referenced in Sections~\ref{sec:onboarding} and~\ref{sec:optional-agent} but are \textit{not} included in the 51-entity lending schema because they hold deployment metadata rather than banking-domain state.

\begin{table}[H]
\centering
\small
\setlength{\tabcolsep}{4pt}
\renewcommand{\arraystretch}{1.12}
\caption[Audit and logging table taxonomy]{Audit and logging table taxonomy: separate relations for on-chain events, governance actions, ML inference, and system security.}
\label{tab:db-audit-taxonomy}
\begin{tabular}{@{}>{\raggedright\arraybackslash}p{3.4cm}>{\raggedright\arraybackslash}p{2.2cm}>{\raggedright\arraybackslash}p{3.2cm}>{\raggedright\arraybackslash}p{4.7cm}@{}}
\toprule
\tableheadcolor
\tableheadfont Table &
\tableheadfont Writer &
\tableheadfont Scope &
\tableheadfont Typical contents \\
\midrule
\texttt{BLOCKCHAIN\_EVENT\_LOG} &
  Event listener &
  On-chain receipts &
  \texttt{tx\_hash}, \texttt{event\_name}, \texttt{raw\_data}; FK target for \texttt{LOAN\_REQUEST}, \texttt{LOAN}, \texttt{INSTALLMENT} \\
\rowcolor{RowShade}
\texttt{LOAN\_RISK\_ASSESSMENT} &
  FastAPI ML service &
  Per-loan inference &
  \texttt{fraud\_probability}, \texttt{shap\_vector}, FK \texttt{model\_id} $\to$ \texttt{MODEL\_REGISTRY} \\
\texttt{AUDIT\_LOGS} &
  Express API &
  Governance / compliance &
  Travel Rule packets, role grants, manual approver actions (Section~\ref{sec:regulatory-compliance}) \\
\rowcolor{RowShade}
\texttt{SECURITY\_EVENT\_LOG} &
  ML / monitoring &
  System-wide security &
  AML alerts, anomaly spikes, agent injection blocks \\
\texttt{AGENT\_ACTION\_LOG} &
  Express agent hooks &
  Agent write tools &
  \texttt{tool\_name}, \texttt{confirmation\_turn\_id}, \texttt{onchain\_tx\_hash} \\
\bottomrule
\end{tabular}
\end{table}

\subsubsection{EER Constructs Applied}

\begin{table}[H]
\centering
\caption[EER constructs applied]{EER constructs applied: specialisation, hierarchy, and constraints.}
\label{tab:eer-constructs}
\begin{tabular}{p{3.5cm}p{3.5cm}p{6.5cm}}
\toprule
\tableheadcolor
\textbf{EER Construct} & \textbf{Applied To} & \textbf{Notes} \\
\midrule
Specialisation (disjoint) & BANK\_USER $\to$ WorldBankUser, NationalBankUser, LocalBankUser & A bank user belongs to exactly one bank type (\texttt{world}, \texttt{national}, or \texttt{local}); enforced by CHECK constraint (Table~\ref{tab:integrity-constraints}). \\
Generalisation / specialisation & INSTITUTION $\to$ WORLD\_BANK, NATIONAL\_BANK, LOCAL\_BANK & ID-based table-per-type inheritance: subtype PK = FK to \texttt{INSTITUTION.institution\_id}; shared identity for cross-tier FKs (Table~\ref{tab:db-entities}). \\
Weak entity & INSTALLMENT depends on LOAN & Existence and identity determined by parent loan. \\
Multi-valued attribute & INCOME\_PROOF per BORROWER & Modelled as separate entity to maintain 1NF. \\
Aggregation & LOAN\_RISK\_ASSESSMENT $\leftarrow$ LOAN\_REQUEST; SECURITY\_EVENT\_LOG $\leftarrow$ runtime monitors & Per-loan ML scores and system-wide security events are separate append-only relations, not a conflated aggregate. \\
Association entity & LOAN\_REQUEST links BORROWER + LOCAL\_BANK & Captures many-to-many borrowing relationships; distinct from \texttt{BORROWER.registered\_local\_bank\_id} (home bank at onboarding). \\
Participation constraint & BORROWER.\texttt{registered\_local\_bank\_id} NOT NULL at onboarding & Home Local Bank is mandatory; \texttt{LOAN\_REQUEST} participation for lending services enforced at application layer. \\
Append-only table (policy) & AGENT\_ACTION\_LOG, AUDIT\_LOGS, SECURITY\_EVENT\_LOG & DB-level INSERT-only role + Row-Level Security policy enforces that no UPDATE or DELETE is permitted on audit records. \\
\bottomrule
\end{tabular}
\TableScreenshotEnd{EER Constructs Applied: Specialization, Hierarchy, and Constraints}
\end{table}

\subsubsection{Normalization}
\label{sec:normalization}

The schema has been normalized to Third Normal Form (3NF) and checked against Boyce-Codd Normal Form (BCNF):

\begin{itemize}[noitemsep]
\item \textbf{1NF:} There are no repeating groups in any of the attributes. Separate entities store multi-valued attributes, such as proof of income.
\item \textbf{2NF:} There are no partial dependencies. The full composite key (\texttt{loan\_id}, \texttt{installment\_number}) determines the non-key attributes of INSTALLMENT.
\item \textbf{3NF:} No dependencies that go through other dependencies. Instead of storing redundant values like \texttt{total\_borrowed} in bank entities, they are calculated at query time.
\item \textbf{BCNF:} All non-trivial functional dependencies have a superkey as their determinant. For example, in \texttt{INTEREST\_RATE\_TIER}, \texttt{tier\_id} $\to$ \{\texttt{base\_rate}, \texttt{kink\_utilisation}, \texttt{rate\_above\_kink}, \texttt{max\_rate}\}: the determinant \texttt{tier\_id} is the primary key, satisfying BCNF. Prior to the normalisation fix in this section, interest-rate parameters were columns on the three bank entities, determined by \texttt{bank\_tier\_id} (a non-key attribute), violating BCNF.
\end{itemize}

A transitive dependency was identified during schema review: \texttt{interest\_rate\_parameters} $\to$ \texttt{bank\_tier\_id} (interest rate values were stored as columns in the World Bank, National Bank, and Local Bank entities, making them dependent on a non-key attribute via the tier relationship). This violates 3NF. The fix is to extract these parameters into a dedicated \texttt{INTEREST\_RATE\_TIER} table keyed by \texttt{tier\_id}, which each bank entity references via FK. This normalisation change ensures that updating a base rate does not require touching multiple bank-tier rows.

\subsubsection{Relational Integrity Constraints}
\label{sec:integrity-constraints}

\begin{table}[H]
\centering
\footnotesize
\setlength{\tabcolsep}{3pt}
\renewcommand{\arraystretch}{1.05}
\caption[Relational integrity constraints]{Relational integrity constraints.}
\label{tab:integrity-constraints}
\resizebox{\textwidth}{!}{%
\begin{tabular}{@{}>{\raggedright\arraybackslash}p{2.85cm}>{\raggedright\arraybackslash}p{10.0cm}@{}}
\toprule
\tableheadcolor
\textbf{Constraint Type} & \textbf{Examples} \\
\midrule
Primary Key & One surrogate or composite key per entity (Tables~\ref{tab:db-entities},~\ref{tab:db-entities-ext2}): e.g.\ \texttt{institution\_id}, \texttt{borrower\_id}, \texttt{request\_id}, (\texttt{loan\_id}, \texttt{installment\_number}), \ldots \\
Foreign Key & Tier subtypes $\to$ \texttt{INSTITUTION}; \texttt{BORROWER.registered\_local\_bank\_id} $\to$ \texttt{LOCAL\_BANK}; \texttt{LOAN\_RISK\_ASSESSMENT} $\to$ \texttt{LOAN\_REQUEST}, \texttt{MODEL\_REGISTRY}; multi-entity stubs $\to$ \texttt{INSTITUTION.institution\_id}. \\
UNIQUE & \texttt{BORROWER.wallet\_address}; \texttt{NATIONAL\_BANK(country\_code)}; \texttt{INSTITUTION.contract\_address}; \texttt{LOAN\_REQUEST.blockchain\_tx\_hash}; \texttt{BORROWING\_LIMIT.borrower\_id}; \texttt{ASSETS.symbol}; \texttt{BLOCKCHAIN\_EVENT\_LOG.tx\_hash}; partial unique on active \texttt{INTEREST\_RATE\_TIER(bank\_tier\_type)}. \\
CHECK & \texttt{WORLD\_BANK}: singleton (\texttt{institution\_id = 1}). \texttt{BANK\_USER}: exactly one parent FK matching \texttt{bank\_type}. \texttt{TRANSACTION}: at least one of \texttt{borrower\_id}, \texttt{origin\_institution\_id}, \texttt{counterparty\_institution\_id} non-null. \\
NOT NULL & \texttt{INSTITUTION.name}, \texttt{INSTITUTION.institution\_type}; \texttt{BORROWER.wallet\_address}, \texttt{registered\_local\_bank\_id}, \texttt{kyc\_level}; \texttt{status}; \texttt{oracle\_state} (default \texttt{NONE}). \\
APPEND-ONLY & \texttt{AGENT\_ACTION\_LOG}, \texttt{CREDIT\_PASSPORT\_HISTORY}, \texttt{SECURITY\_EVENT\_LOG}, \texttt{AUDIT\_LOGS}: INSERT-only role + RLS; no UPDATE/DELETE. \\
PROJECTION (RLS) & \texttt{CREDIT\_PASSPORT}, \texttt{BORROWING\_LIMIT}, \texttt{COLLATERAL\_POSITION}: \texttt{event\_listener} role only may INSERT/UPDATE; Express API SELECT-only. \\
FOREIGN KEY & \texttt{LOAN} $\to$ \texttt{ASSETS} (collateral, loan asset); \texttt{AGENT\_ACTION\_LOG} $\to$ \texttt{SESSIONS}, \texttt{AGENT\_CONVERSATION\_TURN}; \texttt{LOAN.disbursement\_event\_log\_id} $\to$ \texttt{BLOCKCHAIN\_EVENT\_LOG}. \\
\bottomrule
\end{tabular}%
}
\TableScreenshotEnd{Relational Integrity Constraints: Referential and Domain Rules}
\end{table}

\vspace{2pt}
\noindent\textit{\small This integrity constraints table summarizes referential and domain constraints that keep loan, repayment, and identity records consistent. These constraints reduce operational risk by preventing orphaned records and invalid lifecycle states, which is critical for auditability in financial systems.}


\subsection{On-Chain and Off-Chain Data Partitioning}
\label{sec:data-partitioning}

\begin{table}[H]
\centering
\caption[Data partitioning between on-chain and off-chain storage]{Data partitioning between on-chain and off-chain storage.}
\label{tab:data-partitioning}
\begin{tabular}{p{4.5cm}p{2.5cm}p{6.5cm}}
\toprule
\tableheadcolor
\textbf{Data Category} & \textbf{Storage} & \textbf{Rationale} \\
\midrule
Reserve balances, loan requests, approval/rejection events, repayment transactions & On-chain & Immutability, public auditability, trustless verification. \\
User profiles, income verification documents, chat messages, loan risk assessments, security event logs & Off-chain (database) & Data privacy, query flexibility, storage cost optimisation. \\
Borrowing limit computations & Off-chain with on-chain enforcement & Complex temporal aggregation; results committed as on-chain constraints via \texttt{LocalBank.updateBorrowingLimit(borrowerId, newLimit)} (see sync protocol below). \\
Cryptocurrency market data (non-Chainlink) & Off-chain (\texttt{MARKET\_DATA}) & CoinGecko and other fallback feeds; high-frequency updates; external API dependency. \\
Chainlink oracle price data & Off-chain (\texttt{CHAINLINK\_PRICE\_ROUND}) & Each \texttt{latestRoundData()} round cached with \texttt{answered\_in\_round} for staleness detection; on-chain \texttt{AggregatorV3Interface} is authoritative. \texttt{MARKET\_DATA} does \textit{not} store Chainlink rounds. \\
Agent action execution log & Off-chain (append-only DB) & Immutable audit trail of every agent write-tool call, linked to on-chain tx hash; not stored on-chain to save gas. \\
Session key material & Off-chain (\texttt{sessions} table) & EIP-7702 session key hash, scope JSON, and TTL stored off-chain; the scope restrictions are enforced on-chain by the session key contract at signing time. \\
Session lineage chains (SESSIONS, message history, compression summaries) & Off-chain (PostgreSQL) & Full transcript history is too large for on-chain storage. The lineage chain enables searchable long-term memory without gas cost. \texttt{AGENT\_ACTION\_LOG} provides the on-chain audit link via confirmed transaction hashes. \\
\bottomrule
\end{tabular}
\TableScreenshotEnd{On-Chain vs.~Off-Chain Data Partitioning: State, Analytics, and Privacy Boundaries}
\end{table}

\paragraph{Borrowing limit sync protocol.}
Borrowing limits are enforced on-chain, but the rolling exposure that feeds those limits is derived from off-chain transaction history. When a \texttt{LoanDisbursed} or \texttt{LoanRepaid} event arrives, the event listener first records it in \texttt{TRANSACTION} and recalculates the borrower's six-month and one-year windows. If the new figure no longer matches \texttt{on\_chain\_enforced\_limit} stored in \texttt{BORROWING\_LIMIT}, the Express API submits an on-chain update through \texttt{LocalBank.updateBorrowingLimit}. Once the chain emits \texttt{BorrowingLimitUpdated}, the listener refreshes the projection row and clears any \texttt{sync\_discrepancy\_flag}, keeping the database aligned with the authoritative contract state.

\paragraph{Projection vs.\ application tables.}
The database splits cleanly between tables the application writes and tables the event listener maintains as read models. Projection tables---\texttt{CREDIT\_PASSPORT}, \texttt{CREDIT\_PASSPORT\_HISTORY}, \texttt{BORROWING\_LIMIT}, and \texttt{COLLATERAL\_POSITION}---are populated only from chain events; Express.js may read them for dashboards and API responses but cannot insert or update them directly, a rule enforced by PostgreSQL row-level security (Table~\ref{tab:integrity-constraints}). Workflow tables such as \texttt{LOAN\_REQUEST}, \texttt{CHAT\_MESSAGE}, and \texttt{PROFILE\_SETTING} remain Express-owned. \texttt{LOAN\_REQUEST} is the one hybrid case: most columns are application-writable, but \texttt{oracle\_state} is reserved for the listener, with column-level privileges preventing the API role from modifying it while still allowing the listener to track commit--reveal oracle progress independently of user-facing application state.

\begin{figure}[H]
\centering
\ThesisIncludeGraphics[width=\linewidth,height=0.34\textheight,keepaspectratio]{fig-data-partitioning}
\caption[On-chain versus off-chain data partitioning]{On-chain versus off-chain data partitioning}
\label{fig:data-partitioning}
\end{figure}

\subsubsection{Financial Data Authority and Lifecycle}
\label{sec:financial-data-lifecycle}

On-chain contract state is the \textbf{authoritative source of truth} for settlement: reserve balances, loan lifecycle transitions, role bindings, and enforced borrowing limits. PostgreSQL holds workflow metadata, PII, ML inference records, and \textbf{read-only projections} synchronised from chain events by the event listener (FR-07; Section~\ref{sec:realtime}). Where the two layers disagree after a confirmed transaction, the on-chain record prevails; projection tables exist for fast queries and dashboards, not for independent financial ledgering.

\paragraph{Retail loan write path.}
The end-to-end flow for a standard retail loan is: (1)~the client submits an application (wallet-signed \texttt{requestLoan} on \texttt{LocalBank}, with optional off-chain document hash in \texttt{LOAN\_REQUEST}); (2)~the ML service scores features off-chain and the commit--reveal oracle records the risk score on-chain (Section~\ref{sec:oracle-architecture}); (3)~the Local Bank approver signs \texttt{approveLoan}, which atomically updates loan state and disburses funds; (4)~the event listener ingests \texttt{LoanRequested}, \texttt{LoanApproved}, and \texttt{LoanDisbursed} events into \texttt{BLOCKCHAIN\_EVENT\_LOG} and updates \texttt{LOAN}, \texttt{INSTALLMENT}, and projection tables (Figure~\ref{fig:seq-loan-flow}). Repayments follow the same pattern: \texttt{payInstallment} on-chain, then listener projection.

\vspace{2pt}
\noindent\textit{\small At MVT scope, loan denomination may use \texttt{MockUSDC} or native ETH on the demonstration testnet (Table~\ref{tab:design-vs-demo}); the authority rules above are unchanged---only the settlement asset differs.}

\begin{figure}[H]
\centering
\ThesisIncludeGraphics[width=\linewidth,height=0.36\textheight,keepaspectratio]{fig-financial-data-lifecycle}
\caption[Financial data authority and lifecycle]{Financial data authority and lifecycle: on-chain settlement state, off-chain workflow and projections, and the event-listener bridge.}
\label{fig:financial-data-lifecycle}
\end{figure}

\subsection{Operating Context}
\label{sec:operating-context}

The Crypto World Bank is specified as a \textbf{blockchain-based institutional finance platform} on EVM-compatible infrastructure, not as a geographic pilot for a single rural market. The primary contribution is programmable hierarchical banking---reserve enforcement, tiered capital flow, on-chain auditability, and modular lending products---deployable on a stability-first prototype testnet (Ethereum Sepolia as primary, Polygon Amoy PoS as secondary). Polygon zkEVM remains an evaluated option in the platform comparison tables rather than the default prototype deployment chain narrative.

Three design choices anchor the operating context. First, \textbf{stablecoin-denominated retail lending} (Section~\ref{sec:stablecoin-first}) limits liability-side volatility when collateral and loans use different asset types. Second, \textbf{low per-transaction gas cost} on zkEVM L2 makes high-frequency installment and reserve operations economically viable relative to Ethereum mainnet. Third, wallet-based onboarding (Section~\ref{sec:onboarding}) keeps retail access simple without changing contract role logic. Solidarity group lending (Section~\ref{sec:banking-products}) is a programmable product module informed by microfinance literature (BRAC, Grameen); it is not scoped to a specific country deployment. Institutional onboarding is planned through academic pilots, open-source replication, and regulatory sandbox partnerships (Section~\ref{sec:regulatory-compliance}).

\subsection{User Taxonomy and Onboarding Flows}
\label{sec:onboarding}
\label{sec:tiered-kyc}
\label{sec:identity}

\noindent On-chain roles are bound to wallet addresses; KYC document hashes stay off-chain. Tables~\ref{tab:user-taxonomy}--\ref{tab:permission-matrix} and Figure~\ref{fig:permission-matrix} define actors and permissions. Groth16 zkKYC is Future Work (Section~\ref{sec:future-work}).

\begin{table}[H]
\centering
\caption[Operational user taxonomy]{Operational user taxonomy: KYC level, wallet type, and data residence by user type.}
\label{tab:user-taxonomy}
\begin{tabular}{p{3.1cm}p{1.4cm}p{2.6cm}p{2.4cm}p{3.4cm}}
\toprule
\tableheadcolor
\textbf{User Type} & \textbf{Tier} & \textbf{KYC Level} & \textbf{Wallet Type} & \textbf{Data Residence} \\
\midrule
Anonymous Visitor & --- & None & None & Session only (Redis) \\
Registered Retail User & Tier 4 & Level 1 (gov.\ ID + selfie) & MetaMask or optional ERC-4337 & Off-chain primary; on-chain role only \\
Group Client & Tier 4 & Level 1 + group verification & ERC-4337 abstracted & Off-chain + on-chain group contract \\
Local Bank Operator & Tier 3 & Full institutional KYC & MetaMask (institutional) & Off-chain + on-chain role binding \\
Local Bank Approver & Tier 3 & Full institutional + background check & MetaMask + Safe co-signer & Off-chain + on-chain role binding \\
National Bank Admin & Tier 2 & Full institutional + regulatory license & Safe multisig mandatory & Off-chain + on-chain role binding \\
World Bank Admin (Governance) & Tier 1 & Founding team + supervisor attestation & Safe 3-of-5 multisig & On-chain governance only \\
\bottomrule
\end{tabular}
\end{table}

\begin{table}[H]
\centering
\caption[Formal actor taxonomy]{Formal actor taxonomy (nine actors).}
\label{tab:actor-taxonomy}
\begin{tabular}{p{1.2cm}p{4.2cm}p{7.9cm}}
A1a & Retail Client (pre-KYC) & Browsing only; no banking transactions. \\
A1b & Retail Client (KYC Tier~1) & Bronze/Silver credit tier; small loans up to the KYC-1 cap. \\
A1c & Retail Client (KYC Tier~2) & Gold/Platinum/Diamond; larger limits, group lending. \\
A2 & Local Bank Admin (Approver) & Loan approver; risk officer; reviews AML alerts. \\
A3 & National Bank Admin & Capital allocator; compliance officer; SAR review; freeze authority. \\
A4 & World Bank Admin (Governance) & Parameter governor; system operator via multi-sig. \\
A5 & AI Agent (retail) & Phase~III Must deliverable; read tools always permitted; write tools require explicit human confirmation gate before execution. \\
\multicolumn{3}{l}{\textit{Secondary Actors (external systems or authorities)}} \\
A6 & Regulatory Authority & Read-only audit access via encrypted data package. \\
A7 & Chainlink DON & Oracle, price feeds, automation. \\
A8 & Blockchain Validator & Polygon zkEVM validity proof generation. \\
A9 & External Auditor & Chainlink PoR verification. \\
\bottomrule
\end{tabular}
\end{table}

\begin{table}[H]
\centering
\caption[Actor permission matrix (tabular summary)]{Actor permission matrix (tabular summary of primary actions).}
\label{tab:permission-matrix}
\begin{tabular}{p{4.4cm}p{1.3cm}p{1.5cm}p{1.5cm}p{1.5cm}p{1.6cm}}
\toprule
\tableheadcolor
\textbf{Action} & \textbf{Client} & \textbf{LB Admin} & \textbf{NB Admin} & \textbf{WB Admin} & \textbf{AI Agent} \\
\midrule
View own loan status & YES & YES & YES & YES & READ \\
Submit loan application & YES & NO & NO & NO & WRITE \\
Pay installment & YES & NO & NO & NO & WRITE \\
Submit KYC documents & YES & NO & NO & NO & NO \\
Approve loan application & NO & YES & NO & NO & NO \\
Reject loan application & NO & YES & NO & NO & NO \\
Freeze client account & NO & YES & YES & YES & NO \\
Set borrowing rate & NO & NO & YES & YES & NO \\
Change reserve ratio & NO & NO & NO & YES & NO \\
Upgrade LB borrowing cap & NO & NO & YES & YES & NO \\
View audit logs (all) & NO & YES & YES & YES & NO \\
Generate SAR report & NO & YES & YES & YES & NO \\
View reserve summary (public) & YES & YES & YES & YES & READ \\
\bottomrule
\end{tabular}
\end{table}

\begin{figure}[H]
\centering
\ThesisIncludeGraphics[width=0.9\linewidth,keepaspectratio]{Ch3_actor-permission-matrix}
\caption[Actor permission matrix (primary actions)]{Actor permission matrix (primary actions)}
\label{fig:permission-matrix}
\end{figure}



\subsection{Banking Modules}
\label{sec:banking-modules}

\subsubsection{Kinked Interest Rates}
\label{sec:kinked-rate}

Borrowing cost follows a kinked utilization model (Compound/Aave pattern~[42]) with optimal utilization $U^* = 80\%$. Below the kink the rate rises gently; above it the rate jumps to incentivise repayment (Formulas~18--19, List of Formulas). Figure~\ref{fig:kinked-rate-curve} shows the curve.

\begin{figure}[H]
\centering
\HalfWidthDiagram{fig-kinked-rate-curve}
\caption[Kinked utilization-based borrowing rate]{Kinked utilization-based borrowing rate}
\label{fig:kinked-rate-curve}
\end{figure}

\subsubsection{Liquidation Engine}
\label{sec:liquidation}

\texttt{LiquidationEngine} monitors health factor (HF) and pays third-party liquidators a 5--8\% bonus when $HF < 1.0$. Retail over-collateralised loans use $HF = (C \times LT)/D$; group pools use pool-level HF; credit-based loans trigger SBT downgrade instead of collateral seizure. Figure~\ref{fig:liquidation-engine} summarises tier models.

\begin{figure}[H]
\centering
\HalfWidthDiagram{fig-liquidation-engine}
\caption[LiquidationEngine tier models and lifecycle]{LiquidationEngine tier models and lifecycle}
\label{fig:liquidation-engine}
\end{figure}

\subsubsection{Savings and Fixed Deposits}
\label{sec:savings-vault}

\texttt{SavingsVault} accepts USDC deposits with utilization-linked yield (ERC-4626). \texttt{FixedDeposit} locks principal for a fixed term (ERC-7540). Interest splits among depositor yield, \texttt{InsuranceFund}, and protocol revenue (Formula NetInterest).

\begin{figure}[H]
\centering
\ThesisIncludeGraphics[width=0.9\linewidth,keepaspectratio]{fig-savings-vault}
\caption[SavingsVault and FixedDeposit closed loop]{SavingsVault and FixedDeposit closed loop}
\label{fig:savings-vault}
\end{figure}

\subsubsection{Credit Passport (SBT)}
\label{sec:credit-passport}

Each KYC-verified client receives a non-transferable SBT~[94] encoding repayment history (\texttt{creditScore}, \texttt{riskTier}). Table~\ref{tab:credit-tiers} governs loan limits by tier; the token upgrades on full repayment and downgrades on default.

\begin{figure}[H]
\centering
\HalfWidthDiagram{fig-credit-passport}
\caption[On-chain Credit Passport SBT lifecycle]{On-chain Credit Passport SBT lifecycle}
\label{fig:credit-passport}
\end{figure}

\begin{table}[H]
\centering
\caption[Credit Passport tier schedule]{Credit Passport tier schedule governing maximum loan limits and interest modifiers.}
\label{tab:credit-tiers}
\begin{tabular}{p{0.8cm}p{2.2cm}p{2.6cm}p{3.2cm}p{3.8cm}}
\toprule
\tableheadcolor
\textbf{Tier} & \textbf{Label} & \textbf{Score Range} & \textbf{Max Loan (USDC)} & \textbf{Interest Modifier} \\
\midrule
1 & Bronze & 0--299 & 50 & Base rate \\
2 & Silver & 300--549 & 250 & Base $-$ 0.5\% \\
3 & Gold & 550--749 & 1,000 & Base $-$ 1.0\% \\
4 & Platinum & 750--899 & 5,000 & Base $-$ 1.5\% \\
5 & Diamond & 900--1000 & 25,000 & Base $-$ 2.0\% \\
\bottomrule
\end{tabular}
\end{table}


\subsection{Multi-Entity and Cross-Tier Capital Operations}
\label{sec:multi-entity}

\begin{figure}[H]
\centering
\ThesisIncludeGraphics[width=\linewidth,height=0.78\textheight,keepaspectratio]{Ch3_multi-entity-cross-tier-operations}
\caption[Multi-entity and cross-tier capital operations]{Multi-entity and cross-tier capital operations}
\label{fig:multi-entity-ops}
\end{figure}


The architecture framing of multi-directional flows promises downward, same-tier, and upward lending plus bank-level syndication for large exposures. Earlier specifications left same-tier and upward flows at design-level only. A paper that claims a \emph{complete} banking system must specify these flows---and the genuinely multi-entity operations they enable---at the same depth as the downward flow. This section closes that gap with five concrete contract-and-flow specifications: a fully-specified InterBankLendingPool (Section~\ref{sec:interbank-pool}), upward surplus repatriation (Section~\ref{sec:upward-flow}), syndicated / club lending (Section~\ref{sec:syndicated-lending}), senior--junior tranched lending pools (Section~\ref{sec:tranched-pools}), and cross-tier treasury FX swap (Section~\ref{sec:treasury-fx}). Multilateral settlement netting (\texttt{NettingEngine}) is scoped to Future Work (Section~\ref{sec:future-work}). The corresponding ERD entities, contract list, implementation phase plan, and Future Work are updated downstream so the database design, software-engineering stages, and architecture remain a single consistent system.

\subsubsection{InterBankLendingPool: Fully Specified Same-Tier Lending}
\label{sec:interbank-pool}

One \texttt{InterBankLendingPool} (IBLP) per tier (\texttt{IBLP\_NB}, \texttt{IBLP\_LB}) supplies same-tier short-tenor (1/7/30-day) USDC liquidity among peer banks, gated by active tier role (Section~\ref{sec:defense-in-depth}).

\paragraph{Rate model.}
$r_{\text{IB}}(U)$ follows the kinked curve of Section~\ref{sec:kinked-rate} (Formulas~18--19) with kink $U^{*}_{\text{IB}} = 0.90$. The interbank rate is capped by the parent-tier downward rate:
\begin{equation}
r_{\text{IB}}(U) \le r_{\text{downward}}(U) - \delta
\end{equation}
\noindent $\delta$ is the governance-set spread that blocks arbitrage against the tier above; checked live via \texttt{getCurrentBorrowRate()} at each rate update.

\paragraph{Default cascade.}
On maturity default: (1) reserve buffer debited, (2) \texttt{InsuranceFund} drawn, (3) parent tier absorbs residual and queues role suspension via TimeLock.

\subsubsection{Upward Surplus Repatriation}
\label{sec:upward-flow}

Banks with reserves above the minimum by 20\% may deposit excess into the parent tier via \texttt{UpwardDepositFacility}, earning parent deposit yield. The upward rate is capped below the downward rate:
\begin{equation}
r_{\text{up}} < r_{\text{down}} - \delta
\end{equation}
\noindent Withdrawals revert if they breach the depositor's minimum reserve ratio; daily withdrawals are capped at 20\% of outstanding upward liabilities; net upward exposure is capped at 30\% of on-chain assets.

\subsubsection{Syndicated and Club Lending}
\label{sec:syndicated-lending}

When single-name exposure exceeds a bank's lending ceiling, multiple banks co-fund one loan through \texttt{SyndicatedLoan}: Lead Arranger, Co-Lenders, and Borrower (client via Local Bank only). Funding requires aggregate commitments $\ge 90\%$ of \texttt{totalAmount} and Co-Lender confirmation $\ge 75\%$ by capital share. Each lender's share is recorded as \texttt{shareBps[lender][syndicateId]}; principal and interest (and default recovery) distribute pro-rata by that share.

\subsubsection{Senior--Junior Tranched Lending Pools}
\label{sec:tranched-pools}

\texttt{TranchedPool} splits capital into senior (lower yield, protected) and junior (higher yield, first-loss) tranches. Repayment waterfall: senior interest $\to$ junior interest $\to$ senior principal $\to$ junior principal; losses reverse (junior written down first). Subordination ratio (junior/senior) is fixed at deposit between 10/90 and 30/70.

\subsubsection{Cross-Tier Treasury FX Swap}
\label{sec:treasury-fx}

\texttt{TreasurySwap} executes oracle-priced atomic cross-asset swaps for banking-role wallets:
\begin{equation}
\text{Amount}_{\text{to}} = \text{Amount}_{\text{from}} \times \frac{P_{\text{from}}}{P_{\text{to}}} \times (1 - \text{spread})
\end{equation}
\noindent $P$ is the Chainlink feed price; treasury spread is 5--10~bps. Post-swap reserve ratio must remain above the tier minimum; swaps above \$1M require Safe two-signature confirmation.

\subsubsection{Multilateral Settlement Netting Engine (Future Work)}
\label{sec:netting-settlement}

\noindent\textit{Future Work.} A \texttt{NettingEngine} contract to compress bilateral IBLP and \texttt{TreasurySwap} obligations into batch settlement---including partial settlement (\texttt{settlePartial}) and challenge-period dispute resolution---is listed in Section~\ref{sec:future-work}. ERD stubs (\texttt{NETTING\_BATCH}, \texttt{NETTING\_ENTRY}) remain in the schema for forward compatibility but are not specified at implementation depth in this report.

\subsubsection{Database, Phase, and Contract Consistency for Multi-Entity Operations}
\label{sec:multi-entity-consistency}

Eight ERD entities (\texttt{INTERBANK\_LOAN}, \texttt{UPWARD\_DEPOSIT}, \texttt{SYNDICATE}, \texttt{SYNDICATE\_MEMBER}, \texttt{TRANCHED\_POOL}, \texttt{TREASURY\_SWAP}, \texttt{NETTING\_BATCH}, \texttt{NETTING\_ENTRY}) extend Figure~\ref{fig:erd-extended}, together with \texttt{GROUP\_CONSENT} (Phase~II, required before \texttt{GroupLendingPool}) and \texttt{TRANCHED\_POOL\_SUBSCRIPTION} (Phase~II). All institution-participant FKs anchor to \texttt{INSTITUTION.institution\_id} (Table~\ref{tab:db-entities-ext2}). Five contracts are specified at implementation depth in this report: \texttt{InterBankLendingPool}, \texttt{UpwardDepositFacility}, \texttt{SyndicatedLoan}, \texttt{TranchedPool}, and \texttt{TreasurySwap}; \texttt{NettingEngine} is Future Work (Section~\ref{sec:future-work}). Extended entities include planned \texttt{on\_chain\_address} and \texttt{last\_synced\_block} columns for event-listener mapping.

\begin{table}[H]
\centering
\caption[FK anchors for extended and multi-entity entities]{FK anchors connecting extended entities to core schema.}
\label{tab:ext-entity-fks}
\small
\begin{tabular}{@{}p{3.5cm}p{8.5cm}@{}}
\toprule
\tableheadcolor
\textbf{Entity} & \textbf{Foreign Keys into Core Schema} \\
\midrule
SAVINGS\_ACCOUNT, FIXED\_DEPOSIT, CURRENT\_ACCOUNT & \texttt{borrower\_id} $\to$ \texttt{BORROWER}; \texttt{local\_bank\_id} $\to$ \texttt{LOCAL\_BANK(institution\_id)} \\
LOAN\_GROUP & \texttt{local\_bank\_id} $\to$ \texttt{LOCAL\_BANK(institution\_id)} \\
GROUP\_MEMBER & \texttt{group\_id} $\to$ \texttt{LOAN\_GROUP}; \texttt{borrower\_id} $\to$ \texttt{BORROWER} \\
INSURANCE\_FUND & \texttt{local\_bank\_id} $\to$ \texttt{LOCAL\_BANK(institution\_id)} \\
INTERBANK\_LOAN & \texttt{lender\_institution\_id}, \texttt{borrower\_institution\_id} $\to$ \texttt{INSTITUTION}; tier-type CHECK per IBLP instance \\
UPWARD\_DEPOSIT & \texttt{depositor\_institution\_id}, \texttt{parent\_institution\_id} $\to$ \texttt{INSTITUTION}; adjacent-tier CHECK \\
SYNDICATE & \texttt{lead\_institution\_id} $\to$ \texttt{INSTITUTION}; \texttt{borrower\_id} $\to$ \texttt{BORROWER} \\
SYNDICATE\_MEMBER & \texttt{syndicate\_id} $\to$ \texttt{SYNDICATE}; \texttt{member\_institution\_id} $\to$ \texttt{INSTITUTION} \\
TRANCHED\_POOL & \texttt{sponsor\_institution\_id} $\to$ \texttt{INSTITUTION} \\
TRANCHED\_POOL\_SUBSCRIPTION & \texttt{pool\_id} $\to$ \texttt{TRANCHED\_POOL}; \texttt{subscriber\_id} $\to$ \texttt{BANK\_USER} \\
TREASURY\_SWAP & \texttt{initiator\_institution\_id} $\to$ \texttt{INSTITUTION}; \texttt{from\_asset\_id}, \texttt{to\_asset\_id} $\to$ \texttt{ASSETS} \\
NETTING\_BATCH & \texttt{coordinator\_id} $\to$ \texttt{BANK\_USER} \\
NETTING\_ENTRY & \texttt{batch\_id} $\to$ \texttt{NETTING\_BATCH}; \texttt{participant\_institution\_id} $\to$ \texttt{INSTITUTION} \\
INSTITUTION\_MESSAGE & \texttt{sender\_institution\_id}, \texttt{receiver\_institution\_id} $\to$ \texttt{INSTITUTION} \\
LOAN\_RISK\_ASSESSMENT & \texttt{loan\_request\_id} $\to$ \texttt{LOAN\_REQUEST}; \texttt{model\_id} $\to$ \texttt{MODEL\_REGISTRY} \\
GROUP\_CONSENT & \texttt{loan\_request\_id} $\to$ \texttt{LOAN\_REQUEST}; \texttt{member\_id} $\to$ \texttt{GROUP\_MEMBER} \\
COLLATERAL\_POSITION & \texttt{borrower\_id} $\to$ \texttt{BORROWER}; \texttt{loan\_id} $\to$ \texttt{LOAN}; \texttt{asset\_id} $\to$ \texttt{ASSETS} \\
BLOCKCHAIN\_EVENT\_LOG & Referenced by: \texttt{LOAN\_REQUEST.submission\_event\_log\_id}, \texttt{approval\_event\_log\_id}; \texttt{LOAN.disbursement\_event\_log\_id}; \texttt{INSTALLMENT.payment\_event\_log\_id} \\
\bottomrule
\end{tabular}
\TableScreenshotEnd{FK Anchors for Extended and Multi-Entity Entities}
\end{table}

\subsection{System Modeling}
\label{sec:system-modeling}

This section presents the system analysis diagrams that model the platform's structure, behavior, and data flow. Loan state transitions are formalised in Section~\ref{sec:oracle-architecture}; the diagrams below show actor interactions and operational flows at the use-case and sequence level.

\subsubsection{Use Case Diagram}
\label{sec:usecase}

The system involves nine actors across eleven representative use cases (Figure~\ref{fig:usecase}). Four \textbf{core primary} actors initiate banking operations: Retail Client (A1, three KYC sub-types), Local Bank Admin (A2), National Bank Admin (A3), and World Bank Admin (A4). A5 (AI Agent) is a \textbf{planned Phase~III primary} actor with restricted write permissions gated by the human confirmation protocol (Table~\ref{tab:mvt-checklist}, item~11). Four secondary actors are external systems or authorities: Regulatory Authority (A6), Chainlink DON (A7), Blockchain Validator (A8), and External Auditor (A9). Note that the primary actor classes in the use-case diagram are a condensation of the formal nine-actor taxonomy defined in Table~\ref{tab:actor-taxonomy}: Visitor and Registered Retail User are pre-conditions for the CLIENT (BORROWER) role rather than independent diagram actors, and Local Bank Operator / Approver are merged into the Bank Approver role at the diagram level.

\begin{figure}[H]
\centering
\BalancedDiagram[width=0.98\linewidth]{Ch3_use-case-nine-actor-taxonomy}
\caption[Use-case diagram (nine-actor taxonomy)]{Use-case diagram (nine-actor taxonomy)}
\label{fig:usecase}
\end{figure}

\subsubsection{Activity Diagrams}

Figures~\ref{fig:act-lending}--\ref{fig:act-aux} present three end-to-end activity flows. Each diagram has a single start and a single end node; decision diamonds route exceptional paths (limit exceeded, rejection, verification retry) back to the terminal state without spawning parallel lifelines.

\begin{figure}[H]
\centering
\ActivityDiagram{fig-activity-lending}
\caption[USDC loan lifecycle activity flow]{USDC loan lifecycle activity flow}
\label{fig:act-lending}
\label{fig:act-loan}\label{fig:act-hierarchy}
\end{figure}

\begin{figure}[H]
\centering
\ActivityDiagram{fig-activity-onboarding-id}
\caption[Retail onboarding and identity activity flow]{Retail onboarding and identity activity flow}
\label{fig:act-onboarding}
\label{fig:act-income}\label{fig:act-profile}
\end{figure}

\begin{figure}[H]
\centering
\ActivityDiagram{fig-activity-aux}
\caption[Client banking session activity flow]{Client banking session activity flow}
\label{fig:act-aux}
\label{fig:act-chat}\label{fig:act-aichatbot}\label{fig:act-market}
\end{figure}

\paragraph{SAR activity flow.}
The Suspicious Activity Report (SAR) workflow is triggered when the Isolation Forest anomaly detector (Section~\ref{sec:threat-model}) flags a wallet with an anomaly score above the 0.75 threshold. The five-step flow is:

\begin{enumerate}[noitemsep]
\item Isolation Forest flags the wallet: \texttt{anomaly\_score > 0.75}.
\item \texttt{LOAN\_RISK\_ASSESSMENT} records the per-loan detection (FK \texttt{model\_id} $\to$ \texttt{MODEL\_REGISTRY}); \texttt{SECURITY\_EVENT\_LOG} records the system-wide AML alert event.
\item Express.js backend emits a planned event-bus message (\texttt{aml-alert} topic); the compliance dashboard consumes it.
\item Bank officer reviews the alert in the admin compliance queue:
    \begin{itemize}[noitemsep]
    \item \textbf{False positive:} Dismiss and document reason (audit trail in \texttt{audit\_logs}).
    \item \textbf{Confirmed:} Generate SAR record in \texttt{audit\_logs}; insert \texttt{INSTITUTION\_MESSAGE} (\texttt{message\_type='SAR\_ESCALATION'}) to notify the tier above (National Bank); freeze wallet via \texttt{LocalBank.freezeAccount(clientId)}.
    \end{itemize}
\item The \texttt{freezeAccount} function is gated by the \texttt{onlyApprover} modifier and is Foundry-tested (Phase~III invariant suite). A frozen wallet cannot initiate new loan disbursements or installment payments until the freeze is lifted by the tier above.
\end{enumerate}

\subsubsection{Data Flow Diagrams}

The core lending data flows are specified in Figures~\ref{fig:dfd-level0-lending} and~\ref{fig:dfd-level1-lending} (Section~\ref{sec:high-level-arch}). Level-0 defines subsystem boundary flows to the Client, Local Bank Approver, National Bank, ML Risk Scoring Service, and Blockchain Ledger. Level-1 expands Process~0 into submit/validate (1.0), risk scoring with commit--reveal oracle (2.0), approver review (3.0), disbursement (4.0), and installment processing (5.0), reading and writing stores D1--D6 (\texttt{LOAN\_REQUEST}, \texttt{LOAN}, \texttt{INSTALLMENT}, \texttt{LOAN\_RISK\_ASSESSMENT}, \texttt{BLOCKCHAIN\_EVENT\_LOG}, \texttt{BORROWING\_LIMIT}).

\subsubsection{Sequence Diagrams}

Figures~\ref{fig:seq-loan-flow}--\ref{fig:seq-chat-bot} consolidate the platform's nine interaction flows into four panels: loan approval with reject alternative, installment and income, hierarchical banking with market data and borrowing-limit, and chat with AI-chatbot.

\SequenceDiagram{fig-seq-loan-flow}
\caption[Loan approval sequence]{Loan approval sequence: application submission, ML scoring, commit--reveal oracle, approver review, and on-chain disbursement or rejection.}
\label{fig:seq-loan-flow}
\label{fig:seq-loan}\label{fig:seq-reject}
\SequenceDiagramEnd

\SequenceDiagram{fig-seq-installment-income}
\caption[Installment and income verification]{Installment payment via wallet-signed \texttt{payInstallment} with listener projection, and income-proof upload for KYC Level~2/3 officer review.}
\label{fig:seq-installment-income}
\label{fig:seq-installment}\label{fig:seq-income}
\SequenceDiagramEnd

\SequenceDiagram{fig-seq-banking-data}
\caption[Hierarchical banking and data sync]{Downward capital allocation across tiers, Chainlink price read with Redis cache, and event-driven borrowing-limit synchronisation to PostgreSQL.}
\label{fig:seq-banking-data}
\label{fig:seq-hierarchy}\label{fig:seq-marketdata}\label{fig:seq-borrowlimit}
\SequenceDiagramEnd

\SequenceDiagram{fig-seq-chat-chatbot}
\caption[Chat and AI agent sequences]{Human client--officer messaging through the Express API, and Phase~III confirmation-gated conversational agent with MCP read/write tools.}
\label{fig:seq-chat-bot}
\label{fig:seq-chat}\label{fig:seq-aichatbot}
\SequenceDiagramEnd

\subsubsection{UML Class Diagram}
\label{sec:uml-class-diagram}

Figure~\ref{fig:uml-class} is a \textbf{Software Engineering} UML class diagram for the \textit{core lending domain}. It is not another system-architecture drawing (those are Figures~\ref{fig:three-layer-arch} and~\ref{fig:component-diagram}). Instead it answers: \emph{which software classes exist, what data they hold, and how they relate}---the same question CSE470 asks when moving from requirements to implementation.

The diagram has three layers of classes:
\begin{itemize}[noitemsep]
\item \textbf{Institution hierarchy} --- \texttt{Institution} is an abstract supertype; \texttt{WorldBankReserve}, \texttt{NationalBank}, and \texttt{LocalBank} are its subtypes (generalisation), matching Figure~\ref{fig:erd-relational} and the Solidity contracts in \texttt{contracts/}.
\item \textbf{Retail lending} --- \texttt{Borrower} and \texttt{LoanRequest} mirror PostgreSQL tables; \texttt{LoanController} is the on-chain state machine that enforces the oracle gate and disbursement rules.
\item \textbf{Off-chain services} --- \texttt{LoanService} (Express.js) and \texttt{ScoringService} (FastAPI) are controller/service classes that the React UI calls; they delegate to contracts and the ML pipeline without duplicating business logic in the frontend.
\end{itemize}

\noindent\textit{Why include it?} Use-case and sequence diagrams (above) show \textit{who talks to whom over time}; the class diagram shows \textit{the static object model} that developers implement. Examiners assessing System Analysis and Design and Software Engineering coursework can trace each class back to a requirement in Table~\ref{tab:functional-requirements}.

\begin{figure}[H]
\centering
\ThesisIncludeGraphics[width=\linewidth,height=0.90\textheight,keepaspectratio]{fig-uml-class}
\caption[UML class diagram (core lending domain)]{UML class diagram (core lending domain): institution hierarchy, loan entities, on-chain controller, and off-chain services.}
\label{fig:uml-class}
\end{figure}

\subsubsection{User Interface Design}
\label{sec:ui-design}

The retail and approver interfaces follow a mobile-first Material Design~3 layout implemented in \texttt{frontend/} (React~18, TypeScript, Tailwind, MUI). Two primary surfaces are specified:

\begin{itemize}[noitemsep]
\item \textbf{Retail loan flow} --- wallet connect, KYC status banner, loan application form with plain-language APR and installment preview, and post-submission status tracking.
\item \textbf{Authority Brief dashboard} --- approver view with composite risk score $s$, SHAP waterfall chart, anomaly deviation table, and approve/decline actions gated on oracle \texttt{SCORE\_REVEALED} state.
\end{itemize}

\noindent High-fidelity Figma prototype exports (screen flows, component layout, and approver dashboard) are collected in \textbf{Appendix~F} (Section~\ref{app:ui-prototypes}) so the main chapters stay specification-focused; Figure~\ref{fig:ui-wireframe-retail} and Figure~\ref{fig:ui-wireframe-approver} will be inserted there before final submission.

\subsubsection{Four-Tier Capital Flow}

Figure~\ref{fig:four-tier} illustrates the hierarchical capital flow and cascading repayment structure.

\begin{figure}[H]
\centering
\BalancedDiagram{fig-hierarchical-banking}
\caption[Four-tier hierarchical capital flow]{Four-tier hierarchical capital flow}
\label{fig:four-tier}
\end{figure}


\subsection{Group Solidarity Lending}
\label{sec:banking-products}

Groups of 3--20 clients borrow via \texttt{GroupLendingPool} with shared collateral. Per-member share is $\text{LoanTotal}/N_{\text{members}}$. Default uses collateral first, then SBT \texttt{riskTier} downgrade~[94].

\paragraph{Over-indebtedness risk control.}
\label{sec:over-indebt}
Max two simultaneous group loans, 30-day cooling-off, and cross-bank check via \texttt{ICreditPassport.getScore()}. Reject if debt-to-income $\text{DTI} > 0.40$ (Formula in List of Formulas).

\begin{figure}[H]
\centering
\BalancedDiagram{fig-group-lending-lifecycle}
\caption[Solidarity group lending lifecycle]{Solidarity group lending lifecycle}
\label{fig:group-lending-lifecycle}
\end{figure}


\subsection{Governance Framework}
\label{sec:governance-framework}
\label{sec:governance-dual-path}
\label{sec:regulatory-compliance}

Parameter changes route through a 48-hour \texttt{TimelockController} with snapshot voting and 60\% quorum. Emergency pause uses a 2-of-3 multisig; production admin keys use Safe multisigs (Table~\ref{tab:governance-summary}). Prototype runs on public testnets; SAR workflow: Isolation Forest $\to$ compliance queue $\to$ \texttt{freezeAccount} (Section~\ref{sec:system-modeling}).

\begin{figure}[H]
\centering
\BalancedDiagram{fig-governance-dual-path}
\caption[Governance dual-path]{Governance dual-path}
\label{fig:governance-dual-path}
\end{figure}

\begin{table}[H]
\centering
\caption[Governance framework summary]{Governance framework summary.}
\label{tab:governance-summary}
\small
\begin{tabular}{@{}p{2.8cm}p{11.5cm}@{}}
\toprule
\tableheadcolor
\textbf{Area} & \textbf{Implementation} \\
\midrule
Network membership & World Bank registers National Banks; National Banks register Local Banks; hierarchical RBAC on-chain \\
Business operations & Reserve management, loan lifecycle, event notifications; testnet best-effort SLA \\
Technology & EVM contracts (decentralised); Express API (centralised); event listener syncs PostgreSQL projections \\
Emergency controls & Per-function pause registry; 2-of-3 emergency multisig; upgrade via UUPS + Safe \\
Regulatory & Audit logs via contract events; sandbox pilots for production; zkKYC Future Work \\
\bottomrule
\end{tabular}
\end{table}

\begin{figure}[H]
\centering
\BalancedDiagram{fig-sar-aml-workflow}
\caption[SAR and AML compliance workflow]{SAR and AML compliance workflow}
\label{fig:sar-aml-workflow}
\end{figure}


\subsection{Five-Layer Defense-in-Depth Security Architecture}
\label{sec:defense-in-depth}
\label{sec:contract-security}
\label{sec:fail-closed}

\begin{figure}[H]
\centering
\BalancedDiagram{fig-defense-in-depth}
\caption[Five-layer defense-in-depth security architecture]{Five-layer defense-in-depth security architecture}
\label{fig:defense-in-depth}
\end{figure}

Security spans five layers (Table~\ref{tab:defense-layers}): L5 operational keys (Safe multisigs), L4 runtime monitoring (Tenderly alerts), L3 ML fraud scoring with human review, L2 API auth and rate limits, L1 contract RBAC + \texttt{ReentrancyGuard} + Foundry fuzz tests. Each layer fails closed: if ML is unreachable, loans route to manual review rather than auto-approval.

\begin{table}[H]
\centering
\caption[Five-layer defense-in-depth model]{Five-layer defense-in-depth model.}
\label{tab:defense-layers}
\small
\begin{tabular}{@{}p{1.2cm}p{3.2cm}p{10.0cm}@{}}
\toprule
\tableheadcolor
\textbf{Layer} & \textbf{Component} & \textbf{CWB control (MVT)} \\
\midrule
L5 & Operational & Safe 2-of-3 for \texttt{WorldBankAdmin} \\
L4 & Monitoring & Tenderly alerts on large disbursement, low reserve, pause \\
L3 & AI/ML & RF + IF + SHAP; commit--reveal oracle; human approver gate \\
L2 & Application & JWT auth, rate limits, prompt-injection scan on agent input \\
L1 & Smart contract & OpenZeppelin RBAC, CEI, pause, Hardhat + Foundry tests \\
\bottomrule
\end{tabular}
\end{table}

\subsection{Threat Model and Security Controls}
\label{sec:threat-model}

Table~\ref{tab:threat-model} maps primary threats to mitigations across the five layers.

\begin{table}[H]
\centering
\caption[Threat model (summary)]{Threat model and mitigations (summary).}
\label{tab:threat-model}
\footnotesize
\begin{tabular}{@{}p{2.4cm}p{1.0cm}p{3.2cm}p{6.8cm}@{}}
\toprule
\tableheadcolor
\textbf{Threat} & \textbf{Layer} & \textbf{Vector} & \textbf{Mitigation} \\
\midrule
Reentrancy & L1 & Recursive external call & \texttt{ReentrancyGuard}, CEI, Foundry tests \\
Oracle manipulation & L1 & Bad price feed & Chainlink feeds; mock feed on testnet \\
Governance capture & L1+L5 & Privileged role abuse & Safe multisig, TimeLock, role expiry \\
JWT / API abuse & L2 & Stolen token, DDoS & Short JWT TTL, \texttt{slowapi} limits \\
ML evasion & L3 & Adversarial features & SHAP review, human approver gate \\
Admin-key theft & L5 & Compromised EOA & Safe multisig; no single-key admin \\
\bottomrule
\end{tabular}
\end{table}




\subsection{Design Decisions and Alternatives}

Table~\ref{tab:design-decisions} demonstrates the chosen stack for frontend, backend, and all dependencies for the stability and completion of the project. The comparison table is shown below:

\begin{table}[H]
\centering
\caption[Design decisions and alternatives considered]{Design decisions and alternatives considered.}
\label{tab:design-decisions}
\begin{tabular}{p{3cm}p{3cm}p{3cm}p{4.5cm}}
\toprule
\tableheadcolor
\textbf{Decision Area} & \textbf{1st Choice} & \textbf{2nd Choice} & \textbf{Key Criterion} \\
\midrule
Methodology & Agile / Scrum & Incremental & Evolving scope, milestones \\
Architecture & DApp + Off-chain AI & Hybrid with Oracle & Gas cost, ML flexibility \\
Frontend & React + TypeScript & Vue + TypeScript & Web3 ecosystem maturity \\
Smart Contract & EVM (Solidity) & Solana (Rust) & Gas, control size, free testnets \\
Fraud Detection & Random Forest & XGBoost & SHAP compatibility, simplicity \\
Anomaly Detection & Isolation Forest & Autoencoder & Unsupervised; no labelled data needed \\
XAI Method & SHAP & LIME & Theoretical guarantees, regulatory fit \\
Database & PostgreSQL & SQLite & 3NF support, async queries \\
Hosting & Vercel + Render & Localhost only & \$0 cost, publicly accessible URL \\
Network (primary) & Ethereum Sepolia (primary) and/or Polygon Amoy PoS (secondary) & Polygon zkEVM Cardona & Stability-first prototype vs.\ ZK-anchored design preference \\
Oracle mechanism & Commit-reveal relay (MVT) & Chainlink Functions (DON) & Always-available audited relay vs.\ DON-based upgrade when supported \\
Agent tool framework & MCP tool server & Direct Express.js API & Defined schema = safety boundary; Phase~III Must \\
Agent session auth & Human confirmation + standard wallet signing & EIP-7702 session keys & MVT relies on explicit consent; scoped session keys are Future Work \\
Prompt construction & Three-tier assembly & Single ad-hoc string & Phase~III Must \\
Session memory & Lineage model & 10-turn window & Phase~III Should; cross-session continuity \\
Tool exposure per turn & Named toolsets & Full 17-tool list & Phase~III Should; attack-surface reduction \\
Write-tool enforcement & Confirmation audit hook & In-pipeline check & Phase~III Must; agent write-path \\
\bottomrule
\end{tabular}
\TableScreenshotEnd{Design Decisions: Evaluated Alternatives and Selection Rationale}
\end{table}

\begin{table}[H]
\centering
\caption[Design alternative evaluation criteria]{Design alternative evaluation criteria (selected vs.\ rejected).}
\label{tab:design-alternatives-criteria}
\footnotesize
\begin{tabular}{@{}p{2.6cm}p{2.2cm}p{2.2cm}p{2.8cm}p{3.4cm}@{}}
\toprule
\tableheadcolor
\textbf{Decision} & \textbf{Selected} & \textbf{Rejected} & \textbf{Criteria (weight)} & \textbf{Outcome} \\
\midrule
Methodology & Agile + incremental phases & Pure Waterfall & Adaptability (40\%), milestone clarity (30\%), team size (30\%) & Agile wins on evolving scope \\
Settlement chain & Sepolia / Amoy (demo) & zkEVM Cardona (day-1) & Stability (50\%), cost (30\%), ZK anchor (20\%) & Demo path stability-first \\
ML delivery & Off-chain FastAPI + commit--reveal & On-chain Chainlink Functions & Latency (35\%), cost (35\%), auditability (30\%) & Off-chain for MVT \\
Fraud model & Random Forest + SHAP & XGBoost + LIME & Explainability consistency (50\%), tabular fit (30\%), imbalance (20\%) & RF exact TreeExplainer \\
Database & PostgreSQL 3NF & SQLite & Concurrency (40\%), window queries (35\%), FK integrity (25\%) & PostgreSQL for lending projections \\
\bottomrule
\end{tabular}
\end{table}

\subsubsection{Justification of Selected Technologies}

\begin{table}[H]
\centering
\caption[Selected technology justifications]{Selected technology justifications (1st choice).}
\label{tab:tech-justification-selected}
\small
\begin{tabular}{@{}p{3.2cm}p{9.8cm}@{}}
\toprule
\tableheadcolor
\textbf{Technology} & \textbf{Justification} \\
\midrule
Agile/Scrum & Supports evolving contract, UI, and ML scope within short iteration cycles. \\
DApp + off-chain AI & ML runs off-chain for speed and cost; critical lending decisions remain on-chain. \\
React + TypeScript & Strong Web3 library support (Wagmi, RainbowKit, Viem, ethers.js). \\
EVM (Solidity) & Mature security tooling; Sepolia/Amoy testnets; OpenZeppelin libraries. \\
MUI (Material Design 3) & Ready-made tables, forms, and dashboards for banking UI. \\
Random Forest & Exact SHAP TreeExplainer; strong on tabular lending data; handles class imbalance. \\
Isolation Forest & Unsupervised anomaly detection with few labeled fraud samples. \\
SHAP & Consistent, theory-grounded feature attribution for explainability. \\
PostgreSQL & ACID compliance, window queries for borrowing limits, JSON for prototyping. \\
Vercel + Render & Free hosting for live supervisor and examiner demos. \\
Sepolia / Amoy testnets & Stability-first deployment path within the thesis timeline. \\
Commit--reveal oracle & Simple MVT path with immutable on-chain ML score audit trail. \\
MCP tool server & Typed tool schema bounds agent reads and writes. \\
EIP-7702 session keys & Future Work: scoped session keys for agent on-chain actions. \\
\bottomrule
\end{tabular}
\end{table}

\begin{table}[H]
\centering
\caption[Alternative technology justifications]{Retained alternatives (2nd choice) and why they were not selected.}
\label{tab:tech-justification-alternatives}
\small
\begin{tabular}{@{}p{3.2cm}p{9.8cm}@{}}
\toprule
\tableheadcolor
\textbf{Alternative} & \textbf{Why retained but not primary} \\
\midrule
Waterfall & Clear phases, but too rigid for research prototyping. \\
Hybrid on-chain oracle & Higher latency and cost; external oracle dependency. \\
Vue + TypeScript & Smaller Web3 ecosystem than React. \\
Solana (Rust) & Steeper learning curve and weaker audit tooling within 16 weeks. \\
Tailwind CSS & More custom UI work than MUI for banking layouts. \\
XGBoost & Approximate SHAP; weaker explanation consistency. \\
Autoencoder & Needs more data and tuning than Isolation Forest. \\
LIME & Less stable explanations than SHAP for audit use. \\
SQLite & Limited concurrency and window-query support. \\
Localhost-only deploy & Blocks remote supervisor and examiner access. \\
Polygon Amoy PoS (2nd network) & Fallback testnet if Sepolia is unstable. \\
Chainlink Functions & Stretch upgrade when supported; higher latency and fees. \\
Direct Express.js API & Weaker safety boundary than MCP tool schema. \\
ERC-4337 paymasters & Weaker scope control than EIP-7702 for agent actions. \\
\bottomrule
\end{tabular}
\end{table}

\subsection{Planned AI/ML Support and Risk-Score Wiring}
\label{sec:aiml-support}

The purpose of the AI/ML layer is to provide a risk score per loan approval that will be delivered on-chain through the commit--reveal oracle (Section~\ref{sec:oracle-architecture}). The oracle gate design has been shown in Figure~\ref{fig:ml-oracle-commit-reveal}. ML scoring and explainability outputs have been shown in Section~\ref{sec:ml-explainability}, and the rest of the benchmarks and evaluations have been reported in Chapter~\ref{ch:evaluation} (Sections~\ref{sec:performance-evaluation} and~\ref{sec:statistical-analysis}).

\textbf{Risk Scoring:} A sequential step combines Random Forest classifier and Isolation Forest to provide fraud probability and anomaly scores, and these scores are evaluated within $s \in [0,1]$.

\textbf{Commit--reveal oracle:} This commits a hash of $(s, \text{nonce})$ to \texttt{LoanController}, then each transaction comes with a risk score.

\textbf{Decision threshold:} A minimum and maximum threshold of the value of score $(s)$ is maintained for loan approval decisions: approval consideration ($s < 0.4$), manual review ($0.4 \le s < 0.7$), or rejection ($s \ge 0.7$).

\begin{figure}[H]
\centering
\BalancedDiagram{Ch4_ml-oracle-commit-reveal}
\caption[ML oracle commit--reveal gate]{ML oracle commit--reveal gate}
\label{fig:ml-oracle-commit-reveal}
\end{figure}



\subsection{Conversational Agent (Phase III--IV)}
\label{sec:optional-agent}

The conversational agent that has been designed to assist the clients to better navigate the platform is planned to be completed by Phase~III. It is an auxiliary feature that works alongside the standard web UI for those who may not be familiar with the banking system of the project and have less domain knowledge and want to understand the platform and how to navigate with it. The schema of the full MCP agent has been shown in Table~\ref{tab:mcp-tools}; explainability for approvers is documented in Section~\ref{sec:ml-explainability}.

\paragraph{Agent six-step pipeline.} Clients are expected to interact with the frontend UI like any other typical user-friendly banking website. The agent works as a guiding assistant and its tool-calling abilities use the same Express.js API endpoints; the same smart contracts are considered for functionalities. This six-step pipeline lets the Qwen3-8B model have access to the live on-chain context injection layer that helps the agent process related instant and consistent information that is relevant to how a client is interacting with the platform. The layers have been described shortly below:

\begin{enumerate}[noitemsep]
\item \textbf{User enters the platform and interacts with the agent:} frontend interaction is recorded and forwarded via the Express.js backend over Server-Sent Events (SSE).
\item \textbf{Agent model with live on-chain context:} The Express.js context injection layer provides the client's information to the model (credit score, active installments, loan status, SBT risk tier, pool utilisation). The model receives this information about the user before the conversation prompt by the user, to process relevant information about the client to assist them.
\item \textbf{Q\&A:} For this part of the conversation between the agent and client, the agent only processes relevant information about the client via RAG (ChromaDB) and maintains formal conversation with the client, and does not call any tools.
\item \textbf{Action request:} After a formal conversation with the client, if the client requests a state-modifying action such as requesting a loan, the agent will choose and prepare MCP write tools, assemble the full parameter set from the injected context, and show a concise summary of action on behalf of the client. Ex: \emph{``You are requesting a loan of 150~USDC for 6~months at 8.5\% APR. Your monthly installment would be 26.25~USDC. Shall I proceed?''}
\item \textbf{Human confirmation gate:} The agent will wait until a confirmation of affirmation for the action that has been prepared to initiate the loan process. The MCP agent writing tools will not be initiated until a positive response has been received from the client side. If a client responds with vague actions three times, the agent does not make the tool calls, pauses the agent service for 10~minutes, and logs this interaction to \texttt{AGENT\_ACTION\_LOG}.
\item \textbf{MCP write tool functionality:} If a successful conversation is established with the client, the agent takes confirmation from the client to initiate the loan request process; with a positive response the agent initiates write tools and accesses the Express.js banking API, then checks for an EIP-7702 session key---if active, it signs the on-chain transaction approval request. The agent gives a confirmation that the request has been forwarded and notifies the client that the loan application has been submitted, shows a reference number, and how long the approval might take.
\end{enumerate}

\paragraph{MCP tool server.} For the MVT scope the agent has been designed for 3--5 tools for a minimal typed tool schema. The specified 17 tools have been designed for future work, and some of these tools may be included in this prototype build if they can be delivered in the allocated time, after the MVT scope has been secured.

\begin{table}[H]
\centering
\caption[MCP tool server (MVT subset)]{MCP tool server (MVT subset). The full 17-tool suite is Future Work.}
\label{tab:mcp-tools}
\begin{tabular}{p{4.8cm}p{1.4cm}p{6.3cm}}
\toprule
\tableheadcolor
\textbf{Tool} & \textbf{Type} & \textbf{Maps to / Description} \\
\midrule
\texttt{get\_account\_state}       & READ & SBT tier, loan count, savings balance \\
\texttt{get\_loan\_status}         & READ & Active loans + next installment \\
\texttt{get\_requirements}         & READ & Documents + KYC needed for a given loan amount \\
\texttt{submit\_loan\_application} & WRITE & POST /api/loan/apply \\
\texttt{pay\_installment}          & WRITE & POST /api/installment/pay \\
\bottomrule
\end{tabular}
\end{table}

\subsubsection{In-product assistant: local large language model (LLM) integration (prototype)}
\label{sec:local-llm}

\paragraph{Purpose.}
The in-product assistant is a \textbf{planned Phase~III Must deliverable} (Section~\ref{sec:optional-agent}), providing a plain-language interface alongside the React dashboard. It combines a locally hosted, privacy-preserving large language model (Qwen3-8B, Apache~2.0~licence) with a Model Context Protocol (MCP) tool server exposing a minimal MVT subset of \textbf{3--5 tools} (Section~\ref{sec:optional-agent}). The full \textbf{17-tool} suite (9 read tools, including a session history search tool, and 8 write tools requiring human confirmation) remains a specified Future Work expansion rather than the graded prototype baseline. The agent can both answer questions---via retrieval-augmented generation (RAG) from policy documents---and execute on-chain banking operations via the MVT write-tool subset, subject to an explicit human confirmation gate before any state-modifying action is taken. Local hosting keeps conversation data on institutional infrastructure, supporting data-sovereignty requirements in regulated pilots. A client who prefers the React UI experiences no reduction in available functionality, since every operation the agent can perform is equally available through the standard dashboard.

\paragraph{Model and runtime (local development).}
The inference endpoint follows the \textbf{OpenAI Chat Completions} contract (\texttt{/v1/chat/completions}) with \texttt{stream: true}, proxied from the project backend to a local service such as \texttt{http://127.0.0.1:1234}. During early development, the assistant used a locally hosted \textbf{Gemma-4-E4B} checkpoint (Google mixture-of-experts model, $\sim$26B total / 4B active parameters) via LM Studio in GGUF form (e.g., \texttt{Q4\_K\_M} quantization) for a laptop-friendly footprint. For the final thesis evaluation, the assistant is replaced by \textbf{Qwen3-8B} (Alibaba Cloud Qwen team, released April~28~2025, Apache~2.0~licence)~[67], a dense 8.19B-parameter causal language model with a 32K-token native context window (extendable to 131K tokens via YaRN scaling) and support for 100+ languages and dialects. A distinguishing architectural feature of Qwen3-8B is its \emph{switchable reasoning mode}: the model operates in \textbf{thinking mode} (chain-of-thought wrapped in \texttt{<think>} tokens, suited to multi-step compliance or policy questions) and \textbf{non-thinking mode} (direct, low-latency responses for general navigation queries), selectable per-turn with \texttt{/think} or \texttt{/no\_think} in the system prompt. For Q\&A and retrieval-augmented generation (RAG) queries, non-thinking mode is the default; thinking mode is activated only for the red-team regulatory-compliance prompts in the evaluation protocol of Section~\ref{sec:optional-agent}. For agent tool-calling, non-thinking mode is the default for read tools and confirmation summaries; thinking mode is activated for complex multi-step operations such as EMI calculation or group signature coordination. The per-user context namespace (see Section~\ref{sec:aiml-support}) is prepended to every system prompt as a structured JSON block, so the model has full account context before reading the user's message. The API model identifier is environment-driven; the integration layer remains identical for both models.

\paragraph{End-to-end behavior.}
The user interface assembles a short transcript (user/assistant messages) and posts it to the backend streaming route. Before the user message reaches the model, the context assembly layer: (a)~selects the active toolset based on intent classification, (b)~applies prompt injection scanning to all user-sourced fields, and (c)~assembles the three-tier system prompt (Stable + Context + Volatile), injecting the compression summary from the parent session into the Volatile tier if this is a continuation session. The active toolset name and any injection scan flag are recorded in \texttt{AGENT\_ACTION\_LOG} for every turn. The backend then opens a streaming request to the local model server, converts upstream token events into \textbf{server-sent events (SSE)} for the browser, and the UI incrementally appends text. The message renderer uses Markdown (including GitHub-flavored extensions), math (KaTeX) for \texttt{\$\$...\$\$} blocks, and line-break rules appropriate for chat transcripts.

For write-tool requests, the agent assembles the full parameter set from the user's request and the injected on-chain state, then presents a confirmation summary. Only after the user sends explicit affirmative consent does the agent call the corresponding MCP write tool via the Express.js banking API layer. The EIP-7702 session key (if active) signs the resulting on-chain transaction within its approved scope. Every executed write tool is logged to \texttt{agent\_action\_log} with the confirmation message ID as the audit reference.

\paragraph{Live on-chain context injection.}
The Express.js backend injects the client's full on-chain state as a structured JSON context block into every system prompt before forwarding the user's message to the model. The per-user context namespace is:
\begin{verbatim}
{
  "tier": "volatile",
  "borrower_id": "550e8400-e29b-41d4-a716-446655440000",
  "wallet_address": "0xABC...",
  "language": "en",
  "credit_tier": "Silver",
  "credit_score": 512,
  "active_loans": [
    {
      "loan_id": "L-4821",
      "amount": 100,
      "next_due": "2026-06-15",
      "installment": 17.5
    }
  ],
  "savings_balance": 250.0,
  "pool_utilisation": 0.67,
  "kyc_tier": 1,
  "session_id": "sess_20260602_xyz",
  "parent_session_id": "sess_20260515_abc",
  "compression_summary": "Client asked about Gold tier eligibility.
  Was told 2 more on-time repayments required. No action taken.",
  "conversation_history": [ "...last 10 turns..." ]
}
\end{verbatim}
The \texttt{compression\_summary} field is injected only when this is a child session (i.e., when \texttt{parent\_session\_id} is non-null). It provides the prior-session context without exceeding the Volatile tier token budget. This context block is injected only when the user is authenticated and the \texttt{eth\_call} round-trip completes within a 200~ms timeout (falling back to a static-context response otherwise). The model uses the injected context to ground its answers in the user's actual account data rather than generic policy text, materially reducing hallucination risk for account-specific queries.


\paragraph{Security and limitations (prototype stance).}
The agent's safety posture is enforced at four levels. (1)~\textbf{Tool-schema boundary:} the agent interacts with CWB exclusively through the MVT MCP tool schema (\textbf{3--5 tools}) and cannot execute arbitrary code, make unapproved API calls, or read data outside that schema. The full 17-tool suite remains a specified but undeployed Future Work expansion. (2)~\textbf{Human confirmation gate:} every write tool requires explicit affirmative consent in the conversation history; the confirmation turn ID is logged as the audit reference. (3)~\textbf{Session key scope:} the EIP-7702 session key is scoped to specific tools, value-capped, time-bound to 24~hours, and revocable at any time. A session that attempts to call an out-of-scope tool is rejected at the key level, not just at the application level. (4)~\textbf{Prompt injection scanning and lifecycle hook middleware:} user-sourced content is scanned for injection patterns before entering the prompt (Layer~2 control, Section~\ref{sec:aiml-support}), and every write-tool HTTP POST is intercepted by a middleware stack that enforces the confirmation invariant at the application layer independently of the model's output. Together, these four levels form a defence-in-depth posture in which no single point of failure can circumvent a critical safety constraint. Hallucination risk is bounded by: (a) the injected on-chain state grounding account-specific answers, and (b) the refusal layer for regulatory and legal queries (100\% target on the red-team set, Section~\ref{sec:optional-agent}).
\paragraph{Architecture diagram.}
Figure~\ref{fig:local-llm-mermaid} shows the compact end-to-end request path; Figure~\ref{fig:local-llm-tikz} expands the same flow with component boundaries and authentication context.

\begin{figure}[H]
\centering
\OnePageDiagram{fig-local-llm-compact}
\caption[AI banking agent request path]{AI banking agent request path}
\label{fig:local-llm-mermaid}
\end{figure}

\begin{figure}[H]
\centering
\OnePageDiagram{fig-local-llm}
\caption[AI agent data flow]{AI agent data flow}
\label{fig:local-llm-tikz}
\end{figure}

\begin{table}[H]
\centering
\caption[Technology stack summary]{Technology stack summary.}
\label{tab:tech-stack}
\begin{tabular}{p{3cm}p{4cm}p{6.5cm}}
\toprule
\tableheadcolor
\textbf{Layer} & \textbf{Technology} & \textbf{Version / Notes} \\
\midrule
Smart Contract & Solidity, OpenZeppelin & 0.8.20; Ownable, ReentrancyGuard \\
Frontend & React, TypeScript & Material UI; Design 3 \\
Wallet Integration & Wagmi, RainbowKit, Viem & EIP-1193 compliant \\
Build and Test & Hardhat & Automated test suite; deployment scripts \\
Backend API & Express.js, Node.js & REST API; middleware architecture \\
Database & PostgreSQL (designed) & 51 entities (31 core + 6 extension + 14 stubs), 3NF; relational integrity \\
Target Network & Ethereum Sepolia (primary) and/or Polygon Amoy PoS (secondary single-chain prototype) & Stable public testnets for delivery; Polygon zkEVM retained as evaluated design preference \\
AI Agent Engine & Qwen3-8B (llama.cpp, Q4\_K\_M) & MCP tool server; per-user context namespace; human confirmation gate \\
MCP Tool Server & Minimal MVT tool subset (3--5 tools); prompt injection scanner; confirmation audit hook middleware & Exposes a small read+write surface to the agent with typed parameter schemas \\
Session Keys & EIP-7702 (Future Work) & Scoped, time-bound, value-capped agent wallet authorisation \\
Oracle Network & Commit--reveal (MVT) + Chainlink Functions (stretch) & Commit--reveal enforces auditable score commitment; DON upgrade when supported \\
Price Feeds & Chainlink AggregatorV3Interface & Fiat/USD pairs (e.g.\ EUR/USD) and ETH/USD; 8-decimal precision \\
Automation & Chainlink Automation & Trustless installment due-date checking; replaces centralised cron job \\
Proof of Reserve & Chainlink PoR & Cryptographically verifiable WorldBankReserve balance \\
Event Indexing & PostgreSQL event listener & Stability-first indexing baseline (GraphQL over The Graph optional later) \\
\bottomrule
\end{tabular}
\TableScreenshotEnd{Local LLM Assistant Integration: Components, Configuration, and Prototype Stance}
\end{table}



\subsection{Real-Time Dashboard Pipeline and Runtime Monitoring}
\label{sec:realtime}

For stability, lightweight data indexing strategies are used to reduce complexities. After the correct implementation of the core contracts and event schemas, an optional query layer is to be added for GraphQL access for better demonstration of analytics. A Node.js service is utilized for WebSocket event listener via Socket.IO. Tenderly monitoring and MVT-graded time and attempt triggers are added for monitoring and inconsistency checks.



\subsubsection{Functional Verification of Design Alternatives}
\label{sec:design-alternative-verification}
% Reserved for Phase~IV Foundry simulation and gate-demo results; see Section~\ref{sec:impl-phase-plan}.




%% CHAPTER 5: EVALUATION, ANALYSIS, AND FINAL DESIGN
\chapter{Evaluation, Analysis, and Final Design}
\label{ch:evaluation}
\label{ch:market-analysis}

\section{Performance Evaluation}
\label{sec:performance-evaluation}
\label{sec:eval-method}
\label{sec:verification-validation}
\label{sec:foundry}

\begin{itemize}[noitemsep]
\item \textbf{Verification:} Hardhat contract tests (\texttt{test/}, 42 cases); Vitest API tests (\texttt{backend/tests/}); Foundry invariants (solvency, role segregation, capital-flow direction) in Phase~IV.
\item \textbf{Validation:} Phase gates G1--G4; BCCC fraud benchmarks (Table~\ref{tab:ml-bccc-results}, Figure~\ref{fig:ml-eval}).
\end{itemize}

\label{sec:implementation-evidence}

\begin{table}[H]
\centering
\caption[Implementation evidence]{Implementation evidence (repository root).}
\label{tab:implementation-evidence}
\small
\begin{tabular}{@{}p{3.2cm}p{4.4cm}p{6.6cm}@{}}
\toprule
\tableheadcolor
\textbf{Component} & \textbf{Path} & \textbf{Status} \\
\midrule
Smart contracts & \texttt{contracts/*.sol} & World/National/Local Bank compilable \\
Contract tests & \texttt{test/*.test.ts} & 42 Hardhat cases (four suites) \\
Backend API & \texttt{backend/src/} & Loans, banks, auth, income routes \\
API tests & \texttt{backend/tests/} & Vitest auth, banks, loans, chat \\
Frontend & \texttt{frontend/src/} & React + Wagmi wallet shell \\
ML pipeline & \texttt{Documentation/.../ML pipeline - Mac M4/} & BCCC subsample; \texttt{metrics.json} \\
ML service & \texttt{ml-service/app/} & FastAPI \texttt{/health}, \texttt{/score} scaffold \\
Deploy & \texttt{scripts/deploy.ts} & Sepolia/Amoy/local ready \\
\bottomrule
\end{tabular}
\end{table}

\section{Analysis of Design Solutions}
\label{sec:analysis-design-solutions}

\begin{table}[H]
\centering
\caption[Design selection vs.\ CO6 factors]{Selected design vs.\ CO6 evaluation factors.}
\label{tab:co6-selection-summary}
\small
\begin{tabular}{@{}p{2.2cm}p{3.6cm}p{4.0cm}p{4.4cm}@{}}
\toprule
\tableheadcolor
\textbf{Factor} & \textbf{Selected} & \textbf{Evidence} & \textbf{Status} \\
\midrule
Cost & \$0 prototype; L2 gas ${\approx}$\$0.011/loan & Tables~\ref{tab:economic-feasibility},~\ref{tab:gas-cost} & Met \\
Efficiency & Off-chain ML; single-chain & NFR-01; Phase~IV gas table & ML trained; live p95 pending \\
Usability & React + wallet; MCP agent & NFR-05; FR-08 & UI shell; agent Phase~III \\
Impact & Tier-transparent flow; group lending & Section~\ref{sec:societal-impact} & Specified \\
Sustainability & PoS testnets; off-chain ML & Section~\ref{sec:environmental-impact} & Qualitative \\
Maintainability & Monorepo; OpenZeppelin; G1--G4 & Table~\ref{tab:implementation-evidence} & 42 tests passing \\
Performance & RF+IF+stacking; F1\,=\,0.891 & Table~\ref{tab:ml-bccc-results} & Benchmark met; E2E Phase~II--IV \\
\bottomrule
\end{tabular}
\end{table}

\subsection{Requirements Traceability Matrix}
\label{sec:traceability}

\begin{table}[H]
\centering
\caption[Requirements traceability (excerpt)]{Requirements traceability matrix (excerpt).}
\label{tab:traceability}
\footnotesize
\begin{tabular}{@{}p{1.1cm}p{2.6cm}p{2.8cm}p{2.6cm}p{3.2cm}@{}}
\toprule
\tableheadcolor
\textbf{Req.} & \textbf{Use case / diagram} & \textbf{Module} & \textbf{Test / evidence} & \textbf{MVT} \\
\midrule
FR-01 & Figure~\ref{fig:four-tier} & \texttt{WorldBankReserve.sol} & \texttt{hierarchy.test.ts} & \#1 \\
FR-02 & Figures~\ref{fig:act-lending},~\ref{fig:seq-loan-flow} & \texttt{LocalBank.sol}, loans API & \texttt{LocalBank.test.ts} & \#2--3 \\
FR-04 & Table~\ref{tab:integrity-constraints} & \texttt{BORROWING\_LIMIT} + SBT & Foundry invariant (planned) & \#5 \\
FR-05 & Figure~\ref{fig:ml-oracle-commit-reveal} & \texttt{ml-service/} & Table~\ref{tab:ml-bccc-results} & \#4, \#8 \\
FR-07 & Figures~\ref{fig:dfd-level0-lending}--\ref{fig:dfd-level1-lending} & Event listener & PostgreSQL tests & DT-I.12--13 \\
FR-08 & Figure~\ref{fig:seq-chat-bot} & MCP agent & Section~\ref{sec:optional-agent} & \#11--12 \\
NFR-02 & Table~\ref{tab:defense-layers} & OpenZeppelin RBAC & \texttt{WorldBankReserve.test.ts} & \#13--14 \\
\bottomrule
\end{tabular}
\end{table}

\subsection{Final Design Adjustments}
\label{sec:final-design-adjustments}
\label{sec:bridge}

Multi-chain bridging (Sepolia $\leftrightarrow$ zkEVM) was rejected due to bridge attack surface. \textbf{Final design:} single-chain deployment on Sepolia/Amoy; Polygon zkEVM Cardona remains the design-preference target without a live cross-chain bridge. Details in Tables~\ref{tab:design-decisions},~\ref{tab:design-alternatives-criteria} (Chapter~\ref{ch:design-alternatives}).

\section{Statistical Analysis}
\label{sec:statistical-analysis}

\subsection{Datasets and Benchmark Results}
\label{sec:ml-workload}
\label{sec:ml-explainability}
\label{sec:ml-evaluation}

\begin{table}[H]
\centering
\caption[BCCC fraud-detection benchmark]{BCCC-DeFiFraudTrans-2025 subsample benchmark (address-grouped hold-out; \texttt{evidence/mac-m4/metrics.json}). Primary dataset; Elliptic~[68] secondary only; on-chain CWB log for future retrain.}
\label{tab:ml-bccc-results}
\begin{tabular}{lrrrrr}
\toprule
\tableheadcolor
\textbf{Split} & \textbf{Precision} & \textbf{Recall} & \textbf{F1} & \textbf{ROC-AUC} & \textbf{$n_{\text{test}}$} \\
\midrule
Held-out test & 0.874 & 0.908 & 0.891 & 0.991 & 5{,}140 \\
\bottomrule
\end{tabular}

\vspace{4pt}
\footnotesize TN\,=\,4{,}621, FP\,=\,60, FN\,=\,42, TP\,=\,417; fraud rate ${\approx}8.7\%$; 21 features; leakage column \texttt{flag} removed.
\end{table}

\begin{figure}[H]
\centering
\small
\begin{minipage}[t]{0.48\linewidth}
\centering
\Diagram[width=\linewidth]{fig-ml-metrics-bars}
\end{minipage}\hfill
\begin{minipage}[t]{0.48\linewidth}
\centering
\Diagram[width=\linewidth]{fig-ml-confusion-matrix}
\end{minipage}

\vspace{6pt}
\Diagram[width=0.92\linewidth]{fig-ml-shap-importance}

\caption[ML evaluation and SHAP]{ML evaluation (F1\,=\,0.891, ROC-AUC\,=\,0.991) and SHAP feature importance.}
\label{fig:ml-eval}
\end{figure}

\section{Comparisons and Relationships}
\label{sec:comparisons-relationships}

\begin{table}[H]
\centering
\caption[Selected vs.\ rejected design choices]{Selected vs.\ rejected design choices and linked metrics.}
\label{tab:comparison-relationships}
\footnotesize
\begin{tabular}{@{}p{2.4cm}p{2.8cm}p{2.8cm}p{5.2cm}@{}}
\toprule
\tableheadcolor
\textbf{Dimension} & \textbf{CWB (selected)} & \textbf{Rejected / baseline} & \textbf{Metric / requirement} \\
\midrule
Architecture & Four-tier hierarchy & Flat Aave/Compound & FR-01/FR-03; Table~\ref{tab:protocol-comparison} \\
Settlement & Single-chain Sepolia/Amoy & Multi-chain bridge & Table~\ref{tab:gas-cost}; Section~\ref{sec:final-design-adjustments} \\
ML delivery & Off-chain RF+IF+SHAP & On-chain-only scoring & F1\,=\,0.891; NFR-01 \\
Oracle & Commit--reveal relay & Chainlink Functions DON & Human approver gate (FR-02) \\
Gas & L2 Cardona (design) & Ethereum mainnet & ${\approx}$\$0.011 vs \$2.40--\$4.80 \\
Human vs ML & Decision support & Auto-disburse & NFR-02; Section~\ref{sec:ethical-issues} \\
\bottomrule
\end{tabular}
\end{table}

\section{Discussion}
\label{sec:discussion}
\label{sec:research-limitations}

\begin{table}[H]
\centering
\caption[Research limitations (summary)]{Research limitations and mitigations (summary).}
\label{tab:research-limitations}
\small
\begin{tabular}{@{}p{1.0cm}p{3.2cm}p{8.0cm}@{}}
\toprule
\tableheadcolor
\textbf{ID} & \textbf{Limitation} & \textbf{Mitigation / scope} \\
\midrule
G1 & Partial DeFi precedents (Maple, Goldfinch) & Claim: combined four-tier + ML + group lending \\
G2 & BCCC flat-pool training $\neq$ hierarchical CWB flows & Baseline F1 only; CWB retrain after testnet traffic \\
G3 & Pre-thesis evidence = spec + benchmarks & MVT testnet E2E and Foundry suite (Phase~IV) \\
G4 & On-chain group lending $\neq$ social solidarity & No BRAC/Grameen repayment claims \\
G5 & Oracle collusion; model staleness & Commit--reveal MVP; Certora proofs in Future Work \\
G6 & Technical default only & No sovereign preferred-creditor enforcement \\
\bottomrule
\end{tabular}
\end{table}



%% CHAPTER 6: CONCLUSION
\chapter{Conclusion}
\label{ch:conclusion}

\section{Conclusion}
\label{sec:conclusion-summary}
\label{sec:research-contribution}

This paper is a complete specification of all the core sections that will be used in development of the project. It introduces a problem--solution view of how decentralized financing can be made adaptable to the institutional layers of our traditional financing system that is not flat, rather layered. Flat DeFi is hard to be made adaptable among the general mass, so our system provides a tier-based architecture for financial flow, with the added benefit of transparency with programmable smart contracts for the betterment of the client, and less control for the capitalistic organizations. The paper shows credible studies from the past that have contributed to the field before, and uses that evidence and those architectures as proof that a tier-based DeFi system can be built with better security layers that can contribute to the development finance sector globally. The core architecture stack follows mostly a traditional DeFi financial platform build, but the AI/ML layer integration and the tier-based layers are completely novel; no other paper or project shows all four-layer integration before. This project idea is enhanced by the existing systems that have two-tier-based banking systems.

A detailed plan has been introduced for how new clients can be onboarded and attracted to the platform, how clients will interact with the system, how banking authorities can approve loans, and how ML and AI layers can communicate in the middle with both clients and authorities to make this whole process of lending decentralized cryptocurrencies easier and more secure. The potential of such a decentralized system to be popular and adaptable is huge if implemented correctly and if the higher authorities show interest and invest heavily in this idea for the good of humanity. The purpose of this project is to push the idea of decentralized currencies forward for future generations to show more interest, and to support the building blocks for better financing systems in the future.

\section{Future Work}
\label{sec:future-work}

This section has been specified in the above chapters; it focuses on the potential of the project and the technical abilities and functional capabilities that can be improved and enhanced in the future. First priority is completing the MVT checklist of the development cycle for the defense of the project, and the rest are mostly added stretch to show how many features can be incorporated in the given time. A small summary of the future work:

Deployment on Polygon zkEVM Cardona, which remains experimental for production grade currently; as zkEVM may prove to be more stable, the platform will be deployed on Polygon zkEVM as it is more cost efficient for transactional activities.

Cross-tier features such as \texttt{InterBankLendingPool} and group lending, upward deposit facilities will be improved and made functional. For security and monitoring, multisig admin for World Bank operations and a live reserve transparency dashboard will be integrated. zkKYC (currently unstable globally), Certora proofs, federated learning, and GraphSAGE (quite complex to implement now) have been specified in the paper and will be more attainable with more time and effort; these have been moved to future work.

This field requires much research, as financial flow is heavily manipulated and misused globally by capitalists; transparency in capital flow is vital for equal access to resources and opportunity to live in harmony. The long-term goal of this project is to evolve this developed system to a more robust level, with more proper and efficient architecture and system design, a regulation-aware system while keeping transparency, equal access, and beneficial to humanity as core design principles.


\appendix

\chapter{Database Schema Reference}
\label{app:db-schema-reference}

This appendix records the physical-schema details moved from Chapter~\ref{ch:design-alternatives} to keep the main design chapter focused on conceptual design.

\section{Full Relational ERD}
\label{sec:appendix-erd-full}

Figure~\ref{fig:erd-relational} is the complete logical data model: 51 documented entities plus two Phase-IV configuration tables (\texttt{CHAIN\_REGISTRY}, \texttt{SESSION\_KEY\_PERMISSION}), with crow's-foot cardinalities, migration-tier colour coding, and a notation key. It is placed here rather than in Chapter~\ref{ch:design-alternatives} because the diagram requires a full landscape page at high magnification to remain legible.

\LandscapeDiagram{ERD_diagram_relational}
\caption[Full relational ERD (51 entities)]{Full relational ERD: 51 documented entities plus two Phase-IV configuration tables (\texttt{CHAIN\_REGISTRY}, \texttt{SESSION\_KEY\_PERMISSION}), grouped by migration tier (M1 core, M2, M3, stub, future work).}
\label{fig:erd-relational}
\LandscapeDiagramEnd

\section{Indexing Strategy}
\label{sec:indexing-strategy}

B-tree indexes (the PostgreSQL default) support efficient retrieval on frequently queried columns:

\begin{table}[H]
\centering
\caption[Indexing strategy]{Indexing strategy: B-tree indexes for time-window and high-frequency queries.}
\label{tab:indexing-strategy}
\begingroup
\footnotesize
\setlength{\tabcolsep}{4pt}
\renewcommand{\arraystretch}{1.2}
\begin{tabular}{@{}%
  >{\raggedright\arraybackslash}p{0.19\linewidth}%
  >{\raggedright\arraybackslash}p{0.41\linewidth}%
  >{\raggedright\arraybackslash}p{0.36\linewidth}@{}}
\toprule
\tableheadcolor
\textbf{Index} & \textbf{Table / Column(s)} & \textbf{Rationale} \\
\midrule
\texttt{idx\_loan\_borrower} &
  \texttt{LOAN\_REQUEST}(\texttt{borrower\_id}) &
  High-frequency lookup for borrower loan history. \\
\texttt{idx\_loan\_status} &
  \texttt{LOAN\_REQUEST}(\texttt{status}, \texttt{oracle\_state}) &
  Efficient filtering of application lifecycle and oracle pipeline states. \\
\texttt{idx\_installment\_due} &
  \texttt{INSTALLMENT}(\texttt{due\_date}, \texttt{status})\newline
  \texttt{WHERE status = 'PENDING'} &
  Partial-index range query for pending overdue installment detection. \\
\texttt{idx\_txn\_created} &
  \texttt{TRANSACTION}(\texttt{created\_at}) &
  Time-window aggregation for 6-month and 1-year borrowing limits. \\
\texttt{idx\_txn\_borrower} &
  \texttt{TRANSACTION}(\texttt{borrower\_id}) &
  Per-borrower transaction history retrieval. \\
\texttt{idx\_market\_symbol} &
  \texttt{MARKET\_DATA}(\texttt{symbol},\newline
  \texttt{recorded\_at}) &
  Composite index for price feed time-series queries. \\
\texttt{idx\_loan\_bank\_status} &
  \texttt{LOAN}(\texttt{local\_bank\_id}, \texttt{status},\newline
  \texttt{created\_at DESC}) &
  Loan lifecycle queries by bank and status. \\
\texttt{idx\_loan\_borrower\_active} &
  \texttt{LOAN}(\texttt{borrower\_id}, \texttt{status})\newline
  \texttt{WHERE status = 'ACTIVE'} &
  Partial index for per-borrower open-loan count. \\
\texttt{idx\_agent\_session\_date} &
  \texttt{AGENT\_ACTION\_LOG}(\texttt{session\_id},\newline
  \texttt{created\_at DESC}) &
  Per-session agent action history retrieval (join to \texttt{SESSIONS} for borrower). \\
\texttt{idx\_agent\_status} &
  \texttt{AGENT\_ACTION\_LOG}(\texttt{status})\newline
  \texttt{WHERE status = 'PENDING'} &
  Partial index for the agent status-monitoring loop. \\
\texttt{idx\_risk\_score} &
  \texttt{LOAN\_RISK\_ASSESSMENT}(\texttt{anomaly\_score DESC})\newline
  \texttt{WHERE anomaly\_score > 0.5} &
  Partial index for AML alert queue (SAR workflow). \\
\texttt{idx\_ib\_loan\_status} &
  \texttt{INTERBANK\_LOAN}(\texttt{status}, \texttt{lender\_institution\_id}) &
  IBLP default cascade and maturity monitoring. \\
\texttt{idx\_institution\_type} &
  \texttt{INSTITUTION}(\texttt{institution\_type}, \texttt{status}) &
  Tier-scoped institution lookups for multi-entity operations. \\
\texttt{idx\_txn\_institution} &
  \texttt{TRANSACTION}(\texttt{origin\_institution\_id}, \texttt{created\_at DESC}) &
  Institution-to-institution ledger history retrieval. \\
\texttt{idx\_chainlink\_feed} &
  \texttt{CHAINLINK\_PRICE\_ROUND}(\texttt{feed\_address},\newline
  \texttt{updated\_at DESC}) &
  Latest round lookup per Chainlink feed for LTV staleness checks. \\
\texttt{idx\_session\_parent} &
  \texttt{SESSIONS}(\texttt{parent\_session\_id})\newline
  \texttt{WHERE parent\_session\_id IS NOT NULL} &
  Partial index for session lineage chain traversal. \\
\texttt{idx\_collateral\_borrower} &
  \texttt{COLLATERAL\_POSITION}(\texttt{borrower\_id}, \texttt{loan\_id}) &
  Health-factor lookup per borrower and per loan. \\
\texttt{idx\_event\_log\_txhash} &
  \texttt{BLOCKCHAIN\_EVENT\_LOG}(\texttt{tx\_hash}) &
  Deduplication guard: event listener checks this before INSERT. \\
\texttt{idx\_event\_log\_block} &
  \texttt{BLOCKCHAIN\_EVENT\_LOG}(\texttt{block\_number}, \texttt{chain\_id}) &
  Catchup-sync replay from a given block. \\
\bottomrule
\end{tabular}
\endgroup
\TableScreenshotEnd{Indexing Strategy: B-tree Indexes for Time-Window and High-Frequency Queries}
\end{table}

\section{Functional Dependencies}

\begin{table}[H]
\centering
\caption[Representative functional dependencies]{Representative functional dependencies.}
\label{tab:functional-deps}
\begingroup
\footnotesize
\setlength{\tabcolsep}{3pt}
\renewcommand{\arraystretch}{1.15}
\begin{tabular}{@{}%
  >{\raggedright\arraybackslash}p{0.17\linewidth}%
  >{\raggedright\arraybackslash}p{0.43\linewidth}%
  >{\raggedright\arraybackslash}p{0.36\linewidth}@{}}
\toprule
\tableheadcolor
\textbf{Relation} & \textbf{Functional Dependency} & \textbf{Notes} \\
\midrule
LOAN\_REQUEST & \texttt{request\_id} $\to$ all attributes & Application row keyed by request identifier. \\
LOAN\_REQUEST & \texttt{borrower\_id, local\_bank\_id} $\to$ \texttt{status}, \texttt{oracle\_state}, \texttt{amount}, \ldots & Lifecycle \texttt{status} and oracle pipeline \texttt{oracle\_state} are separate columns. \\
LOAN & \texttt{loan\_id} $\to$ all attributes & Disbursement record keyed by loan identifier. \\
LOAN & \texttt{loan\_request\_id} $\to$ \texttt{loan\_id} & 1:1 link from approved request to disbursed loan. \\
BORROWING\_LIMIT & \texttt{borrower\_id} $\to$ \texttt{six\_month\_limit}, \texttt{one\_year\_limit}, \newline \texttt{on\_chain\_enforced\_limit}, \texttt{last\_synced\_block}, \ldots & 1:1 with BORROWER; event-listener owned projection. \\
CREDIT\_PASSPORT & \texttt{passport\_id} $\to$ all attributes;\newline \texttt{borrower\_id} $\to$ \texttt{passport\_id} & Read-model projection of on-chain SBT; event-listener owned. \\
BORROWER & \texttt{wallet\_address} $\to$ \texttt{borrower\_id}, all attributes;\newline \texttt{registered\_local\_bank\_id} $\to$ home Local Bank;\newline \texttt{kyc\_level} $\to$ borrowing-cap tier (Section~\ref{sec:tiered-kyc}) & Unique candidate key in the prototype; Phase~IV may use \texttt{(wallet\_address, chain\_id)}. \\
INSTITUTION & \texttt{institution\_id} $\to$ \texttt{institution\_type}, \texttt{name}, \texttt{wallet\_address}, \texttt{contract\_address}, \texttt{status} & Supertype for cross-tier FK anchors. \\
NATIONAL\_BANK & \texttt{country\_code} $\to$ \texttt{institution\_id} & One national reserve per nation (\texttt{UNIQUE}). \\
MODEL\_REGISTRY & \texttt{model\_id} $\to$ \texttt{model\_name}, \texttt{version}, \texttt{training\_date}, \texttt{metrics}, \texttt{is\_active} & Model lineage for explainability audit. \\
LOAN\_RISK\_ASSESSMENT & \texttt{assessment\_id} $\to$ all attributes;\newline \texttt{model\_id} $\to$ \texttt{fraud\_probability}, \texttt{anomaly\_score}, \texttt{shap\_vector} & Per-loan ML inference; FK to \texttt{MODEL\_REGISTRY}. \\
INSTALLMENT & \texttt{loan\_id, installment\_number} $\to$ \texttt{amount}, \texttt{due\_date}, \texttt{status}, \texttt{payment\_event\_log\_id} & Composite primary key. \\
SESSIONS & \texttt{session\_id} $\to$ \texttt{borrower\_id}, \texttt{wallet\_address}, \texttt{session\_key\_hash}, \newline \texttt{session\_key\_scope}, \texttt{session\_key\_expires\_at}, \newline \texttt{expires\_at}, \texttt{revoked} & Retail-client agent sessions only at MVT; bank-authority sessions are Future Work. \\
AGENT\_ACTION\_LOG & \texttt{action\_id} $\to$ \texttt{session\_id}, \texttt{tool\_name}, \newline \texttt{parameters}, \texttt{confirmation\_turn\_id}, \newline \texttt{onchain\_tx\_hash}, \texttt{status}, \texttt{created\_at} & Links every write-tool execution to the user's consent turn in \texttt{AGENT\_CONVERSATION\_TURN}. \\
ASSETS & \texttt{asset\_id} $\to$ \texttt{symbol}, \texttt{name}, \texttt{asset\_type}, \newline \texttt{network}, \texttt{decimals}, \texttt{is\_active};\newline \texttt{symbol} $\to$ \texttt{asset\_id} & \texttt{symbol} is a unique candidate key. \\
INTEREST\_RATE\_TIER & \texttt{tier\_id} $\to$ \texttt{base\_rate}, \texttt{kink\_utilization}, \newline \texttt{rate\_above\_kink}, \texttt{max\_rate}, \texttt{bank\_tier\_type}, \texttt{effective\_from}, \texttt{effective\_to}, \texttt{is\_active} & Eliminates transitive dependency; temporal versioning for governance changes. \\
\bottomrule
\end{tabular}
\endgroup
\TableScreenshotEnd{Functional Dependencies in the Core Lending and Identity Schema}
\end{table}

\section{Materialized Views and Agent Query Optimisation}

The MCP tool \texttt{get\_account\_state} joins \texttt{BORROWER}, \texttt{CREDIT\_PASSPORT}, \texttt{COLLATERAL\_POSITION}, active \texttt{LOAN} rows, pending \texttt{INSTALLMENT} rows, \texttt{SAVINGS\_ACCOUNT}, and \texttt{BORROWING\_LIMIT}. To avoid six-to-seven table joins on every agent call, the schema defines a materialized view refreshed by the event listener after any event that modifies component tables:

\begin{verbatim}
CREATE MATERIALIZED VIEW mv_client_dashboard AS
SELECT
    b.borrower_id,
    b.wallet_address,
    b.kyc_level,
    cp.credit_score,
    cp.risk_tier,
    bl.six_month_remaining,
    bl.one_year_remaining,
    count(l.loan_id) FILTER (WHERE l.status = 'ACTIVE') AS active_loan_count,
    max(i.due_date) FILTER (WHERE i.status = 'PENDING') AS next_installment_due,
    sum(sa.balance) AS total_savings_balance,
    sum(col.locked_amount) AS total_collateral_locked
FROM borrower b
LEFT JOIN credit_passport cp USING (borrower_id)
LEFT JOIN borrowing_limit bl USING (borrower_id)
LEFT JOIN loan l USING (borrower_id)
LEFT JOIN installment i USING (loan_id)
LEFT JOIN savings_account sa USING (borrower_id)
LEFT JOIN collateral_position col USING (borrower_id)
GROUP BY b.borrower_id, b.wallet_address, b.kyc_level,
         cp.credit_score, cp.risk_tier,
         bl.six_month_remaining, bl.one_year_remaining;

CREATE UNIQUE INDEX ON mv_client_dashboard (borrower_id);
\end{verbatim}

\noindent Refresh strategy: \texttt{REFRESH MATERIALIZED VIEW CONCURRENTLY mv\_client\_dashboard} triggered by the event listener after processing events on projection or lending tables (Phase~III).

\section{Representative SQL Queries}
\label{sec:sql-queries}

The following queries illustrate Database Design coursework competencies (joins, aggregates, subqueries, and views) on the prototype schema. They map to FR-02, FR-04, and FR-07 in Table~\ref{tab:functional-requirements}.

\paragraph{Q1 --- Active loans per Local Bank with borrower KYC level (INNER JOIN + aggregate).}
\begin{verbatim}
SELECT lb.name AS local_bank,
       COUNT(l.loan_id) AS active_loans,
       AVG(l.principal_amount) AS avg_principal
FROM loan l
JOIN local_bank lb ON l.local_bank_id = lb.institution_id
JOIN borrower b ON l.borrower_id = b.borrower_id
WHERE l.status = 'ACTIVE'
GROUP BY lb.institution_id, lb.name
ORDER BY active_loans DESC;
\end{verbatim}

\paragraph{Q2 --- Overdue installments with borrower contact hash (JOIN + filter).}
\begin{verbatim}
SELECT i.loan_id, i.installment_number, i.due_date, i.amount,
       b.wallet_address, b.kyc_level
FROM installment i
JOIN loan l ON i.loan_id = l.loan_id
JOIN borrower b ON l.borrower_id = b.borrower_id
WHERE i.status = 'PENDING' AND i.due_date < CURRENT_DATE
ORDER BY i.due_date ASC;
\end{verbatim}

\paragraph{Q3 --- Borrowers exceeding six-month rolling limit (subquery).}
\begin{verbatim}
SELECT b.borrower_id, b.wallet_address,
       bl.six_month_limit,
       (SELECT COALESCE(SUM(t.amount), 0)
        FROM transaction t
        WHERE t.borrower_id = b.borrower_id
          AND t.transaction_type = 'LOAN_DISBURSEMENT'
          AND t.created_at >= NOW() - INTERVAL '6 months') AS disbursed_6mo
FROM borrower b
JOIN borrowing_limit bl ON b.borrower_id = bl.borrower_id
WHERE (SELECT COALESCE(SUM(t.amount), 0) FROM transaction t
       WHERE t.borrower_id = b.borrower_id
         AND t.transaction_type = 'LOAN_DISBURSEMENT'
         AND t.created_at >= NOW() - INTERVAL '6 months')
      > bl.six_month_limit;
\end{verbatim}

\paragraph{Q4 --- High-risk assessments awaiting approver review (JOIN + ORDER BY).}
\begin{verbatim}
SELECT lr.request_id, lra.composite_score, lra.anomaly_score,
       lra.created_at, mr.model_name, mr.version
FROM loan_request lr
JOIN loan_risk_assessment lra ON lr.request_id = lra.loan_request_id
JOIN model_registry mr ON lra.model_id = mr.model_id
WHERE lr.status = 'UNDER_REVIEW'
  AND lra.composite_score BETWEEN 0.4 AND 0.7
ORDER BY lra.composite_score DESC;
\end{verbatim}

\paragraph{Q5 --- Client dashboard via materialized view (simplified SELECT).}
\begin{verbatim}
SELECT borrower_id, wallet_address, credit_score, risk_tier,
       active_loan_count, next_installment_due, total_savings_balance
FROM mv_client_dashboard
WHERE borrower_id = $1;
\end{verbatim}

\chapter{Agent Harness Reference}
\label{app:agent-harness}

\section{Per-User Context and Prompt Assembly}

The Express.js context injection layer:

\begin{verbatim}
{
  "borrower_id": "550e8400-e29b-41d4-a716-446655440000",
  "wallet_address": "0xABC...",
  "language": "en",
  "credit_tier": "Silver",
  "credit_score": 512,
  "active_loans": [
    {
      "loan_id": "L-4821",
      "amount": 100,
      "next_due": "2026-06-15",
      "installment": 17.5
    }
  ],
  "savings_balance": 250.0,
  "pool_utilisation": 0.67,
  "kyc_tier": 1,
  "conversation_history": [ "...last 10 turns..." ]
}
\end{verbatim}

\section{Session Lineage, Compression, and Toolset Scoping}

\begin{table}[H]
\centering
\caption[MCP toolset scoping]{MCP toolset scoping.}
\label{tab:agent-toolsets}
\begin{tabular}{p{3.5cm}p{9cm}}
\toprule
\tableheadcolor
\textbf{Toolset} & \textbf{Tools visible} \\
\midrule
\texttt{read\_only} & Full-suite taxonomy (Future Work): all 9 read tools; the MVT deploys only the subset needed for the selected 3--5 tools. \\
\texttt{loan\_actions} & \texttt{read\_only} $+$ loan, installment, and group write tools. \\
\texttt{account\_management} & \texttt{read\_only} $+$ deposit, KYC, and reminder write tools. \\
\bottomrule
\end{tabular}
\end{table}

\section{Lifecycle Hooks and Optional Capabilities}

\paragraph{Lifecycle hook middleware.}
Write-tool POST requests pass through a hook stack: (1)~confirmation audit hook (Phase~II) verifies a logged confirmation turn or returns HTTP~403; (2)~session-key scope pre-check (Phase~III); (3)~AML pre-check against \texttt{LOAN\_RISK\_ASSESSMENT} and \texttt{SECURITY\_EVENT\_LOG} (Phase~IV).

\paragraph{Agent capabilities.}
\begin{itemize}[noitemsep]
\item \textbf{Proactive installment reminders} via \texttt{get\_installment\_schedule} and \texttt{pay\_installment}.
\item \textbf{Credit Passport coaching} via \texttt{get\_credit\_score} and Table~\ref{tab:credit-tiers}.
\item \textbf{Loan calculator} with live FX from \texttt{get\_market\_data}.
\item \textbf{Group lending coordination} and \textbf{KYC tier upgrade guidance} via write/read tools.
\item \textbf{Cross-session context retrieval} via \texttt{get\_session\_history}.
\end{itemize}

\chapter{Smart Contract Capabilities}
\label{app:smart-contract-capabilities}

\begin{table}[H]
\centering
\caption[Smart contract capabilities summary]{Smart contract capabilities summary (17 modular contracts).}
\label{tab:appendix-contract-capabilities}
\footnotesize
\setlength{\tabcolsep}{3.5pt}
\renewcommand{\arraystretch}{1.08}
\begin{tabular}{@{}>{\raggedright\arraybackslash}p{2.4cm}>{\raggedright\arraybackslash}p{1.6cm}>{\centering\arraybackslash}p{0.9cm}>{\raggedright\arraybackslash}p{1.5cm}>{\raggedright\arraybackslash}p{6.2cm}@{}}
\toprule
\tableheadcolor
\textbf{Contract} & \textbf{Module} & \textbf{Ph.} & \textbf{Status} & \textbf{Primary capability} \\
\midrule
WorldBankReserve & Tier~1 & I & Specified & Reserve custody, national-bank registration, downward allocation \\
\rowcolor{RowShade}
NationalBank & Tier~2 & I & Specified & Local-bank registration, borrow upward, allocate downward \\
LocalBank & Tier~3 & I & Specified & Client onboarding, approvers, \texttt{LoanController}, \texttt{freezeAccount} \\
\rowcolor{RowShade}
LoanController & Lending & I--II & Shell / Ph.~II & Request, approve, disburse, installment state machine \\
SavingsVault & Deposits & II & Specified & Variable-yield savings; reserve-gated withdrawals \\
\rowcolor{RowShade}
FixedDeposit & Deposits & II+ & Stub & Term deposits; early-withdrawal penalty \\
GroupLendingPool & Group & II & Specified & Group consent, shared liability, per-member disbursement \\
\rowcolor{RowShade}
InsuranceFund & Risk & I--II & Should & 5\% interest capture; default coverage; PoR attestation \\
CurrentAccount & Payments & II+ & Stub & Atomic transfers; recurring payments \\
\rowcolor{RowShade}
InterBankLendingPool & Multi-entity & II & Specified & Same-tier peer pools (1 / 7 / 30 day) \\
UpwardDepositFacility & Multi-entity & II & Specified & Surplus repatriation to parent tier \\
\rowcolor{RowShade}
SyndicatedLoan & Multi-entity & III & Planned & Multi-bank co-funding and consent vote \\
TranchedPool & Multi-entity & III & Planned & Senior / junior tranches; waterfall recovery \\
\rowcolor{RowShade}
TreasurySwap & Multi-entity & III & Planned & Institutional oracle-priced asset swap \\
FXModule & Treasury / FX & III & Planned & Retail FX conversion and audit trail \\
\rowcolor{RowShade}
LiquidationEngine & Risk & III & Planned & Health-factor monitor; liquidator bonus \\
NettingEngine & Multi-entity & --- & Future Work & Multilateral settlement netting (stub only) \\
\bottomrule
\end{tabular}
\end{table}


%% APPENDIX D: SMART CONTRACT INTERFACE
\chapter*{Appendix D: WorldBankReserve Contract Interface}
\phantomsection
\addcontentsline{toc}{chapter}{Appendix D: WorldBankReserve Contract Interface}

The following Solidity interface documents the public API of the \texttt{WorldBankReserve} contract---the formally specified Tier~1 contract (Phase~I deployment target). Three design annotations follow the listing:

\begin{lstlisting}[language=Solidity,basicstyle=\ttfamily\scriptsize,breaklines=true,
    commentstyle=\color{GrayText},keywordstyle=\color{PrimaryBlue}\bfseries,
    frame=single,rulecolor=\color{AccentBlue!40},captionpos=b,
    caption={WorldBankReserve.sol -- Tier 1 public interface. Full implementation in project repository.}]
pragma solidity ^0.8.20;

import "@openzeppelin/contracts/security/ReentrancyGuard.sol";
import "@openzeppelin/contracts/access/AccessControl.sol";

contract WorldBankReserve is ReentrancyGuard, AccessControl {

    bytes32 public constant NATIONAL_BANK_ROLE = keccak256("NATIONAL_BANK_ROLE");
    bytes32 public constant AUDITOR_ROLE       = keccak256("AUDITOR_ROLE");

    uint256 public totalReserve;
    uint256 public minimumReserveRatio;
    mapping(address => uint256) public allocatedTo;

    event CapitalAllocated(address indexed nationalBank, uint256 amount);
    event RepaymentReceived(address indexed nationalBank, uint256 amount);
    event ReserveRatioUpdated(uint256 newRatio);

    function allocateCapital(address nationalBank, uint256 amount)
        external onlyRole(DEFAULT_ADMIN_ROLE) nonReentrant;

    function recordRepayment(uint256 amount)
        external onlyRole(NATIONAL_BANK_ROLE) nonReentrant;

    function utilizationRate() external view returns (uint256);

    function isReserveAdequate() external view returns (bool);

    function getReserveSummary() external view returns (
        uint256 totalDeposited,
        uint256 totalLoaned,
        uint256 reserveRatio,
        uint256 insuranceFundBalance
    );
}
\end{lstlisting}

\chapter*{Appendix E: Internal Industry Research Notes}
\phantomsection
\addcontentsline{toc}{chapter}{Appendix E: Internal Industry Research Notes}
\label{app:internal-research}

\noindent Repository: \href{https://github.com/Yuvraajrahman/Cryto-World-Bank}{github.com/Yuvraajrahman/Cryto-World-Bank}; path \texttt{Documentation/2nd Phase (Bokhtiar)/Improvements/}.

\begin{table}[H]
\centering
\caption[Internal industry research corpus]{Internal industry research corpus.}
\label{tab:internal-research}
\begingroup
\footnotesize
\setlength{\tabcolsep}{4pt}
\renewcommand{\arraystretch}{1.15}
\begin{tabular}{@{}%
  >{\raggedright\arraybackslash}p{0.07\linewidth}%
  >{\raggedright\arraybackslash}p{0.36\linewidth}%
  >{\raggedright\arraybackslash}p{0.53\linewidth}@{}}
\toprule
\tableheadcolor
\textbf{ID} & \textbf{Document} & \textbf{Thesis use} \\
\midrule
IR-1 & \texttt{Binance\_FTX\_Full\_Report.md} & CEX vs.\ CWB comparison; FTT/BNB token lesson; SAFU and Binance Earn analogues; PoR and commingling safeguards \\
IR-2 & \texttt{blockchain platforms research.csv} & Five-exchange matrix (Binance, Bybit, Coinbase, OKX, Kraken) \\
IR-3 & \texttt{Binance Business Analysis.md} & Exchange revenue taxonomy (fees, listing, OTC) \\
IR-4 & \texttt{binance\_software\_engineering\_architecture.md} & Nine-actor taxonomy (Section~\ref{sec:onboarding}); use-case and sequence-diagram conventions \\
IR-5 & \texttt{binance\_architecture\_research.md} & Off-chain PostgreSQL/Redis patterns (background; not cited in body text) \\
\bottomrule
\end{tabular}
\endgroup
\end{table}

\chapter*{Appendix F: UI Prototype Wireframes (Figma)}
\phantomsection
\addcontentsline{toc}{chapter}{Appendix F: UI Prototype Wireframes (Figma)}
\label{app:ui-prototypes}

\paragraph{Retail loan application flow.}
\begin{figure}[H]
\centering
\fbox{\parbox{0.88\linewidth}{\centering\vspace{3.2cm}\textit{[Insert Figma export: retail loan application wireframe / prototype]}\vspace{3.2cm}}}
\caption[Retail loan application prototype (Figma)]{Retail loan application prototype (Figma).}
\label{fig:ui-wireframe-retail}
\end{figure}

\paragraph{Authority Brief approver dashboard.}
\begin{figure}[H]
\centering
\fbox{\parbox{0.88\linewidth}{\centering\vspace{3.2cm}\textit{[Insert Figma export: Authority Brief approver dashboard]}\vspace{3.2cm}}}
\caption[Authority Brief approver dashboard prototype (Figma)]{Authority Brief approver dashboard prototype (Figma).}
\label{fig:ui-wireframe-approver}
\end{figure}

\paragraph{Client onboarding funnel (KYC levels).}
\begin{figure}[H]
\centering
\fbox{\parbox{0.88\linewidth}{\centering\vspace{3.2cm}\textit{[Insert Figma export: five-stage onboarding funnel / KYC status screens]}\vspace{3.2cm}}}
\caption[Client onboarding funnel prototype (Figma)]{Client onboarding funnel prototype (Figma).}
\label{fig:ui-wireframe-onboarding}
\end{figure}

\paragraph{Reserve transparency dashboard (Tier~1).}
\begin{figure}[H]
\centering
\fbox{\parbox{0.88\linewidth}{\centering\vspace{3.2cm}\textit{[Insert Figma export: World Bank reserve balances and allocation overview]}\vspace{3.2cm}}}
\caption[Reserve transparency dashboard prototype (Figma)]{Reserve transparency dashboard prototype (Figma).}
\label{fig:ui-wireframe-reserve}
\end{figure}

\paragraph{Conversational banking agent (chat).}
\begin{figure}[H]
\centering
\fbox{\parbox{0.88\linewidth}{\centering\vspace{3.2cm}\textit{[Insert Figma export: MCP chat assistant with confirmation-gated actions]}\vspace{3.2cm}}}
\caption[Conversational banking agent prototype (Figma)]{Conversational banking agent prototype (Figma).}
\label{fig:ui-wireframe-agent}
\end{figure}

\paragraph{Group solidarity lending application.}
\begin{figure}[H]
\centering
\fbox{\parbox{0.88\linewidth}{\centering\vspace{3.2cm}\textit{[Insert Figma export: group loan request and member share breakdown]}\vspace{3.2cm}}}
\caption[Group solidarity lending prototype (Figma)]{Group solidarity lending prototype (Figma).}
\label{fig:ui-wireframe-group-lending}
\end{figure}

\paragraph{Compliance and AML alert review.}
\begin{figure}[H]
\centering
\fbox{\parbox{0.88\linewidth}{\centering\vspace{3.2cm}\textit{[Insert Figma export: SAR/AML alert queue and officer review screen]}\vspace{3.2cm}}}
\caption[Compliance and AML alert review prototype (Figma)]{Compliance and AML alert review prototype (Figma).}
\label{fig:ui-wireframe-compliance}
\end{figure}

\paragraph{Mobile navigation and responsive layout.}
\begin{figure}[H]
\centering
\fbox{\parbox{0.88\linewidth}{\centering\vspace{3.2cm}\textit{[Insert Figma export: mobile nav, breakpoints, and component library sheet]}\vspace{3.2cm}}}
\caption[Mobile navigation and responsive layout prototype (Figma)]{Mobile navigation and responsive layout prototype (Figma).}
\label{fig:ui-wireframe-mobile}
\end{figure}

%% REFERENCES
\chapter*{References}
\phantomsection
\addcontentsline{toc}{chapter}{References}

{\small
\begin{enumerate}[label={[\arabic*]}]
\item S.~M. Werner, D.~Perez, L.~Gudgeon, A.~Klages-Mundt, D.~Harz, and W.~J. Knottenbelt. 2022. SoK: Decentralized Finance (DeFi). \emph{arXiv preprint arXiv:2101.08778}. \url{https://arxiv.org/abs/2101.08778}


\item M.~Bastankhah, V.~Nadkarni, C.~Jin, S.~Kulkarni, and P.~Viswanath. 2024. Thinking Fast and Slow: Data-Driven Adaptive DeFi Borrow-Lending Protocol. In \emph{Proc. 6th Conf. Advances in Financial Technologies (AFT)} (LIPIcs, Vol.~316). Article 27, 23 pages. \url{https://doi.org/10.4230/LIPIcs.AFT.2024.27}


\item G.~Palaiokrassas, S.~Scherrers, I.~Ofeidis, and L.~Tassiulas. 2024. Leveraging Machine Learning for Multichain DeFi Fraud Detection. In \emph{Proc. IEEE Int. Conf. Blockchain and Cryptocurrency (ICBC)}. \url{https://doi.org/10.1109/ICBC59979.2024.10634350}


\item D.~Adom, E.~O. Acheampong, and M.~Boateng. 2022. LIME and SHAP: A Comparison of Model-Agnostic Approaches to Explainability in Loan Approval Systems. \emph{International Journal of Advanced Computer Science and Applications} 13, 11. \url{https://thesai.org/Downloads/Volume13No11/Paper_90-LIME_and_SHAP_A_Comparison.pdf}


\item B.~W. Tan. 2023. Central Bank Digital Currency and Financial Inclusion. \emph{IMF Working Paper WP/23/69}. \url{https://www.imf.org/en/Publications/WP/Issues/2023/03/27/Central-Bank-Digital-Currency-and-Financial-Inclusion-531273}


\item F.~T. Liu, K.~M. Ting, and Z.-H. Zhou. 2008. Isolation Forest. In \emph{Proc. IEEE Int. Conf. Data Mining (ICDM)}, pages 413--422. \url{https://doi.org/10.1109/ICDM.2008.17}


\item N.~Atzei, M.~Bartoletti, and T.~Cimoli. 2017. A Survey of Attacks on Ethereum Smart Contracts (SoK). In \emph{Proc. Int. Conf. Principles of Security and Trust (POST)}, pages 164--186. \url{https://doi.org/10.1007/978-3-662-54455-6_8}


\item DefiLlama. 2024. Lending Protocols --- DeFi TVL and Protocol Rankings. \url{https://defillama.com/protocols/Lending}


\item World Bank. 2021. The Global Findex Database 2021: Financial Inclusion, Digital Payments, and Resilience. \url{https://www.worldbank.org/en/publication/globalfindex}


\item Aave. 2024. Aave Protocol Documentation. \url{https://docs.aave.com/}


\item Compound. 2024. Compound Protocol Documentation. \url{https://docs.compound.finance/}


\item World Bank. 2022. World Development Report 2022: Finance for an Equitable Recovery. \url{https://www.worldbank.org/en/publication/wdr2022}


\item International Finance Corporation (IFC). 2017. MSME Finance Gap: Assessment of the Shortfalls and Opportunities in Financing Micro, Small, and Medium Enterprises in Emerging Markets. World Bank. \url{https://openknowledge.worldbank.org/entities/publication/ff4c9839-21ac-5676-a23a-7cf6f745df0c}


\item Bank for International Settlements. 2024. DeFi Lending: Intermediation Without Information?. BIS Working Paper No.~1183. \url{https://www.bis.org/publ/work1183.pdf}


\item Committee on Payments and Market Infrastructures and Bank for International Settlements. 2016. Correspondent Banking --- Technical Report. CPMI Papers No.~147. \url{https://www.bis.org/cpmi/publ/d147.pdf}


\item Ripple Labs. 2025. RLUSD Stablecoin Product Overview. Ripple Documentation. \url{https://ripple.com/solutions/stablecoin/}


\item World Bank. 2024. Remittance Prices Worldwide: Making Markets More Transparent. Migration and Development Brief~40. \url{https://remittanceprices.worldbank.org/}


\item DefiLlama. 2026. Aave V3 TVL, Fees, Revenue \& Income Statement. March. \url{https://defillama.com/protocol/aave-v3}


\item DefiLlama. 2026. Compound V3 TVL, Fees, Revenue \& Income Statement. March. \url{https://defillama.com/protocol/compound-v3}


\item Fensory. 2026. Sky Protocol Projects \$611M Revenue in 2026 as USDS Supply Targets \$20.6 Billion. February. \url{https://www.fensory.com/intelligence/rwa/sky-protocol-tokenization-regatta-solana-february-2026}


\item Fensory. 2026. Maple Finance --- Institutional DeFi Lending Protocol. \url{https://www.fensory.com/insights/protocols/maple-finance}


\item Fensory. 2026. DeFi Credit Platforms Hit \$2.4B Milestone. March. \url{https://www.fensory.com/intelligence/rwa/defi-private-credit-platforms-2-4-billion-loans-ethereum}


\item Financial Stability Board. 2020. Enhancing Cross-border Payments: Stage 3 Roadmap. \url{https://www.fsb.org/2020/10/enhancing-cross-border-payments-stage-3-roadmap/}


\item GlobeNewswire. 2025. R3's Corda Leads Tokenized RWA Market with Over \$10 Billion in On-chain Assets. February. \url{https://www.globenewswire.com/news-release/2025/02/13/3025637/0/en/R3-s-Corda-leads-tokenized-RWA-market-with-over-10-billion-in-on-chain-assets-and-unrivalled-industry-adoption.html}


\item World Bank. 2025. World Bank Group Tracks Project Funds with New Blockchain Tool. Press Release, September. \url{https://www.worldbank.org/en/news/press-release/2025/09/26/world-bank-group-tracks-project-funds-with-new-blockchain-tool}


\item JPMorgan. 2025. Kinexys Digital Payments: Real-Time Multicurrency Payments. \url{https://www.jpmorgan.com/kinexys/index}


\item World Economic Forum. 2022. Why Decentralized Finance Is a Leapfrog Technology for the 1.1 Billion People Who Are Unbanked. September. \url{https://www.weforum.org/stories/2022/09/decentralized-finance-a-leapfrog-technology-for-the-unbanked/}


\item Opera Newsroom. 2026. 160M CELO Allocation Proposal to Grow Opera from Distribution Partner into Key Network Stakeholder. March. \url{https://press.opera.com/2026/03/19/opera-celo-partnership-2026/}


\item Morpho. 2026. Network Data. March. \url{https://data.morpho.org/}


\item Stellar. 2025. End of Year 2025 Report --- Execution at Scale. \url{https://www.stellar.org/blog/foundation-news/2025-year-in-review}


\item Y.~Chen, S.~Bin, and H.~He. 2024. Design and Implementation of a Multi-Chain Lending Model in Blockchain. In \emph{Proc. 3rd Int. Conf. Artificial Intelligence and Computer Information Technology (AICIT)}, pages 97--102.


\item S.~Yang and W.~Cui. 2023. An Evaluation System for DeFi Lending Protocols. \emph{arXiv preprint arXiv:2303.01022}. \url{https://arxiv.org/abs/2303.01022}


\item J.~Hartmann and O.~Hasan. 2021. A Social-Capital Based Approach to Blockchain-Enabled Peer-to-Peer Lending. In \emph{Proc. IEEE Int. Conf. Blockchain Computing and Applications (BCCA)}, pages 105--110. \url{https://hal.science/hal-03371872}


\item P.~T.~Hasan, H.~K.~T.~Akhter, M.~R.~Tahsinur, A.~B.~Haque, A.~K.~M.~N.~Islam, and R.~M.~Rahman. 2022. Blockchain and Machine Learning for Fraud Detection: A Privacy-Preserving and Adaptive Incentive Based Approach. \emph{IEEE Access}, vol.~10, pages 87115--87131. \url{https://ieeexplore.ieee.org/document/9857827}


\item Bank for International Settlements et al. 2020. Central Bank Digital Currencies: Foundational Principles and Core Features. BIS and Central Banks Joint Report. \url{https://www.bis.org/publ/othp33.htm}


\item M.~Asaduzzaman, F.~Hasib, and Z.~B.~Hafiz. 2020. Towards Using Blockchain Technology for Microcredit Industry in Bangladesh. In \emph{Proc. 23rd Int. Conf. Computer and Information Technology (ICCIT)}, pages 1--6. \url{https://doi.org/10.1109/ICCIT51783.2020.9392730}


\item H.~Kim and D.~Kim. 2024. Optimal Gas Fee Minimization in DeFi: Enhancing Efficiency and Security on the Ethereum Blockchain. \emph{IEEE Access}, vol.~12, pages 173810--173823. \url{https://ieeexplore.ieee.org/document/10750187}


\item X.~Yuan. 2024. SHAP-based Interpretable Models for Credit Default Assessment Using Machine Learning. In \emph{Proc. 14th Int. Conf. Software Technology and Engineering (ICSTE)}, pages 213--217. \url{https://doi.org/10.1109/ICSTE63875.2024.00044}


\item T.~Dao, T.~Trinh, and V.~Pham. 2025. Optimizing Credit Scoring Models for Decentralized Financial Applications. In \emph{Information and Communication Technology}, Springer. \url{https://doi.org/10.1007/978-981-96-4282-3_36}


\item A.~Beniiche. 2020. A Study of Blockchain Oracles. \emph{arXiv preprint arXiv:2004.07140}. \url{https://arxiv.org/pdf/2004.07140}


\item A.~Pasdar, Y.~C.~Lee, and Z.~Dong. 2023. Connect API with Blockchain: A Survey on Blockchain Oracle Implementation. \emph{ACM Computing Surveys}, 55, 10. \url{https://doi.org/10.1145/3567582}


\item L.~Gudgeon, S.~M.~Werner, D.~Perez, and W.~J.~Knottenbelt. 2020. DeFi Protocols for Loanable Funds: Interest Rates, Liquidity and Market Efficiency. In \emph{Proc. 2nd ACM Conf. Advances in Financial Technologies (AFT)}, pages 92--112. \url{https://doi.org/10.1145/3419614.3423254}


\item K.~W.~Wu. 2024. Strengthening DeFi Security: A Static Analysis Approach to Flash Loan Vulnerabilities. \emph{arXiv preprint arXiv:2411.01230}.


\item Y.~Li et al. 2025. A Comprehensive Study of Exploitable Patterns in Smart Contracts: From Vulnerability to Defense. \emph{arXiv preprint arXiv:2504.21480}. \url{https://arxiv.org/abs/2504.21480}


\item F.~Piper, K.~Wolf, and J.~Heiss. 2025. Privacy-Preserving On-chain Permissioning for KYC-Compliant Decentralized Applications. TU Berlin, \emph{arXiv preprint arXiv:2510.05807}.


\item N.~Decker. 2025. Zero-Knowledge Proofs: Cryptographic Model for Financial Compliance and Global Banking Security. \emph{SSRN Working Paper No.~5170068}.


\item Grameen Bank. 2024. Grameen Bank at a Glance. Grameen Communications. \url{https://www.grameen-info.org/}


\item A.~Carstens et al. 2021. Stablecoins: Risks, Potential and Regulation. \emph{BIS Working Papers No.~905}, Bank for International Settlements. \url{https://www.bis.org/publ/work905.pdf}


\item J.~A.~Ocampo and K.~Gallagher. 2024. The Role of Multilateral Development Banks and Development Assistance in the Provision of Global Public Goods. UNDP Background Paper. \url{https://hdr.undp.org/system/files/documents/background-paper-document/2024jaocampokdgonzaleztheroleofmultilateraldevlmntbanks.pdf}


\item F.~A.~Aponte-Novoa, A.~L.~S.~Orozco, R.~Villanueva-Polanco, and P.~Wightman. 2021. The 51\% Attack on Blockchains: A Mining Behavior Study. \emph{IEEE Access}, vol.~9, pages 140549--140564. \url{https://doi.org/10.1109/ACCESS.2021.3119291}


\item Sygnum Bank. 2026. Institutional DeFi in 2025: The Disconnect Between Infrastructure and Allocation. Sygnum Blog, May 2025. \url{https://www.sygnum.com/blog/2025/05/30/institutional-defi-in-2025-the-disconnect-between-infrastructure-and-allocation/}


\item M.~J.~Page, J.~E.~McKenzie, P.~M.~Bossuyt, I.~Boutron, T.~C.~Hoffmann, C.~D.~Mulrow, et~al. 2021. The PRISMA 2020 statement: an updated guideline for reporting systematic reviews. \emph{BMJ}, vol.~372, n71. \url{https://doi.org/10.1136/bmj.n71}


\item D.~Cheng, X.~Wang, Y.~Zhang, and L.~Zhang. 2022. Graph Neural Network for Fraud Detection via Spatial-Temporal Attention. \emph{IEEE Trans. Knowledge and Data Engineering}, 34, 8, pages 3800--3813. \url{https://ieeexplore.ieee.org/document/9356317}


\item D.~Wang, B.~Wu, X.~Yuan, L.~Wu, Y.~Zhou, and H.~Cui. 2024. DeFiGuard: A Price Manipulation Detection Service in DeFi Using Graph Neural Networks. \emph{arXiv preprint arXiv:2406.11157}. \url{https://arxiv.org/abs/2406.11157}


\item B.~McMahan, E.~Moore, D.~Ramage, S.~Hampson, and B.~A.~y Arcas. 2017. Communication-Efficient Learning of Deep Networks from Decentralized Data. In \emph{Proc. Int. Conf. Artificial Intelligence and Statistics (AISTATS)}.


\item Z.~Chen, J.~Chen, M.~Sra, et~al. 2025. Standard Benchmarks Fail --- Auditing LLM Agents in Finance Must Prioritize Risk. \emph{arXiv preprint arXiv:2502.15865}. \url{https://arxiv.org/abs/2502.15865}


\item MicroSave Consulting. 2025. Shaping a Sustainable Microfinance Sector in Bangladesh. MSC Signature Project, September. \url{https://www.microsave.net/signature_projects/shaping-a-sustainable-microfinance-sector-in-bangladesh-through-mscs-transformation-feasibility-study/}


\item European Parliament and Council of the European Union. 2024. Regulation (EU) 2023/1114 of 31 May 2023 on Markets in Crypto-Assets (MiCA). \emph{Official Journal of the European Union}, L~150/40, June 2023; fully applicable from 30 December. \url{https://eur-lex.europa.eu/eli/reg/2023/1114/oj}


\item European Banking Authority. 2024. Guidelines on the Minimum Content of the Governance Arrangements for Issuers of Asset-Referenced Tokens under MiCAR. EBA/GL//06, June. \url{https://www.eba.europa.eu/activities/single-rulebook/regulatory-activities/asset-referenced-and-e-money-tokens-micar/guidelines-internal-governance-arrangements-issuers-arts-under-micar}


\item United States Congress. 2025. Guiding and Establishing National Innovation for U.S. Stablecoins (GENIUS) Act. Pub.~L.~No.~119-27, July~18. \url{https://www.congress.gov/bill/119th-congress/senate-bill/919}


\item Bank for International Settlements Innovation Hub. 2024. Project mBridge: Connecting Economies through CBDC --- Reaches Minimum Viable Product. BIS Innovation Hub Report, June. \url{https://www.bis.org/about/bisih/topics/cbdc/mcbdc_bridge.htm}


\item Bank for International Settlements and Central Banks of France, Japan, Korea, Mexico, Switzerland, the UK, and the US. 2024. Project Agora: Tokenization of Cross-Border Payments using Unified Ledgers. BIS Innovation Hub Working Paper, April. \url{https://www.bis.org/about/bisih/topics/fmis/agora.htm}


\item Bangladesh Bank. 2025. FE Circular No.~06: Electronic Communication for Letters of Credit and Trade Finance Instruments. Foreign Exchange Policy Department, January. \url{https://www.bb.org.bd/mediaroom/circulars/fepd/jan172025fepd06.pdf}


\item Bangladesh Bank. 2025. Financial Inclusion in Bangladesh: Progress and Gender-Disaggregated Outlook 2024--25. Special Publication SP-02, Financial Inclusion Department, June. \url{https://www.bb.org.bd/pub/special/SP2025-02.pdf}


\item BRAC. 2025. BRAC Microfinance Annual Performance Report 2025: Women-Centric Lending Outcomes. BRAC Research and Evaluation Division. \url{https://www.brac.net/program/microfinance/}


\item Ethereum Improvement Proposals. 2025. EIP-4626: Tokenized Vault Standard. Ethereum Foundation, April 2022. ; ``EIP-7540: Asynchronous ERC-4626 Tokenized Vaults,'' Ethereum Foundation, 2023. ; ``EIP-3643: T-REX -- Token for Regulated EXchanges,'' Ethereum Foundation, 2021.  See also: Spectral Finance, ``On-Chain Credit Scoring for DeFi Lending,'' 2023. ; RociFi Labs, ``Undercollateralized DeFi Lending via On-Chain Credit Score,'' 2023. ; Qwen Team (Alibaba Cloud), ``Qwen3 Technical Report,'' \emph{arXiv preprint arXiv:2505.09388}. \url{https://eips.ethereum.org/EIPS/eip-4626}


\item Ethereum Improvement Proposals. 2022--2025. EIP-4626, EIP-7540, EIP-3643; Spectral Finance and RociFi documentation; Qwen Team. 2025. Qwen3 Technical Report (\emph{arXiv:2505.09388}). \url{https://eips.ethereum.org/EIPS/eip-4626}


\item Elliptic. 2019. The Elliptic Data Set. Kaggle. \url{https://www.kaggle.com/datasets/ellipticco/elliptic-data-set}


\item M.~Alizadeh, M.~Gholampour, A.~Imani, and J.~Schulz. 2026. Simple Prompt Injection Attacks Can Leak Personal Data Observed by LLM Agents During Task Execution. \emph{arXiv preprint arXiv:2506.01055}, University of Zurich / IPM, June.


\item OWASP Foundation. 2026. OWASP Top 10 for Large Language Model Applications and Agents, 2026 Edition. OWASP Foundation. \url{https://owasp.org/www-project-top-10-for-large-language-model-applications/}


\item E.~Gamma, R.~Helm, R.~Johnson, and J.~Vlissides. 1994. \emph{Design Patterns: Elements of Reusable Object-Oriented Software}. Addison-Wesley.


\item G.~Aspris and J.~Svec. 2025. Institutionalizing Risk Curation in Decentralized Credit. \emph{arXiv preprint arXiv:2512.11976}. \url{https://arxiv.org/abs/2512.11976}


\item S.~Jaffer. 1999. Microfinance and the Mechanics of Solidarity Lending. Harvard Law School, \emph{SSRN Working Paper No.~162548}. \url{https://papers.ssrn.com/sol3/papers.cfm?abstract_id=162548}


\item R.~Sherratt. 2016. From Solidarity to Coercion: The Dynamics of Group Liability. \emph{Oxford Economic Papers}, 68, 4, pages 955--974. \url{https://doi.org/10.1093/oep/gpw018}


\item H.~Fazliani, A.~Pasdar, and Y.~C.~Lee. 2025. Ensemble-Based Semi-Supervised Learning for Illicit Account Detection in Ethereum DeFi Transactions. \emph{arXiv preprint arXiv:2412.02408}, rev.~October. \url{https://arxiv.org/abs/2412.02408}


\item H.~Alhasawi, A.~Almutairi, and S.~Alharbi. 2025. FedFraud: A Federated Approach to Scalable and Trustworthy Financial Fraud Detection. \emph{Security and Privacy}, Wiley. \url{https://doi.org/10.1002/spy2.70012}


\item European Parliament and Council of the European Union. 2026. Regulation (EU) 2024/1689 laying down harmonised rules on artificial intelligence (AI Act). \emph{Official Journal of the European Union}, L~2024/1689, July 2024; high-risk credit-scoring provisions applicable from August. \url{https://eur-lex.europa.eu/eli/reg/2024/1689/oj}


\item L.~Dong, Y.~Qiu, and F.~Xu. 2023. Blockchain-Enabled Deep-Tier Supply Chain Finance. \emph{Manufacturing \& Service Operations Management}, 25, 6, pages 2021--2037. \url{https://doi.org/10.1287/msom.2022.1123}


\item U.~Bindseil. 2020. Tiered CBDC and the Financial System. \emph{ECB Working Paper Series}, no.~2351, January. \url{https://www.ecb.europa.eu/pub/pdf/scpwps/ecb.wp2351~c8c18bbd60.en.pdf}


\item J.~Kiff, J.~Alwazir, S.~Davidovic, A.~Farias, A.~Khan, T.~Khiaonarong, M.~Malaika, H.~Monroe, N.~Sugimoto, H.~Tourpe, and P.~Zhou. 2020. A Survey of Research on Retail Central Bank Digital Currency. \emph{IMF Working Paper WP/20/104}, June. \url{https://www.imf.org/en/Publications/WP/Issues/2020/06/26/A-Survey-of-Research-on-Retail-Central-Bank-Digital-Currency-49069}


\item R.~Auer and R.~B\"{o}hme. 2020. The Technology of Retail Central Bank Digital Currency. \emph{BIS Quarterly Review}, March; see also BIS Working Paper No.~948. \url{https://www.bis.org/publ/qtrpdf/r_qt2003j.htm}


\item L.~Ante. 2025. From Adoption to Continuance: Stablecoins in Cross-Border Remittances and the Role of Digital and Financial Literacy. \emph{Telematics and Informatics}, vol.~97, art.~102230. \url{https://doi.org/10.1016/j.tele.2024.102230}


\item Y.~Li, Y.~Xiang, Q.~Wang, T.~H.~Yuen, A.~Deppeler, and J.~Yu. 2026. SoK: Stablecoins in Retail Payments. \emph{arXiv preprint arXiv:2601.00196}. \url{https://arxiv.org/abs/2601.00196}


\item W.~Du, C.~Huang, and D.~Scharfstein. 2026. Competing Rails for Cross-Border Payments: Banks, Fintechs, and Stablecoins. Harvard Business School Working Paper, February. \url{https://www.hbs.edu/faculty/Pages/item.aspx?num=66992}


\item M.~Fr\"{o}hlich, J.~A.~Vega Vermehren, F.~Alt, and A.~Schmidt. 2022. Implementation and Evaluation of a Point-Of-Sale Payment System Using Bitcoin Lightning. In \emph{Proc. NordiCHI}. \url{http://www.medien.ifi.lmu.de/pubdb/publications/pub/froehlich2022nordichi/froehlich2022nordichi.pdf}


\item K.~Zhang, T.-M.~Choi, S.-H.~Chung, Y.~Dai, and X.~Wen. 2024. Blockchain Adoption in Retail Operations: Stablecoins and Traceability. \emph{European Journal of Operational Research}, 315, 1, pages 147--160. \url{https://doi.org/10.1016/j.ejor.2023.11.026}


\item J.~Lin, M.~Liu, S.~Li, and X.~Wang. 2025. SecurePay: Enabling Secure and Fast Payment Processing for Platform Economy. In \emph{Proc. IEEE/ACM Int. Symp. Quality of Service (IWQoS)}; \emph{arXiv:2505.17466}. \url{https://arxiv.org/abs/2505.17466}


\item Committee on Payments and Market Infrastructures. 2023. Considerations for the Use of Stablecoin Arrangements in Cross-Border Payments. Bank for International Settlements, July. \url{https://www.bis.org/cpmi/publ/d220.pdf}


\item A.~Kosse, J.~Lisack, and N.~Boakye-Agyeman. 2023. The Impact of Stablecoins on the International Monetary and Financial System. \emph{BIS Papers}, no.~170, February. \url{https://www.bis.org/publ/bppdf/bispap170.pdf}


\item E.~Cerutti, M.~Firat, and H.~P\'{e}rez-Saiz. 2025. Estimating the Impact of Digital Money on Cross-Border Flows: Scenario Analysis Covering the Intensive Margin. \emph{IMF Fintech Notes} 2025/002, February. \url{https://www.imf.org/-/media/files/publications/ftn063/2025/english/ftnea2025002.pdf}


\item S.~Chaliasos, M.~Charalambous, L.~Zhou, R.~Galanopoulou, A.~Gervais, D.~Mitropoulos, and B.~Livshits. 2024. Smart Contract and DeFi Security: Insights from Tool Evaluations and Practitioner Surveys. In \emph{Proc. IEEE/ACM 46th Int. Conf. Software Engineering (ICSE)}, pages 160:1--160:13. \url{https://doi.org/10.1145/3597503.3623302}


\item K.~Qin, L.~Zhou, B.~Livshits, and A.~Gervais. 2021. Attacking the DeFi Ecosystem with Flash Loans for Fun and Profit. In \emph{Financial Cryptography and Data Security (FC)}, pages 3--32. \url{https://doi.org/10.1007/978-3-662-63514-3_1}


\item T.~Mackinga, T.~Nadahalli, and R.~Wattenhofer. 2024. SoK: Attacks on DAOs. \emph{arXiv preprint arXiv:2406.15071}. \url{https://arxiv.org/abs/2406.15071}


\item E.~G.~Weyl, P.~Ohlhaver, and V.~Buterin. 2022. Decentralized Society: Finding Web3's Soul. \emph{SSRN Working Paper No.~4105763}. \url{https://papers.ssrn.com/sol3/papers.cfm?abstract_id=4105763}


\item SmartContractSecurity. 2020. SWC Registry: Smart Contract Weakness Classification and Test Cases. \url{https://swcregistry.io/}


\item G.~Caldarelli. 2025. Can Artificial Intelligence Solve the Blockchain Oracle Problem? Unpacking the Challenges and Possibilities. \emph{Frontiers in Blockchain}, vol.~8, art.~1682623. \url{https://doi.org/10.3389/fbloc.2025.1682623}


\item M.~Gurupriya, R.~Gayathri, C.~Yadav, and U.~Umamaheshwar. 2025. Blockchain-Powered Platform for Microloans and Lending Solutions. In \emph{Proc. Int. Conf. on Intelligent Computing, Instrumentation and Control Technologies (ICICT)}, pages 1--6. \url{https://doi.org/10.1109/ICTEST64710.2025.11042350}


\item V.~Buterin, Y.~Weiss, K.~Gazso, N.~Patel, D.~Tirosh, S.~Nacson, and T.~Hess. 2021. ERC-4337: Account Abstraction Using Alt Mempool. Ethereum Improvement Proposals. \url{https://eips.ethereum.org/EIPS/eip-4337}


\item Behaviour-Centric Cybersecurity Center (BCCC), York University. 2025. BCCC-DeFiFraudTrans-2025: A Labelled Ethereum DeFi Fraud Transaction Dataset. Accessed June 29, 2026. \url{https://www.yorku.ca/research/bccc/ucs-technical/cybersecurity-datasets-cds/defi-fraud-transactions-bccc-defifraudtrans-2025/}

\item P.~K.~Ozili. 2024. Central Bank Digital Currency Adoption Challenges. \emph{Journal of Central Banking Theory and Practice}, 13, 1, pages 5--24. \url{https://doi.org/10.2478/jcbtp-2024-0001}


\item McKinsey \& Company. 2018. Global Payments 2018: A Dynamic Industry Continues to Defy Predictions. McKinsey Global Payments Map. \url{https://www.mckinsey.com/industries/financial-services/our-insights/global-payments-a-dynamic-industry-continues-to-defy-predictions}


\item Y.~Wang, Q.~Zhang, K.~Li, Y.~Tang, J.~Chen, X.~Luo, and T.~Chen. 2021. iBatch: Saving Ethereum Fees via Secure and Cost-Effective Batching of Smart-Contract Invocations. In \emph{Proc. ACM Joint Eur. Software Engineering Conf. and Symp. Foundations of Software Engineering (ESEC/FSE)}, pages 777--789. \url{https://doi.org/10.1145/3468264.3468574}


\item World Inequality Lab. 2026. Exorbitant Privilege. \emph{World Inequality Report 2026}. \url{https://wir2026.wid.world/insight/exorbitant-privilege/}


\item G.~Wood. 2023. Ethereum: A Secure Decentralised Generalised Transaction Ledger (Yellow Paper). Ethereum Foundation. \url{https://ethereum.github.io/yellowpaper/paper.pdf}


\item Y.~Wang, J.~Liu, and Z.~Chen. 2020. ContractWard: Automated Vulnerability Detection Models for Ethereum Smart Contracts. \emph{IEEE Trans. Network Science and Engineering}, vol.~8, no.~2, pages 1133--1144. \url{https://doi.org/10.1109/TNSE.2020.2968505}


\item J.-W.~Liao, T.-T.~Tsai, C.-K.~He, and C.-W.~Tien. 2019. SoliAudit: Smart Contract Vulnerability Assessment Based on Machine Learning and Fuzz Testing. In \emph{Proc. 6th Int. Conf. Internet of Things: Systems, Management and Security (IOTSMS)}, pages 458--465. \url{https://doi.org/10.1109/iotsms48152.2019.8939256}


\item Federal Reserve Bank of New York and Bank for International Settlements Innovation Hub. 2025. Project Pine: Central Bank Open Market Operations with Smart Contracts. BIS Innovation Hub / NY Fed Joint Report, May. \url{https://www.bis.org/publ/othp95.htm}

\end{enumerate}
}





\end{document}
